% ---------- Page and fonts ----------
\documentclass[11pt,a4paper,openany,oneside]{book}
\usepackage[T1]{fontenc}
\usepackage[utf8]{inputenc}
\usepackage[sc]{mathpazo}
\linespread{1.08}
\usepackage[scaled=0.9]{helvet}
\usepackage{amsmath,amssymb,mathtools}
\usepackage{bm}
\usepackage{cancel}
\usepackage[a4paper,top=24mm,bottom=24mm,inner=24mm,outer=24mm,headsep=8mm,footskip=12mm]{geometry}
\usepackage{microtype}
\usepackage[dvipsnames,svgnames]{xcolor}
\usepackage{graphicx}
\usepackage{booktabs,array,tabularx}
\usepackage{multicol}
\usepackage{enumitem}
\setlist{nosep,leftmargin=1.4em,topsep=3pt}
\usepackage{tikz}
\usetikzlibrary{angles,quotes,arrows.meta,calc,decorations.pathreplacing,patterns,positioning,shapes.geometric,decorations.markings,3d,intersections}
\usepackage{pgfplots}
\pgfplotsset{compat=1.17}
\usepgfplotslibrary{fillbetween,groupplots}
\usepackage{tikz-3dplot}
\usepackage[most]{tcolorbox}
\usepackage{titlesec}
\usepackage{fancyhdr}
\usepackage{etoolbox}
\usepackage[hidelinks,unicode,psdextra,bookmarksnumbered=true,bookmarksopen=true,bookmarksopenlevel=1,pdfpagemode=UseOutlines]{hyperref}
\usepackage{bookmark}
\hypersetup{pdftitle={Classical Mechanics & Gravitation: The Math Book (Modules 1-4)},pdfauthor={ASD-STE100 Technical Reference},pdfsubject={Momentum, Rotation, Angular Momentum, and Gravitation}}

% ---------- Colours ----------
\definecolor{ink}{HTML}{1F2A33}
\definecolor{teal}{HTML}{0F6E6E}
\definecolor{tealsoft}{HTML}{E8F3F2}
\definecolor{amber}{HTML}{B4690E}
\definecolor{ambersoft}{HTML}{FFF6E8}
\definecolor{plum}{HTML}{6A3D7A}
\definecolor{plumsoft}{HTML}{F5EFF7}
\definecolor{rose}{HTML}{B03A48}
\definecolor{rosesoft}{HTML}{FCEEF0}
\definecolor{slate}{HTML}{5B6770}
\definecolor{slatesoft}{HTML}{F3F5F6}
\definecolor{bluec}{HTML}{2459A8}
\color{ink}

% ---------- Headings ----------
\newcommand{\sffamilybold}{\fontfamily{phv}\selectfont\bfseries}
\titleformat{\part}[display]{\centering\sffamilybold\color{teal}}{\Large\MakeUppercase{\thepart}}{8pt}{\Huge}
\renewcommand{\thepart}{Module \arabic{part}}
\usepackage{titletoc}
\titlecontents{part}[0pt]{\addvspace{12pt}\fontfamily{phv}\selectfont\bfseries\color{teal}}{\thecontentslabel\quad}{}{\hfill\contentspage}[\addvspace{3pt}]
\titleformat{\chapter}[display]
  {\sffamilybold\color{ink}}
  {\color{teal}\large CHAPTER \thechapter}{4pt}{\huge}[\vspace{2pt}{\color{teal}\titlerule[1.2pt]}]
\titlespacing*{\chapter}{0pt}{0pt}{18pt}
\titleformat{\section}{\sffamilybold\Large\color{teal}}{\thesection}{0.7em}{}
\titleformat{\subsection}{\sffamilybold\large\color{ink}}{\thesubsection}{0.7em}{}
\titleformat{\subsubsection}{\sffamilybold\normalsize\color{slate}}{}{0em}{}
\titlespacing*{\section}{0pt}{16pt}{7pt}
\titlespacing*{\subsection}{0pt}{13pt}{5pt}
\titlespacing*{\subsubsection}{0pt}{9pt}{3pt}
\setcounter{secnumdepth}{2}
\setcounter{tocdepth}{2}

% ---------- Header / footer ----------
\pagestyle{fancy}
\fancyhf{}
\renewcommand{\headrulewidth}{0.3pt}
\renewcommand{\chaptermark}[1]{\markboth{#1}{}}
\renewcommand{\sectionmark}[1]{\markright{\thesection\ \ #1}}
\fancyhead[L]{\small\fontfamily{phv}\selectfont\color{slate}\leftmark}
\fancyhead[R]{\small\fontfamily{phv}\selectfont\color{slate}\rightmark}
\fancyfoot[C]{\small\fontfamily{phv}\selectfont\color{slate}\thepage}
\fancypagestyle{plain}{\fancyhf{}\renewcommand{\headrulewidth}{0pt}\fancyfoot[C]{\small\fontfamily{phv}\selectfont\color{slate}\thepage}}

\setlength{\parindent}{0pt}
\setlength{\parskip}{5pt}

% ---------- Boxes ----------
\tcbset{common/.style={enhanced,breakable,boxrule=0pt,arc=2mm,left=3mm,right=3mm,top=2mm,bottom=2mm,
  fonttitle=\fontfamily{phv}\selectfont\bfseries\small,coltitle=white,before skip=7pt,after skip=7pt}}
% Goal of a derivation
\newtcolorbox{goal}{common,colback=tealsoft,colframe=teal,borderline west={3pt}{0pt}{teal},
  title={GOAL},attach boxed title to top left={xshift=3mm,yshift=-2mm},boxed title style={colback=teal,arc=1mm}}
% Plus-two base facts
\newtcolorbox{base}[1][]{common,colback=ambersoft,colframe=amber,borderline west={3pt}{0pt}{amber},
  title={FROM YOUR +2 BASE#1},attach boxed title to top left={xshift=3mm,yshift=-2mm},boxed title style={colback=amber,arc=1mm}}
% Final result
\newtcolorbox{result}[1][]{common,breakable=false,colback=white,colframe=teal,boxrule=1.2pt,
  title={RESULT#1},attach boxed title to top left={xshift=3mm,yshift=-2mm},boxed title style={colback=teal,arc=1mm}}
% Meaning
\newtcolorbox{means}{common,colback=plumsoft,colframe=plum,borderline west={3pt}{0pt}{plum},
  title={WHAT IT MEANS},attach boxed title to top left={xshift=3mm,yshift=-2mm},boxed title style={colback=plum,arc=1mm}}
% Exam trap
\newtcolorbox{trap}{common,colback=rosesoft,colframe=rose,borderline west={3pt}{0pt}{rose},
  title={WATCH OUT},attach boxed title to top left={xshift=3mm,yshift=-2mm},boxed title style={colback=rose,arc=1mm}}
% Worked example
\newtcolorbox{example}[2]{common,breakable=false,colback=slatesoft,colframe=slate,borderline west={3pt}{0pt}{slate},
  before upper={\belowpdfbookmark{Book Example #1}{ex#1}},
  title={BOOK EXAMPLE\ \ #1\quad #2},attach boxed title to top left={xshift=3mm,yshift=-2mm},boxed title style={colback=slate,arc=1mm}}
% Idea / picture first
\newtcolorbox{idea}{common,colback=white,colframe=teal!50,boxrule=0.6pt,borderline west={3pt}{0pt}{teal!50},
  title={THE IDEA},coltitle=teal,attach boxed title to top left={xshift=3mm,yshift=-2mm},boxed title style={colback=white,colframe=teal!50,boxrule=0.6pt,arc=1mm}}
% Syllabus strip
\newtcolorbox{syll}{enhanced,colback=slatesoft,colframe=slatesoft,arc=1mm,left=3mm,right=3mm,top=1.5mm,bottom=1.5mm,fontupper=\small\fontfamily{phv}\selectfont}

% ---------- Derivation helpers ----------
% right-side reason in align: use  &&\why{reason}
\newcommand{\why}[1]{\text{\small\color{slate}\itshape #1}}
\newcommand{\step}[1]{\par\smallskip{\fontfamily{phv}\selectfont\bfseries\small\color{teal}#1}\enspace}
\newcommand{\term}[1]{\textbf{\color{teal}#1}}
\newcommand{\eqbox}[1]{\boxed{\,#1\,}}
\newcommand{\vect}[1]{\mathbf{#1}}
\newcommand{\uv}[1]{\hat{\mathbf{#1}}}
\newcommand{\ihat}{\hat{\mathbf{i}}}
\newcommand{\jhat}{\hat{\mathbf{j}}}
\newcommand{\khat}{\hat{\mathbf{k}}}
\newcommand{\dd}{\,\mathrm{d}}
\newcommand{\diff}{\mathrm{d}}
\newcommand{\Eloc}{E_{\text{loc}}}
\newcommand{\kT}{k_BT}
\newcommand{\eps}{\varepsilon}
\newcommand{\bookref}[1]{{\small\color{slate}[Book: #1]}}
\newcounter{figno}[chapter]
\newcommand{\figcap}[1]{\par\nopagebreak[4]\smallskip\refstepcounter{figno}{\centering\small\color{slate}\fontfamily{phv}\selectfont \textbf{Fig.\ \thechapter.\arabic{figno}}\quad #1\par}\medskip}

% figure environment wrapper, no float (keeps math next to picture)
\newenvironment{pic}{\par\medskip\noindent\begin{minipage}{\linewidth}\centering}{\end{minipage}\par\medskip}
\usepackage{needspace}
\pretocmd{\section}{\needspace{8\baselineskip}}{}{}
\pretocmd{\subsection}{\needspace{6\baselineskip}}{}{}
\pretocmd{\subsubsection}{\needspace{4\baselineskip}}{}{}

% equation numbering per chapter
\numberwithin{equation}{chapter}

% TikZ styles
\tikzset{
  atom/.style={circle,fill=teal,inner sep=0pt,minimum size=5pt},
  atomB/.style={circle,fill=amber,inner sep=0pt,minimum size=5pt},
  latt/.style={circle,fill=ink,inner sep=0pt,minimum size=3pt},
  vec/.style={-{Stealth[length=2.6mm]},thick},
  ax/.style={-{Stealth[length=2mm]},thin,slate},
  dim/.style={{Stealth[length=1.6mm]}-{Stealth[length=1.6mm]},thin,slate},
  lbl/.style={font=\small},
}

% ---------- Toolkit items (T1, T2, ...) ----------
\newcounter{tool}
\newtcolorbox{toolbox}[2]{common,breakable=false,colback=white,colframe=amber!70,boxrule=0.6pt,
  borderline west={3pt}{0pt}{amber},
  title={\refstepcounter{tool}\label{#2}T\thetool\quad #1},coltitle=amber!80!black,
  before upper={\belowpdfbookmark{T\thetool\ \ #1}{tool\thetool}},
  attach boxed title to top left={xshift=3mm,yshift=-2mm},
  boxed title style={colback=white,colframe=amber!70,boxrule=0.6pt,arc=1mm}}
\newcommand{\tref}[1]{\hyperref[#1]{\textbf{\color{amber!80!black}T\ref*{#1}}}}

% ---------- 3D oblique projection for cells ----------
\tikzset{
  cub/.style={x={(1cm,0cm)},y={(0.42cm,0.30cm)},z={(0cm,1cm)}},
}
\newcommand{\cubeedges}{%
    \draw[thick,ink] (0,0,0)--(1,0,0)--(1,0,1)--(0,0,1)--cycle;
    \draw[thick,ink] (1,0,0)--(1,1,0)--(1,1,1)--(1,0,1);
    \draw[thick,ink] (0,0,1)--(0,1,1)--(1,1,1);
    \draw[ink,dashed] (0,0,0)--(0,1,0)--(1,1,0);
    \draw[ink,dashed] (0,1,0)--(0,1,1);}
\newcommand{\cubecorners}{%
    \foreach \x in {0,1}{\foreach \y in {0,1}{\foreach \z in {0,1}{\fill[teal] (\x,\y,\z) circle (2.2pt);}}}}
\raggedbottom


% ---------- Formula sheets ----------
\newcommand{\formulasheet}[1]{%
  \cleardoublepage\phantomsection
  \chapter*{#1}\addcontentsline{toc}{chapter}{#1}\markboth{#1}{}%
}
\newtcolorbox{fs}[1]{enhanced,breakable=false,colback=white,colframe=teal!60,boxrule=0.5pt,arc=1.5mm,
  left=2mm,right=2mm,top=1mm,bottom=1mm,before skip=4pt,after skip=4pt,
  title={#1},fonttitle=\fontfamily{phv}\selectfont\bfseries\footnotesize,colbacktitle=tealsoft,coltitle=teal}

\newtcolorbox{exampledemo}{common,colback=slatesoft,colframe=slate,borderline west={3pt}{0pt}{slate},
  title={BOOK EXAMPLE\ \ 1.1},attach boxed title to top left={xshift=3mm,yshift=-2mm},boxed title style={colback=slate,arc=1mm}}

\newtcolorbox{also}{common,colback=tealsoft!40,colframe=teal!70,borderline west={3pt}{0pt}{teal!70},
  title={ALSO ACCEPTED},attach boxed title to top left={xshift=3mm,yshift=-2mm},boxed title style={colback=teal!70,arc=1mm}}
\newcommand{\bk}[1]{\texorpdfstring{\enspace{\normalfont\small\color{slate}[Book #1]}}{ [Book #1]}}
\newcommand{\extra}{\texorpdfstring{\enspace{\normalfont\small\color{slate}(extra)}}{ (extra)}}
\newcommand{\qual}{\texorpdfstring{\enspace{\normalfont\small\color{slate}(qualitative)}}{ (qualitative)}}
% second view: x toward viewer (left-down), y right, z up; origin hidden at back
\tikzset{cubB/.style={x={(-0.55cm,-0.45cm)},y={(1cm,0cm)},z={(0cm,1cm)}}}
\newcommand{\cubeedgesB}[1][1]{%
  \draw[thick,ink] (1,0,0)--(1,#1,0)--(1,#1,1)--(1,0,1)--cycle;
  \draw[thick,ink] (1,#1,0)--(0,#1,0)--(0,#1,1)--(1,#1,1);
  \draw[thick,ink] (1,0,1)--(0,0,1)--(0,#1,1);
  \draw[ink,dashed] (0,0,0)--(1,0,0); \draw[ink,dashed] (0,0,0)--(0,#1,0); \draw[ink,dashed] (0,0,0)--(0,0,1);}

\begin{document}

% ============================================================================
% TITLE PAGE
% ============================================================================
\begin{titlepage}
\centering
\vspace*{1.5cm}
{\fontfamily{phv}\selectfont\bfseries\Large\color{teal} ASD-STE100 SIMPLIFIED TECHNICAL ENGLISH EDITION\par}
\vspace{0.6cm}
{\fontfamily{phv}\selectfont\bfseries\Huge\color{ink} CLASSICAL MECHANICS \& GRAVITATION:\\[0.3em] THE MATH BOOK\par}
\vspace{0.5cm}
{\fontfamily{phv}\selectfont\Large\color{slate} Complete Step-by-Step Derivations from $+2$ Base Knowledge,\\Vector Diagrams, All Book Examples, and Unique Solved Problems\par}
\vspace{1.2cm}
\begin{tikzpicture}[scale=1.05]
  % Central decorative SVG-style vector graphic: Orbit + Collision + Gyroscope
  \draw[teal!30,fill=tealsoft!50,rounded corners=4mm] (-5.5,-2.2) rectangle (5.5,2.2);
  % Left: 2D Collision
  \draw[ax] (-4.8,-1.2) -- (-1.8,-1.2) node[right,font=\scriptsize]{$x$};
  \draw[ax] (-3.5,-1.8) -- (-3.5,1.5) node[above,font=\scriptsize]{$y$};
  \fill[teal] (-4.5,0) circle (5pt) node[above=4pt,font=\scriptsize\bfseries]{$m_1$};
  \draw[vec,teal] (-4.5,0) -- (-3.7,0) node[midway,below,font=\scriptsize]{$\vect{v}_1$};
  \fill[amber] (-3.5,0) circle (6pt);
  \draw[vec,teal] (-3.5,0) -- (-2.3,1.0) node[right,font=\scriptsize]{$\vect{v}_1'$};
  \draw[vec,amber] (-3.5,0) -- (-2.2,-0.8) node[right,font=\scriptsize]{$\vect{v}_2'$};
  % Center: Rigid body & Angular momentum
  \draw[thick,plum,fill=plumsoft] (0,0) ellipse (1.1cm and 0.42cm);
  \draw[vec,plum,ultra thick] (0,-1.3) -- (0,1.4) node[above,font=\scriptsize\bfseries]{$\vect{L}=I\bm{\omega}$};
  \draw[vec,rose] (0,0) -- (0.95,-0.2) node[below,font=\scriptsize]{$\vect{r}_i$};
  \fill[ink] (0.95,-0.2) circle (2.5pt);
  % Right: Kepler Orbit & Gravity
  \draw[thick,teal,dashed] (3.4,0) ellipse (1.6cm and 1.0cm);
  \fill[amber] (2.4,0) circle (6pt) node[below=5pt,font=\scriptsize\bfseries]{$M_S$};
  \fill[teal] (4.7,0.58) circle (3.5pt) node[above right,font=\scriptsize]{$m$};
  \draw[vec,rose] (4.7,0.58) -- (3.7,0.33) node[above,font=\scriptsize]{$\vect{F}_g$};
  \draw[vec,bluec] (4.7,0.58) -- (4.35,1.25) node[above,font=\scriptsize]{$\vect{v}$};
\end{tikzpicture}
\vspace{1.2cm}

\begin{tcolorbox}[common,colback=slatesoft,colframe=teal,borderline west={3pt}{0pt}{teal},title={SYLLABUS \& SOURCE COVERAGE}]
\begin{itemize}
  \item \textbf{Module 1 (PDF 1):} Linear Momentum, Impulse, Collisions, Center of Mass, Reduced Mass, and Rocket Propulsion \emph{(Young \& Freedman Ch.\ 8; Analytical Mechanics Ch.\ 6)}.
  \item \textbf{Module 2 (PDF 4):} Rotation of Rigid Bodies, Rotational Kinematics, Rotational Kinetic Energy, and Moment-of-Inertia Integrals \emph{(Young \& Freedman Ch.\ 9)}.
  \item \textbf{Module 3 (PDF 3):} Dynamics of Rotational Motion, Rolling Motion, Angular Momentum, Gyroscopes, and Axial Vectors \emph{(Young \& Freedman Ch.\ 10; Analytical Mechanics Ch.\ 7)}.
  \item \textbf{Module 4 (PDF 2):} Newton's Law of Gravitation, Satellites, Kepler's Laws, Spherical Shells, Gauss's Law for Gravity, Poisson and Laplace Equations, Black Holes, and General Relativity \emph{(Young \& Freedman Ch.\ 13; Analytical Mechanics Ch.\ 9)}.
\end{itemize}
\end{tcolorbox}
\vfill
{\small\fontfamily{phv}\selectfont\color{slate} Written in strict compliance with ASD-STE100 Simplified Technical English.\par}
\end{titlepage}

\tableofcontents

% ============================================================================
% CHAPTER 0: PLUS-TWO MATHEMATICAL TOOLKIT
% ============================================================================
\chapter*{Plus-Two Mathematical Toolkit}
\addcontentsline{toc}{chapter}{Plus-Two Mathematical Toolkit}
\markboth{Plus-Two Mathematical Toolkit}{Plus-Two Mathematical Toolkit}

\begin{syll}
Use these six foundational tools (\textbf{T1}--\textbf{T6}) from $+2$ mathematics and physics to follow every derivation in this book.
\end{syll}

\begin{toolbox}{Derivatives, Product Rule, and Kinematic Rates}{tool:deriv}
Let position be $\vect{r}(t) = x(t)\ihat + y(t)\jhat + z(t)\khat$.
\begin{align*}
  \vect{v} &= \frac{\diff\vect{r}}{\diff t}, \qquad \vect{a} = \frac{\diff\vect{v}}{\diff t} = \frac{\diff^2\vect{r}}{\diff t^2}, \qquad \frac{\diff}{\diff t}(uv) = \frac{\diff u}{\diff t}v + u\frac{\diff v}{\diff t} &&\why{kinematic derivatives and product rule}\\
  \frac{\diff}{\diff x}\!\left(x^n\right) &= nx^{n-1}, \qquad \frac{\diff}{\diff x}(\sin x) = \cos x, \qquad \frac{\diff}{\diff x}(\cos x) = -\sin x &&\why{standard $+2$ derivatives}
\end{align*}
\end{toolbox}

\begin{toolbox}{Definite Integrals and Logarithms}{tool:integ}
Integration is the inverse of differentiation. It sums infinitesimal parts from limit $a$ to limit $b$:
\begin{align*}
  \int_a^b x^n\dd x &= \left[\frac{x^{n+1}}{n+1}\right]_a^b = \frac{b^{n+1}-a^{n+1}}{n+1}\quad (n\neq -1), \qquad \int_a^b \frac{\diff x}{x} = [\ln x]_a^b = \ln\!\left(\frac{b}{a}\right)
\end{align*}
\end{toolbox}

\begin{toolbox}{Scalar (Dot) Product and Work Integral}{tool:dot}
For two vectors $\vect{A}$ and $\vect{B}$ with angle $\phi$ between them:
\begin{align*}
  \vect{A}\cdot\vect{B} &= AB\cos\phi = A_x B_x + A_y B_y + A_z B_z, \qquad \vect{A}\cdot\vect{A} = A^2 &&\why{dot product}\\
  W &= \int_{\vect{r}_1}^{\vect{r}_2} \vect{F}\cdot\diff\vect{r} = K_2 - K_1 = -\Delta U \quad(\text{for conservative }\vect{F}) &&\why{work--energy theorem}
\end{align*}
\end{toolbox}

\begin{toolbox}{Vector (Cross) Product and BAC--CAB Identity}{tool:cross}
The cross product $\vect{C} = \vect{A}\times\vect{B}$ is perpendicular to both $\vect{A}$ and $\vect{B}$ by the right-hand rule:
\begin{align*}
  |\vect{A}\times\vect{B}| &= AB\sin\phi, \qquad \vect{A}\times\vect{B} = -\vect{B}\times\vect{A}, \qquad \vect{A}\times\vect{A} = \vect{0}\\
  \vect{A}\times(\vect{B}\times\vect{C}) &= \vect{B}(\vect{A}\cdot\vect{C}) - \vect{C}(\vect{A}\cdot\vect{B}) &&\why{BAC--CAB vector triple product}
\end{align*}
\end{toolbox}

\begin{toolbox}{Vector Calculus Operators: Gradient, Divergence, Curl, and Gauss's Theorem}{tool:vec_calc}
Define the del operator $\nabla = \ihat\frac{\partial}{\partial x} + \jhat\frac{\partial}{\partial y} + \khat\frac{\partial}{\partial z}$.
\begin{align*}
  \text{Gradient: } \nabla\Phi &= \ihat\frac{\partial\Phi}{\partial x} + \jhat\frac{\partial\Phi}{\partial y} + \khat\frac{\partial\Phi}{\partial z}, \qquad \text{Laplacian: } \nabla^2\Phi = \nabla\cdot(\nabla\Phi) = \frac{\partial^2\Phi}{\partial x^2}+\frac{\partial^2\Phi}{\partial y^2}+\frac{\partial^2\Phi}{\partial z^2}\\
  \text{Curl of gradient: } \nabla\times(\nabla\Phi) &= \vect{0}, \qquad \text{Gauss's Divergence Theorem: } \oint_S \vect{A}\cdot\uv{n}\dd S = \int_V (\nabla\cdot\vect{A})\dd\tau
\end{align*}
\end{toolbox}

\begin{toolbox}{Coordinate Elements in Cylindrical and Spherical Geometry}{tool:coords}
\begin{itemize}
  \item \textbf{Polar / Cylindrical $(r,\phi,z)$:} Ring circumference $= 2\pi r$, disk area element $\dd A = r\dd r\dd\phi$ (or $2\pi r\dd r$), cylindrical shell volume $\dd V = 2\pi r L\dd r$.
  \item \textbf{Spherical $(r,\theta,\phi)$:} Ring radius on sphere of radius $R$ is $R\sin\theta$, ring width is $R\dd\theta$, zone area element $\dd A = 2\pi R^2\sin\theta\dd\theta$, solid angle element $\dd\Omega = \frac{\dd S\cos\theta}{r^2}$ (total sphere $\oint\dd\Omega = 4\pi\text{ sr}$).
\end{itemize}
\end{toolbox}

% ============================================================================
% MODULE 1: LINEAR MOMENTUM, IMPULSE, COLLISIONS, AND ROCKET PROPULSION
% ============================================================================
\part{Linear Momentum, Impulse, Collisions, and Rockets}

\chapter{Momentum, Impulse, and Conservation Laws}
\begin{syll}
\textbf{Syllabus Coverage (PDF 1):} Young \& Freedman Chapter 8 (Sections 8.1--8.2) and Analytical Mechanics Chapter 6 (Sections 6.1, 6.5, 6.6). Newton's second law in momentum form, impulse--momentum theorem, and conservation of total momentum for $N$-particle systems.
\end{syll}

\section{Linear Momentum and Newton's Second Law\bk{8.1, 6.1}}

\begin{goal}
Define linear momentum $\vect{p}$ and express Newton's second law for both constant-mass and variable-mass particles.
\end{goal}

\begin{base}
A particle of mass $m$ moves with velocity $\vect{v} = \frac{\diff\vect{r}}{\diff t}$ and acceleration $\vect{a} = \frac{\diff\vect{v}}{\diff t}$. Newton's second law for constant mass is $\sum\vect{F} = m\vect{a}$.
\end{base}

\step{Step 1: Define linear momentum.}
Multiply the particle mass $m$ by its velocity vector $\vect{v}$:
\begin{equation}\label{eq:p_def}
  \eqbox{\vect{p} = m\vect{v}} \qquad (p_x = mv_x,\quad p_y = mv_y,\quad p_z = mv_z)
\end{equation}

\step{Step 2: Differentiate $\vect{p}$ with respect to time $t$.}
Apply the product rule (\tref{tool:deriv}) to $\vect{p} = m\vect{v}$:
\begin{align}
  \frac{\diff\vect{p}}{\diff t} &= \frac{\diff}{\diff t}(m\vect{v}) = m\frac{\diff\vect{v}}{\diff t} + \vect{v}\frac{\diff m}{\diff t} &&\why{general variable-mass form (Eq.\ 6.1)}\label{eq:F_varmass}
\end{align}
If the mass $m$ is constant, $\frac{\diff m}{\diff t} = 0$. Substitute $\frac{\diff\vect{v}}{\diff t} = \vect{a}$:
\begin{align}
  \sum\vect{F} &= m\frac{\diff\vect{v}}{\diff t} = \frac{\diff}{\diff t}(m\vect{v}) = \frac{\diff\vect{p}}{\diff t} &&\why{constant mass $m$}\label{eq:N2_momentum}
\end{align}

\step{Step 3: Relate kinetic energy $K$ to momentum magnitude $p$.}
Start from $K = \frac{1}{2}mv^2$ and substitute $v = \frac{p}{m}$:
\begin{equation}\label{eq:K_p_rel}
  K = \frac{1}{2}m\!\left(\frac{p}{m}\right)^2 = \eqbox{\frac{p^2}{2m}} \iff p = \sqrt{2mK}
\end{equation}

\begin{result}[Newton's Second Law in Terms of Momentum]
\begin{equation*}
  \sum\vect{F} = \frac{\diff\vect{p}}{\diff t} \quad (\text{constant } m), \qquad \vect{F} = \frac{\diff\vect{p}}{\diff t} = m\frac{\diff\vect{v}}{\diff t} + \vect{v}\frac{\diff m}{\diff t} \quad (\text{variable } m)
\end{equation*}
\end{result}

\section{The Impulse--Momentum Theorem\bk{8.1, 6.5}}

\begin{goal}
Derive the relation between the time integral of net force (impulse $\vect{J}$) and the change in linear momentum $\Delta\vect{p}$.
\end{goal}

\begin{pic}
\begin{tikzpicture}[scale=0.95]
  \begin{axis}[
    width=8.5cm, height=4.5cm,
    axis lines=middle,
    xlabel={$t$}, ylabel={$\sum F_x$},
    xmin=0, xmax=5, ymin=0, ymax=3.2,
    xtick={1,4}, xticklabels={$t_1$,$t_2$},
    ytick={1.35}, yticklabels={$(F_{\text{av}})_x$},
    clip=false
  ]
    \addplot[name path=curve,smooth,thick,rose,domain=1:4,samples=60] {2.8*sin(deg((x-1)*pi/3))^2};
    \path[name path=axis] (axis cs:1,0) -- (axis cs:4,0);
    \addplot[rosesoft] fill between[of=curve and axis];
    \draw[thick,teal,fill=tealsoft,fill opacity=0.45] (axis cs:1,0) rectangle (axis cs:4,1.4);
    \node[font=\footnotesize,rose!80!black] at (axis cs:2.5,2.1) {$\text{Area} = J_x = \int_{t_1}^{t_2}\sum F_x\dd t$};
    \node[font=\footnotesize,teal!80!black] at (axis cs:2.5,0.65) {$(F_{\text{av}})_x(t_2-t_1)$};
  \end{axis}
\end{tikzpicture}
\figcap{The area under the curve of net force $\sum F_x$ versus time $t$ equals the impulse $J_x$.}
\end{pic}

\step{Step 1: Integrate Newton's second law over time.}
Start from $\sum\vect{F} = \frac{\diff\vect{p}}{\diff t}$ and multiply both sides by $\diff t$:
\begin{align}
  \sum\vect{F}\dd t &= \diff\vect{p} &&\why{differential form}\\
  \int_{t_1}^{t_2}\sum\vect{F}\dd t &= \int_{\vect{p}_1}^{\vect{p}_2}\diff\vect{p} = \vect{p}_2 - \vect{p}_1 = \Delta\vect{p} &&\why{integrate from $t_1$ to $t_2$ (\tref{tool:integ})}
\end{align}

\step{Step 2: Define impulse $\vect{J}$ and average force $\vect{F}_{\text{av}}$.}
Define the impulse $\vect{J}$ as the time integral of the net force:
\begin{equation}\label{eq:impulse_def}
  \eqbox{\vect{J} = \int_{t_1}^{t_2}\sum\vect{F}\dd t = \vect{p}_2 - \vect{p}_1 = m\vect{v}_2 - m\vect{v}_1}
\end{equation}
If $\sum\vect{F}$ is constant, $\vect{J} = \sum\vect{F}(t_2 - t_1) = \sum\vect{F}\Delta t$. For a varying force, define the average net force $\vect{F}_{\text{av}}$ so that it gives the same impulse:
\begin{equation}\label{eq:impulse_av}
  \vect{J} = \vect{F}_{\text{av}}(t_2 - t_1), \qquad
  \begin{cases}
    J_x = \int_{t_1}^{t_2}\sum F_x\dd t = (F_{\text{av}})_x(t_2-t_1) = mv_{2x} - mv_{1x}\\[4pt]
    J_y = \int_{t_1}^{t_2}\sum F_y\dd t = (F_{\text{av}})_y(t_2-t_1) = mv_{2y} - mv_{1y}
  \end{cases}
\end{equation}

\begin{means}
Impulse $\vect{J} = \int\sum\vect{F}\dd t$ measures force acting over \textbf{time} and changes \textbf{momentum} ($\Delta\vect{p} = \vect{J}$). Work $W = \int\sum\vect{F}\cdot\diff\vect{r}$ measures force acting over \textbf{distance} and changes \textbf{kinetic energy} ($\Delta K = W$).
\end{means}

\section{Conservation of Momentum for a System of Particles\bk{8.2, 6.6}}

\begin{goal}
Prove that internal forces cancel in pairs and that the total linear momentum $\vect{P}_{\text{tot}}$ of a system is conserved when the net external force is zero.
\end{goal}

\begin{pic}
\begin{tikzpicture}[scale=1.0]
  \fill[teal] (0,0) circle (5pt) node[below=5pt,lbl]{Particle $i$ ($m_i$)};
  \fill[amber] (3.5,0.8) circle (5pt) node[below=5pt,lbl]{Particle $j$ ($m_j$)};
  \draw[dashed,slate] (0,0) -- (3.5,0.8);
  \draw[vec,rose,ultra thick] (0,0) -- (1.3,0.3) node[above,font=\small]{$\vect{F}_{ij}$};
  \draw[vec,rose,ultra thick] (3.5,0.8) -- (2.2,0.5) node[above,font=\small]{$\vect{F}_{ji} = -\vect{F}_{ij}$};
  \draw[vec,bluec] (0,0) -- (-1.1,0.7) node[left,font=\small]{$\vect{F}_i^{(e)}$};
  \draw[vec,bluec] (3.5,0.8) -- (4.5,1.6) node[right,font=\small]{$\vect{F}_j^{(e)}$};
\end{tikzpicture}
\figcap{Internal action--reaction pair $\vect{F}_{ij} = -\vect{F}_{ji}$ and external forces $\vect{F}_i^{(e)}, \vect{F}_j^{(e)}$ on particles $i$ and $j$.}
\end{pic}

\step{Step 1: Write Newton's second law for particle $i$ in an $N$-particle system.}
Let $\vect{F}_i^{(e)}$ be the external force on particle $i$, and let $\vect{F}_{ij}$ be the internal force exerted on particle $i$ by particle $j$ ($j\neq i$):
\begin{equation}
  \frac{\diff\vect{p}_i}{\diff t} = \vect{F}_i^{(e)} + \sum_{\substack{j=1\\j\neq i}}^N \vect{F}_{ij}
\end{equation}

\step{Step 2: Sum over all $N$ particles.}
Define total momentum $\vect{P}_{\text{tot}} = \sum_{i=1}^N \vect{p}_i = \sum_{i=1}^N m_i\vect{v}_i$:
\begin{equation}
  \sum_{i=1}^N \frac{\diff\vect{p}_i}{\diff t} = \frac{\diff\vect{P}_{\text{tot}}}{\diff t} = \sum_{i=1}^N \vect{F}_i^{(e)} + \sum_{i=1}^N\sum_{\substack{j=1\\j\neq i}}^N \vect{F}_{ij}
\end{equation}

\step{Step 3: Apply Newton's third law to the double sum.}
For every pair $(i,j)$, Newton's third law gives $\vect{F}_{ij} = -\vect{F}_{ji}$, so $\vect{F}_{ij} + \vect{F}_{ji} = \vect{0}$. Thus the entire double sum over internal forces vanishes:
\begin{equation}\label{eq:dPtot_dt}
  \frac{\diff\vect{P}_{\text{tot}}}{\diff t} = \sum_{i=1}^N \vect{F}_i^{(e)} = \vect{F}_{\text{tot}}^{(e)}
\end{equation}

\step{Step 4: State the conservation condition.}
If the vector sum of external forces is zero ($\vect{F}_{\text{tot}}^{(e)} = \vect{0}$), then $\frac{\diff\vect{P}_{\text{tot}}}{\diff t} = \vect{0}$:
\begin{equation}\label{eq:P_conserv}
  \eqbox{\vect{P}_{\text{tot}} = m_A\vect{v}_A + m_B\vect{v}_B + \cdots = \text{constant} \qquad (\text{when }\sum\vect{F}_{\text{ext}} = \vect{0})}
\end{equation}

\begin{trap}
Momentum is a \textbf{vector}. You must conserve each component separately: $P_x = \text{const}$, $P_y = \text{const}$, $P_z = \text{const}$. Never add momentum magnitudes directly unless the vectors are parallel.
\end{trap}

% ============================================================================
% CHAPTER 2: COLLISIONS, CENTER OF MASS, AND ROCKET PROPULSION
% ============================================================================
\chapter{Collisions, Center of Mass, and Rocket Propulsion}
\begin{syll}
\textbf{Syllabus Coverage (PDF 1):} Young \& Freedman Chapter 8 (Sections 8.3--8.6) and Analytical Mechanics Chapter 6 (Sections 6.2--6.4, 6.7). Elastic, inelastic, and completely inelastic collisions; 2D glancing collisions; coefficient of restitution; center of mass; reduced mass; and rocket propulsion.
\end{syll}

\section{Completely Inelastic Collisions\bk{8.3, 6.1}}

\begin{goal}
Derive the common final velocity $\vect{v}_2$ and the fraction of kinetic energy lost when two colliding bodies stick together.
\end{goal}

\step{Step 1: Apply momentum conservation.}
Bodies $m_A$ and $m_B$ stick together after impact ($\vect{v}_{A2} = \vect{v}_{B2} = \vect{v}_2$):
\begin{equation}\label{eq:comp_inel_v}
  m_A\vect{v}_{A1} + m_B\vect{v}_{B1} = (m_A + m_B)\vect{v}_2 \implies \eqbox{\vect{v}_2 = \frac{m_A\vect{v}_{A1} + m_B\vect{v}_{B1}}{m_A + m_B}}
\end{equation}

\step{Step 2: Specialize to target $B$ initially at rest ($v_{B1x} = 0$).}
\begin{equation}\label{eq:v2x_rest}
  v_{2x} = \frac{m_A}{m_A + m_B}v_{A1x}
\end{equation}
Compute the initial and final kinetic energies $K_1$ and $K_2$:
\begin{align}
  K_1 &= \frac{1}{2}m_A v_{A1x}^2, \qquad
  K_2 = \frac{1}{2}(m_A + m_B)v_{2x}^2 = \frac{1}{2}(m_A + m_B)\!\left(\frac{m_A}{m_A + m_B}\right)^2 v_{A1x}^2 = \frac{m_A}{m_A + m_B}K_1
\end{align}
\begin{equation}\label{eq:K_ratio_inel}
  \eqbox{\frac{K_2}{K_1} = \frac{m_A}{m_A + m_B} < 1, \qquad \Delta K = K_2 - K_1 = -\frac{1}{2}\frac{m_A m_B}{m_A + m_B}v_{A1x}^2}
\end{equation}

\section{One-Dimensional Elastic Collisions\bk{8.4, 6.3}}

\begin{goal}
Derive the relative-velocity reversal law and the final velocities $v_{A2x}, v_{B2x}$ for a 1D elastic collision.
\end{goal}

\begin{pic}
\begin{tikzpicture}[scale=0.95]
  % Before
  \node[font=\footnotesize\bfseries,slate] at (-1.2,0.8) {BEFORE:};
  \draw[ax] (0,0.8) -- (5.2,0.8) node[right,font=\scriptsize]{$x$};
  \fill[teal] (1.2,0.8) circle (7pt) node[above=6pt,font=\footnotesize]{$m_A$};
  \draw[vec,teal] (1.2,0.8) -- (2.3,0.8) node[midway,below,font=\footnotesize]{$v_{A1x}$};
  \fill[amber] (3.6,0.8) circle (9pt) node[above=8pt,font=\footnotesize]{$m_B$};
  \draw[vec,amber] (3.6,0.8) -- (2.8,0.8) node[midway,below,font=\footnotesize]{$v_{B1x}$};
  % After
  \node[font=\footnotesize\bfseries,slate] at (-1.2,-0.6) {AFTER:};
  \draw[ax] (0,-0.6) -- (5.2,-0.6) node[right,font=\scriptsize]{$x$};
  \fill[teal] (1.2,-0.6) circle (7pt) node[above=6pt,font=\footnotesize]{$m_A$};
  \draw[vec,teal] (1.2,-0.6) -- (0.3,-0.6) node[midway,below,font=\footnotesize]{$v_{A2x}$};
  \fill[amber] (3.6,-0.6) circle (9pt) node[above=8pt,font=\footnotesize]{$m_B$};
  \draw[vec,amber] (3.6,-0.6) -- (4.7,-0.6) node[midway,below,font=\footnotesize]{$v_{B2x}$};
\end{tikzpicture}
\figcap{One-dimensional elastic collision along the $x$-axis.}
\end{pic}

\step{Step 1: Write kinetic energy and momentum conservation.}
\begin{align}
  \frac{1}{2}m_A v_{A1x}^2 + \frac{1}{2}m_B v_{B1x}^2 &= \frac{1}{2}m_A v_{A2x}^2 + \frac{1}{2}m_B v_{B2x}^2 &&\why{kinetic energy conserved}\label{eq:el_K}\\
  m_A v_{A1x} + m_B v_{B1x} &= m_A v_{A2x} + m_B v_{B2x} &&\why{$x$-momentum conserved}\label{eq:el_P}
\end{align}

\step{Step 2: Group terms by mass $m_A$ on the left and $m_B$ on the right.}
From Eq.\ \eqref{eq:el_K}, multiply by $2$ and factor using $a^2 - b^2 = (a-b)(a+b)$:
\begin{align}
  m_A(v_{A1x} - v_{A2x})(v_{A1x} + v_{A2x}) &= m_B(v_{B2x} - v_{B1x})(v_{B2x} + v_{B1x})\label{eq:el_fact1}\\
  m_A(v_{A1x} - v_{A2x}) &= m_B(v_{B2x} - v_{B1x}) &&\why{from Eq.\ \eqref{eq:el_P}}\label{eq:el_fact2}
\end{align}

\step{Step 3: Divide Eq.\ \eqref{eq:el_fact1} by Eq.\ \eqref{eq:el_fact2} (for a non-miss collision $v_{A1x}\neq v_{A2x}$).}
\begin{equation}\label{eq:rel_vel_1D}
  v_{A1x} + v_{A2x} = v_{B2x} + v_{B1x} \implies \eqbox{v_{B2x} - v_{A2x} = -(v_{B1x} - v_{A1x})}
\end{equation}

\step{Step 4: Solve for $v_{A2x}$ and $v_{B2x}$ when target $B$ is initially at rest ($v_{B1x} = 0$).}
With $v_{B1x} = 0$, Eq.\ \eqref{eq:rel_vel_1D} gives $v_{B2x} = v_{A1x} + v_{A2x}$. Substitute this into $m_A v_{A1x} = m_A v_{A2x} + m_B v_{B2x}$:
\begin{align}
  m_A v_{A1x} &= m_A v_{A2x} + m_B(v_{A1x} + v_{A2x}) = (m_A + m_B)v_{A2x} + m_B v_{A1x}\\
  v_{A2x} &= \frac{m_A - m_B}{m_A + m_B}v_{A1x}, \qquad v_{B2x} = v_{A1x}\!\left(1 + \frac{m_A - m_B}{m_A + m_B}\right) = \frac{2m_A}{m_A + m_B}v_{A1x}
\end{align}

\begin{result}[1D Elastic Collision with Stationary Target ($v_{B1x}=0$)]
\begin{equation}\label{eq:elastic_rest}
  v_{A2x} = \frac{m_A - m_B}{m_A + m_B}v_{A1x}, \qquad v_{B2x} = \frac{2m_A}{m_A + m_B}v_{A1x}
\end{equation}
If $v_{B1x}\neq 0$, the general solution is:
\begin{equation}\label{eq:elastic_gen}
  v_{A2x} = \frac{m_A - m_B}{m_A + m_B}v_{A1x} + \frac{2m_B}{m_A + m_B}v_{B1x}, \qquad
  v_{B2x} = \frac{2m_A}{m_A + m_B}v_{A1x} + \frac{m_B - m_A}{m_A + m_B}v_{B1x}
\end{equation}
\end{result}

\section{Two-Dimensional Glancing Elastic Collisions\bk{8.4, 6.3}}

\begin{goal}
Derive the general scattering equation for a mass $M_1$ striking a stationary target $M_2$ at a glancing angle $\theta_1$, and prove that $\theta_1 + \theta_2 = \frac{\pi}{2}$ when $M_1 = M_2$.
\end{goal}

\begin{pic}
\begin{tikzpicture}[scale=0.95]
  \draw[ax] (-2.8,0) -- (3.8,0) node[right,font=\small]{$x$};
  \draw[ax] (0,-1.8) -- (0,2.0) node[above,font=\small]{$y$};
  \fill[teal] (-2.2,0) circle (5.5pt) node[above=5pt,font=\small]{$M_1$};
  \draw[vec,teal] (-2.2,0) -- (-0.9,0) node[midway,below,font=\small]{$\vect{V}_1$};
  \fill[amber] (0,0) circle (5.5pt) node[below left,font=\small]{$M_2$};
  \draw[vec,teal] (0,0) -- (2.2,1.5) node[above right,font=\small]{$\vect{V}_1'$};
  \fill[teal] (1.32,0.9) circle (5.5pt);
  \draw[vec,amber] (0,0) -- (2.4,-1.3) node[below right,font=\small]{$\vect{V}_2'$};
  \fill[amber] (1.44,-0.78) circle (5.5pt);
  \draw[slate] (0.75,0) arc (0:34.3:0.75) node[midway,right,font=\footnotesize]{$\theta_1$};
  \draw[slate] (0.85,0) arc (0:-28.4:0.85) node[midway,right,font=\footnotesize]{$\theta_2$};
\end{tikzpicture}
\figcap{Glancing collision of incoming particle $M_1$ with stationary target $M_2$ at the origin.}
\end{pic}

\step{Step 1: Write $x$-momentum, $y$-momentum, and kinetic energy conservation (Eqs.\ 6.4--6.6).}
\begin{align}
  M_1 V_1 &= M_1 V_1'\cos\theta_1 + M_2 V_2'\cos\theta_2 &&\why{$P_{xi} = P_{xf}$}\label{eq:gl_px}\\
  0 &= M_1 V_1'\sin\theta_1 - M_2 V_2'\sin\theta_2 &&\why{$P_{yi} = P_{yf}$}\label{eq:gl_py}\\
  \frac{1}{2}M_1 V_1^2 &= \frac{1}{2}M_1 V_1'^2 + \frac{1}{2}M_2 V_2'^2 &&\why{$T_i = T_f$}\label{eq:gl_E}
\end{align}

\step{Step 2: Eliminate $\theta_2$ using $\cos^2\theta_2 + \sin^2\theta_2 = 1$.}
Isolate the $M_2 V_2'$ terms in Eqs.\ \eqref{eq:gl_px} and \eqref{eq:gl_py}, square both, and add:
\begin{align}
  (M_1 V_1 - M_1 V_1'\cos\theta_1)^2 + (M_1 V_1'\sin\theta_1)^2 &= M_2^2 V_2'^2(\cos^2\theta_2 + \sin^2\theta_2)\nonumber\\
  M_1^2 V_1^2 - 2M_1^2 V_1 V_1'\cos\theta_1 + M_1^2 V_1'^2 &= M_2^2 V_2'^2\label{eq:gl_sq}
\end{align}

\step{Step 3: Eliminate $V_2'^2$ using Eq.\ \eqref{eq:gl_E}.}
From Eq.\ \eqref{eq:gl_E}, $M_2^2 V_2'^2 = M_1 M_2(V_1^2 - V_1'^2)$. Substitute into Eq.\ \eqref{eq:gl_sq} and divide by $M_1$:
\begin{equation}\label{eq:gl_quad}
  (M_1 + M_2)V_1'^2 - (2M_1 V_1\cos\theta_1)V_1' + (M_1 - M_2)V_1^2 = 0
\end{equation}
Apply the quadratic formula for $V_1'/V_1$ (Eq.\ 6.11):
\begin{equation}\label{eq:gl_ratio}
  \eqbox{\frac{V_1'}{V_1} = \frac{M_1}{M_1 + M_2}\!\left[\cos\theta_1 \pm \sqrt{\cos^2\theta_1 - 1 + \left(\frac{M_2}{M_1}\right)^2}\right]}
\end{equation}

\step{Step 4: Analyze special cases.}
\begin{itemize}
  \item \textbf{Equal masses ($M_1 = M_2$):} Eq.\ \eqref{eq:gl_ratio} gives $V_1' = V_1\cos\theta_1$. Then Eq.\ \eqref{eq:gl_E} gives $V_2'^2 = V_1^2(1-\cos^2\theta_1) = V_1^2\sin^2\theta_1 \implies V_2' = V_1\sin\theta_1$. Substituting into Eq.\ \eqref{eq:gl_py} yields $\cos\theta_1\sin\theta_1 = \sin\theta_1\sin\theta_2 \implies \cos\theta_1 = \sin\theta_2 = \cos\!\left(\frac{\pi}{2}-\theta_2\right)$, so:
  \begin{equation}\label{eq:90deg_law}
    \eqbox{\theta_1 + \theta_2 = \frac{\pi}{2} = 90^\circ \qquad (M_1 = M_2)}
  \end{equation}
  \item \textbf{Heavy projectile ($M_1 > M_2$):} For $V_1'$ to be real, the radicand in Eq.\ \eqref{eq:gl_ratio} must be non-negative: $\cos^2\theta_1 \ge 1 - \frac{M_2^2}{M_1^2}$, or $\sin\theta_1 \le \frac{M_2}{M_1}$. Thus the maximum scattering angle is:
  \begin{equation}\label{eq:theta_max}
    \eqbox{\theta_{\max} = \cos^{-1}\!\left(\sqrt{1 - \frac{M_2^2}{M_1^2}}\right) = \sin^{-1}\!\left(\frac{M_2}{M_1}\right) < \frac{\pi}{2}}
  \end{equation}
\end{itemize}

\section{Q-Value and Coefficient of Restitution\bk{6.4}}

\begin{goal}
Define the reaction energy $Q$ and Newton's coefficient of restitution $e$.
\end{goal}

\begin{equation}\label{eq:Q_and_e}
  \eqbox{Q = T_f - T_i}, \qquad \eqbox{e = \frac{|V_2' - V_1'|}{|V_2 - V_1|}} \qquad
  \begin{cases}
    e = 1, \ Q = 0 & \text{(perfectly elastic)}\\
    0 < e < 1, \ Q < 0 & \text{(inelastic / endoergic)}\\
    e = 0, \ Q < 0 & \text{(completely inelastic)}
  \end{cases}
\end{equation}

\section{Center of Mass and Reduced Mass\bk{8.5, 6.6, 6.7}}

\begin{goal}
Derive the center-of-mass equation of motion and reduce the two-body problem to an equivalent one-body problem with reduced mass $\mu$.
\end{goal}

\step{Step 1: Center of mass position, velocity, and acceleration.}
For a system of total mass $M = \sum_i m_i = \int\dd m$:
\begin{align}
  \vect{r}_{\text{cm}} &= \frac{\sum_i m_i\vect{r}_i}{M} = \frac{1}{M}\int\vect{r}\dd m &&\why{position of CM}\label{eq:r_cm}\\
  M\vect{v}_{\text{cm}} &= \sum_i m_i\vect{v}_i = \vect{P}_{\text{tot}} &&\why{differentiate once with respect to $t$}\label{eq:v_cm}\\
  M\vect{a}_{\text{cm}} &= \sum_i m_i\vect{a}_i = \frac{\diff\vect{P}_{\text{tot}}}{\diff t} = \sum\vect{F}_{\text{ext}} &&\why{differentiate twice with respect to $t$}\label{eq:a_cm}
\end{align}

\step{Step 2: Relative coordinate $\vect{r}$ and reduced mass $\mu$ (Sec.\ 6.7).}
Consider an isolated two-body system ($m_1$ at $\vect{r}_1$ and $m_2$ at $\vect{r}_2$). Let $\vect{F}$ be the force on $m_2$ due to $m_1$. By Newton's third law, the force on $m_1$ due to $m_2$ is $-\vect{F}$:
\begin{align}
  m_1\ddot{\vect{r}}_1 &= -\vect{F} \implies \ddot{\vect{r}}_1 = -\frac{\vect{F}}{m_1}, \qquad
  m_2\ddot{\vect{r}}_2 = +\vect{F} \implies \ddot{\vect{r}}_2 = +\frac{\vect{F}}{m_2}
\end{align}
Define the relative position vector $\vect{r} = \vect{r}_2 - \vect{r}_1$ and subtract the accelerations:
\begin{align}
  \ddot{\vect{r}} &= \ddot{\vect{r}}_2 - \ddot{\vect{r}}_1 = \left(\frac{1}{m_1} + \frac{1}{m_2}\right)\vect{F} = \frac{\vect{F}}{\mu}
\end{align}
\begin{result}[Reduced Mass and Relative Equation of Motion]
\begin{equation}\label{eq:reduced_mass}
  \mu\ddot{\vect{r}} = \vect{F}, \qquad \text{where}\quad \frac{1}{\mu} = \frac{1}{m_1} + \frac{1}{m_2} \iff \mu = \frac{m_1 m_2}{m_1 + m_2}
\end{equation}
\end{result}

\section{Rocket Propulsion\bk{8.6, 6.2}}

\begin{goal}
Derive the rocket thrust equation and the Tsiolkovsky velocity equation in free space and in a uniform gravitational field.
\end{goal}

\begin{pic}
\begin{tikzpicture}[scale=0.95]
  % Time t
  \node[font=\footnotesize\bfseries,slate] at (-1.5,0.9) {At time $t$:};
  \draw[thick,teal,fill=tealsoft] (0,0.6) -- (2.2,0.6) -- (2.7,0.9) -- (2.2,1.2) -- (0,1.2) -- cycle;
  \node[font=\footnotesize] at (1.2,0.9) {Mass $m$};
  \draw[vec,teal] (2.8,0.9) -- (3.8,0.9) node[right,font=\footnotesize]{$v$};
  % Time t + dt
  \node[font=\footnotesize\bfseries,slate] at (-1.5,-0.5) {At $t+\diff t$:};
  \draw[thick,amber,fill=ambersoft,dashed] (-0.2,-0.75) rectangle (0.6,-0.25);
  \node[font=\scriptsize] at (0.2,-1.0) {$-\diff m = \diff M_F$};
  \draw[vec,amber] (-0.2,-0.5) -- (-1.0,-0.5) node[above,font=\scriptsize]{$v_{\text{fuel}} = v - v_{\text{ex}}$};
  \draw[thick,teal,fill=tealsoft] (1.1,-0.8) -- (3.1,-0.8) -- (3.6,-0.5) -- (3.1,-0.2) -- (1.1,-0.2) -- cycle;
  \node[font=\scriptsize] at (2.1,-0.5) {$m+\diff m$};
  \draw[vec,teal] (3.7,-0.5) -- (4.9,-0.5) node[right,font=\footnotesize]{$v+\diff v$};
\end{tikzpicture}
\figcap{Conservation of momentum for a rocket ejecting burned fuel mass $-\diff m = \diff M_F > 0$ at relative exhaust speed $v_{\text{ex}}$.}
\end{pic}

\step{Step 1: Write momentum at time $t$ and at time $t + \diff t$.}
At time $t$, the rocket has mass $m$ and velocity $v$, so $P_1 = mv$. During $\diff t$, a positive mass of burned fuel $-\diff m = \diff M_F > 0$ (since $\diff m < 0$ is the change in rocket mass) is ejected backward with speed $v_{\text{ex}}$ relative to the rocket. The fuel velocity in the inertial frame is $v_{\text{fuel}} = v - v_{\text{ex}}$. At $t + \diff t$:
\begin{equation}
  P_2 = (m + \diff m)(v + \diff v) + (-\diff m)(v - v_{\text{ex}})
\end{equation}

\step{Step 2: Apply momentum conservation ($P_1 = P_2$) in gravity-free space.}
\begin{align}
  mv &= mv + m\dd v + v\dd m + \cancel{\diff m\dd v} - v\dd m + v_{\text{ex}}\dd m &&\why{neglect 2nd-order differential $\diff m\dd v$}\\
  m\dd v &= -v_{\text{ex}}\dd m = +v_{\text{ex}}\dd M_F &&\why{divide by $\diff t$}\label{eq:rocket_diff}
\end{align}
\begin{equation}\label{eq:rocket_thrust}
  \eqbox{F_{\text{thrust}} = m\frac{\diff v}{\diff t} = -v_{\text{ex}}\frac{\diff m}{\diff t}, \qquad a = \frac{\diff v}{\diff t} = -\frac{v_{\text{ex}}}{m}\frac{\diff m}{\diff t}}
\end{equation}

\step{Step 3: Integrate Eq.\ \eqref{eq:rocket_diff} from $(m_0, v_0)$ to $(m, v)$.}
\begin{align}
  \int_{v_0}^v \diff v' &= -v_{\text{ex}}\int_{m_0}^m \frac{\diff m'}{m'} = -v_{\text{ex}}\ln\!\left(\frac{m}{m_0}\right) = v_{\text{ex}}\ln\!\left(\frac{m_0}{m}\right)\label{eq:rocket_vel}
\end{align}
If a uniform downward gravitational force $-mg$ also acts on the rocket ($F_{\text{ext}} = -mg$), then $m\frac{\diff v}{\diff t} = -v_{\text{ex}}\frac{\diff m}{\diff t} - mg$, which integrates to:
\begin{equation}\label{eq:rocket_grav}
  \eqbox{v - v_0 = v_{\text{ex}}\ln\!\left(\frac{m_0}{m}\right) - gt}
\end{equation}

% ============================================================================
% SOLVED BOOK EXAMPLES — MODULE 1
% ============================================================================
\section{Solved Book Examples (Module 1)}

\begin{example}{8.1}{Momentum versus kinetic energy\bk{p.\ 245}}
\textbf{Problem:} Two iceboats of masses $m$ and $2m$ start from rest on a frictionless lake. The wind exerts the same constant horizontal force $\vect{F}$ on both over the same distance $s$. Which boat crosses the finish line with greater momentum, and by what factor?
\tcblower
\textbf{Solution:} Both boats have the same kinetic energy at the finish line because $K = W_{\text{tot}} = Fs$ is identical for both. Using Eq.\ \eqref{eq:K_p_rel}, $p = \sqrt{2mK}$:
\begin{equation*}
  \frac{p_{2m}}{p_m} = \frac{\sqrt{2(2m)K}}{\sqrt{2mK}} = \eqbox{\sqrt{2}}
\end{equation*}
The boat of mass $2m$ takes $\sqrt{2}$ times longer to cover distance $s$, receives $\sqrt{2}$ times the impulse $\vect{J}=\vect{F}\Delta t$, and has $\sqrt{2}$ times the momentum.
\end{example}

\begin{example}{8.2}{A ball hits a wall\bk{p.\ 245}}
\textbf{Problem:} A $0.40\text{-kg}$ ball hits a brick wall horizontally to the left at $30\text{ m/s}$ and rebounds horizontally to the right at $20\text{ m/s}$. (a) Find the impulse of the net force. (b) If the contact time is $0.010\text{ s}$, find the average horizontal force.
\tcblower
\textbf{Solution:} Take $+x$ to the right: $v_{1x} = -30\text{ m/s}$ and $v_{2x} = +20\text{ m/s}$.
\begin{align*}
  p_{1x} &= (0.40\text{ kg})(-30\text{ m/s}) = -12\text{ kg}\cdot\text{m/s}, \qquad p_{2x} = (0.40\text{ kg})(+20\text{ m/s}) = +8.0\text{ kg}\cdot\text{m/s}\\
  \text{(a)}\quad J_x &= p_{2x} - p_{1x} = 8.0 - (-12) = \eqbox{+20\text{ N}\cdot\text{s}} \quad(\text{to the right})\\
  \text{(b)}\quad (F_{\text{av}})_x &= \frac{J_x}{\Delta t} = \frac{20\text{ N}\cdot\text{s}}{0.010\text{ s}} = \eqbox{2000\text{ N}}
\end{align*}
\end{example}

\begin{example}{8.3}{Kicking a soccer ball\bk{p.\ 246}}
\textbf{Problem:} A $0.40\text{-kg}$ soccer ball moves left at $20\text{ m/s}$. After a kick lasting $\Delta t = 0.010\text{ s}$, it moves at $45^\circ$ upward and right at $30\text{ m/s}$. Find the impulse components and the magnitude and direction of the average net force.
\tcblower
\textbf{Solution:} Initial components: $v_{1x} = -20\text{ m/s}$, $v_{1y} = 0$. Final components: $v_{2x} = v_{2y} = (30\text{ m/s})\cos 45^\circ = 21.2\text{ m/s}$.
\begin{align*}
  J_x &= m(v_{2x}-v_{1x}) = 0.40[21.2 - (-20)] = \eqbox{16.5\text{ kg}\cdot\text{m/s}}, \qquad (F_{\text{av}})_x = \frac{J_x}{\Delta t} = 1650\text{ N}\\
  J_y &= m(v_{2y}-v_{1y}) = 0.40(21.2 - 0) = \eqbox{8.5\text{ kg}\cdot\text{m/s}}, \qquad (F_{\text{av}})_y = \frac{J_y}{\Delta t} = 850\text{ N}\\
  F_{\text{av}} &= \sqrt{(1650)^2 + (850)^2} = \eqbox{1.9\times 10^3\text{ N}}, \qquad \theta = \arctan\!\left(\frac{850}{1650}\right) = \eqbox{27^\circ}
\end{align*}
\end{example}

\begin{example}{8.4}{Recoil of a rifle\bk{p.\ 249}}
\textbf{Problem:} A rifle of mass $m_R = 3.00\text{ kg}$ fires a bullet of mass $m_B = 5.00\text{ g} = 0.00500\text{ kg}$ at $v_{Bx} = 300\text{ m/s}$. Find the recoil velocity $v_{Rx}$ of the rifle and the final momentum and kinetic energy of each.
\tcblower
\textbf{Solution:} Total initial momentum is $P_x = 0 = m_B v_{Bx} + m_R v_{Rx}$:
\begin{align*}
  v_{Rx} &= -\frac{m_B}{m_R}v_{Bx} = -\frac{0.00500}{3.00}(300\text{ m/s}) = \eqbox{-0.500\text{ m/s}}\\
  p_{Bx} &= (0.00500)(300) = \eqbox{+1.50\text{ kg}\cdot\text{m/s}}, \qquad p_{Rx} = (3.00)(-0.500) = \eqbox{-1.50\text{ kg}\cdot\text{m/s}}\\
  K_B &= \frac{1}{2}(0.00500)(300)^2 = \eqbox{225\text{ J}}, \qquad K_R = \frac{1}{2}(3.00)(-0.500)^2 = \eqbox{0.375\text{ J}}
\end{align*}
\end{example}

\begin{example}{8.5}{Collision along a straight line\bk{p.\ 249}}
\textbf{Problem:} Glider $A$ ($m_A = 0.50\text{ kg}$, $v_{A1x} = +2.0\text{ m/s}$) and glider $B$ ($m_B = 0.30\text{ kg}$, $v_{B1x} = -2.0\text{ m/s}$) collide on a frictionless air track. After collision, $v_{B2x} = +2.0\text{ m/s}$. Find $v_{A2x}$ and the changes in momentum and velocity.
\tcblower
\textbf{Solution:} Total $x$-momentum before collision: $P_x = (0.50)(2.0) + (0.30)(-2.0) = +0.40\text{ kg}\cdot\text{m/s}$.
\begin{align*}
  v_{A2x} &= \frac{P_x - m_B v_{B2x}}{m_A} = \frac{0.40 - (0.30)(2.0)}{0.50} = \eqbox{-0.40\text{ m/s}}\\
  \Delta p_{Ax} &= 0.50(-0.40 - 2.0) = -1.2\text{ kg}\cdot\text{m/s}, \qquad \Delta p_{Bx} = 0.30[2.0 - (-2.0)] = +1.2\text{ kg}\cdot\text{m/s}\\
  \Delta v_{Ax} &= -2.4\text{ m/s}, \qquad \Delta v_{Bx} = +4.0\text{ m/s}
\end{align*}
\end{example}

\begin{example}{8.6}{Collision in a horizontal plane\bk{p.\ 250}}
\textbf{Problem:} Robot $A$ ($m_A = 20\text{ kg}$) moves at $v_{A1x} = 2.0\text{ m/s}$ along $+x$ and hits stationary robot $B$ ($m_B = 12\text{ kg}$). After impact, $A$ moves at $v_{A2} = 1.0\text{ m/s}$ at $\alpha = 30^\circ$ above $+x$. Find the final velocity $\vect{v}_{B2}$ of robot $B$.
\tcblower
\textbf{Solution:} Apply $x$- and $y$-momentum conservation:
\begin{align*}
  v_{B2x} &= \frac{m_A v_{A1x} - m_A v_{A2}\cos 30^\circ}{m_B} = \frac{20(2.0) - 20(1.0)\cos 30^\circ}{12} = 1.89\text{ m/s}\\
  v_{B2y} &= \frac{0 - m_A v_{A2}\sin 30^\circ}{m_B} = \frac{-20(1.0)\sin 30^\circ}{12} = -0.83\text{ m/s}\\
  v_{B2} &= \sqrt{(1.89)^2 + (-0.83)^2} = \eqbox{2.1\text{ m/s}}, \qquad \beta = \arctan\!\left(\frac{-0.83}{1.89}\right) = \eqbox{-24^\circ}
\end{align*}
\end{example}

\begin{example}{8.7}{A completely inelastic collision\bk{p.\ 252}}
\textbf{Problem:} Repeat Example 8.5 ($m_A = 0.50\text{ kg}, v_{A1x} = 2.0\text{ m/s}$; $m_B = 0.30\text{ kg}, v_{B1x} = -2.0\text{ m/s}$) when the gliders stick together. Find $v_{2x}$, $K_1$, and $K_2$.
\tcblower
\textbf{Solution:}
\begin{align*}
  v_{2x} &= \frac{m_A v_{A1x} + m_B v_{B1x}}{m_A + m_B} = \frac{0.40\text{ kg}\cdot\text{m/s}}{0.80\text{ kg}} = \eqbox{0.50\text{ m/s}}\\
  K_1 &= \frac{1}{2}(0.50)(2.0)^2 + \frac{1}{2}(0.30)(-2.0)^2 = 1.0 + 0.60 = \eqbox{1.6\text{ J}}\\
  K_2 &= \frac{1}{2}(0.50 + 0.30)(0.50)^2 = \eqbox{0.10\text{ J}} \quad\left(\frac{15}{16}\text{ of }K_1\text{ is converted to other energy forms}\right)
\end{align*}
\end{example}

\begin{example}{8.8}{The ballistic pendulum\bk{p.\ 253}}
\textbf{Problem:} A bullet of mass $m_B = 5.00\text{ g}$ embeds into a suspended wooden block of mass $m_W = 2.00\text{ kg}$, which swings up to maximum height $y = 3.00\text{ cm} = 0.0300\text{ m}$. Derive the formula for initial bullet speed $v_1$ and evaluate it.
\tcblower
\textbf{Solution:} Stage 1 (collision, momentum conserved): $m_B v_1 = (m_B + m_W)v_2 \implies v_1 = \frac{m_B + m_W}{m_B}v_2$.
Stage 2 (pendulum swing, mechanical energy conserved): $\frac{1}{2}(m_B + m_W)v_2^2 = (m_B + m_W)gy \implies v_2 = \sqrt{2gy}$.
\begin{align*}
  v_1 &= \eqbox{\frac{m_B + m_W}{m_B}\sqrt{2gy}} = \frac{0.00500 + 2.00}{0.00500}\sqrt{2(9.80)(0.0300)} = \eqbox{307\text{ m/s}}, \qquad v_2 = 0.767\text{ m/s}
\end{align*}
\end{example}

\begin{example}{8.9}{An automobile collision\bk{p.\ 254}}
\textbf{Problem:} A $1000\text{-kg}$ car traveling north ($+y$) at $15\text{ m/s}$ collides with a $2000\text{-kg}$ truck traveling east ($+x$) at $10\text{ m/s}$. They lock together ($M = 3000\text{ kg}$). Find the velocity $\vect{V}$ of the wreckage just after impact.
\tcblower
\textbf{Solution:}
\begin{align*}
  P_x &= (2000\text{ kg})(10\text{ m/s}) = 2.0\times 10^4\text{ kg}\cdot\text{m/s}, \qquad P_y = (1000\text{ kg})(15\text{ m/s}) = 1.5\times 10^4\text{ kg}\cdot\text{m/s}\\
  P &= \sqrt{P_x^2 + P_y^2} = 2.5\times 10^4\text{ kg}\cdot\text{m/s}, \qquad \tan\theta = \frac{P_y}{P_x} = 0.75 \implies \theta = \eqbox{37^\circ\text{ north of east}}\\
  V &= \frac{P}{M} = \frac{2.5\times 10^4}{3000} = \eqbox{8.3\text{ m/s}}
\end{align*}
\end{example}

\begin{example}{8.10}{An elastic straight-line collision\bk{p.\ 257}}
\textbf{Problem:} Repeat Example 8.5 ($m_A = 0.50\text{ kg}, v_{A1x} = 2.0\text{ m/s}$; $m_B = 0.30\text{ kg}, v_{B1x} = -2.0\text{ m/s}$) for an elastic collision. Find $v_{A2x}$ and $v_{B2x}$.
\tcblower
\textbf{Solution:} Relative velocity Eq.\ \eqref{eq:rel_vel_1D} and momentum conservation give:
\begin{align*}
  v_{B2x} - v_{A2x} &= -(v_{B1x} - v_{A1x}) = -(-2.0 - 2.0) = 4.0\text{ m/s} \implies v_{B2x} = v_{A2x} + 4.0\\
  0.50v_{A2x} + 0.30v_{B2x} &= 0.40\text{ m/s} \implies 0.50v_{A2x} + 0.30(v_{A2x} + 4.0) = 0.40\\
  0.80v_{A2x} &= -0.80 \implies \eqbox{v_{A2x} = -1.0\text{ m/s}}, \qquad \eqbox{v_{B2x} = +3.0\text{ m/s}}
\end{align*}
\end{example}

\begin{example}{8.11}{Moderating fission neutrons in a nuclear reactor\bk{p.\ 257}}
\textbf{Problem:} A neutron ($m_n = 1.0\text{ u}$, $v_{n1x} = 2.6\times 10^7\text{ m/s}$) makes a head-on elastic collision with a carbon nucleus ($m_C = 12\text{ u}$) initially at rest. Find the final velocities $v_{n2x}$ and $v_{C2x}$.
\tcblower
\textbf{Solution:} Apply Eq.\ \eqref{eq:elastic_rest}:
\begin{align*}
  v_{n2x} &= \frac{m_n - m_C}{m_n + m_C}v_{n1x} = \frac{1.0 - 12}{1.0 + 12}(2.6\times 10^7\text{ m/s}) = -\frac{11}{13}(2.6\times 10^7) = \eqbox{-2.2\times 10^7\text{ m/s}}\\
  v_{C2x} &= \frac{2m_n}{m_n + m_C}v_{n1x} = \frac{2(1.0)}{13}(2.6\times 10^7\text{ m/s}) = \eqbox{+0.4\times 10^7\text{ m/s}}
\end{align*}
\end{example}

\begin{example}{8.12}{A two-dimensional elastic collision\bk{p.\ 258}}
\textbf{Problem:} Puck $A$ ($m_A = 0.500\text{ kg}$, $v_{A1} = 4.00\text{ m/s}$ along $+x$) collides elastically with stationary puck $B$ ($m_B = 0.300\text{ kg}$). After collision, $v_{A2} = 2.00\text{ m/s}$ at angle $\alpha$ above $+x$, and $B$ moves at speed $v_{B2}$ at angle $\beta$ below $+x$. Find $v_{B2}$, $\alpha$, and $\beta$.
\tcblower
\textbf{Solution:} From energy conservation $\frac{1}{2}m_A v_{A1}^2 = \frac{1}{2}m_A v_{A2}^2 + \frac{1}{2}m_B v_{B2}^2$:
\begin{equation*}
  v_{B2} = \sqrt{\frac{m_A(v_{A1}^2 - v_{A2}^2)}{m_B}} = \sqrt{\frac{0.500(16.0 - 4.00)}{0.300}} = \sqrt{20.0} = \eqbox{4.47\text{ m/s}}
\end{equation*}
Momentum conservation along $x$ and $y$:
\begin{align*}
  (0.300)(4.47)\cos\beta &= (0.500)(4.00) - (0.500)(2.00)\cos\alpha = 2.00 - 1.00\cos\alpha\\
  (0.300)(4.47)\sin\beta &= (0.500)(2.00)\sin\alpha = 1.00\sin\alpha
\end{align*}
Square and add both equations using $\cos^2\beta + \sin^2\beta = 1$:
\begin{align*}
  (1.342)^2 &= 1.80 = (2.00 - \cos\alpha)^2 + \sin^2\alpha = 4.00 - 4.00\cos\alpha + 1.00\\
  \cos\alpha &= \frac{3.20}{4.00} = 0.800 \implies \eqbox{\alpha = 36.9^\circ}, \qquad \sin\beta = \frac{1.00\sin 36.9^\circ}{1.342} = 0.447 \implies \eqbox{\beta = 26.6^\circ}
\end{align*}
\end{example}

\begin{example}{8.13}{Center of mass of a water molecule\bk{p.\ 259}}
\textbf{Problem:} An $\text{H}_2\text{O}$ molecule has oxygen ($16.0\text{ u}$) at the origin and two hydrogen atoms ($1.0\text{ u}$ each) at distance $d = 9.57\times 10^{-11}\text{ m}$ at angles $\pm 105^\circ/2 = \pm 52.5^\circ$ from the $+x$-axis. Find $(x_{\text{cm}}, y_{\text{cm}})$.
\tcblower
\textbf{Solution:} By symmetry, $y_{\text{cm}} = 0$. Along the $x$-axis:
\begin{equation*}
  x_{\text{cm}} = \frac{1.0d\cos 52.5^\circ + 1.0d\cos 52.5^\circ + 16.0(0)}{1.0 + 1.0 + 16.0} = 0.068d = \eqbox{6.5\times 10^{-12}\text{ m}}
\end{equation*}
\end{example}

\begin{example}{8.14}{A tug-of-war on the ice\bk{p.\ 260}}
\textbf{Problem:} James ($m_J = 90.0\text{ kg}$ at $x_1 = -10.0\text{ m}$) and Ramon ($m_R = 60.0\text{ kg}$ at $x_2 = +10.0\text{ m}$) pull on a light rope on frictionless ice. A mug is at $x=0$. When James moves $6.0\text{ m}$ toward the mug (to $x_1' = -4.0\text{ m}$), how far and in what direction has Ramon moved?
\tcblower
\textbf{Solution:} No external horizontal force acts, so $x_{\text{cm}}$ stays fixed:
\begin{align*}
  x_{\text{cm}} &= \frac{90.0(-10.0) + 60.0(+10.0)}{90.0 + 60.0} = -2.0\text{ m}\\
  -2.0\text{ m} &= \frac{90.0(-4.0\text{ m}) + 60.0 x_2'}{150.0} \implies x_2' = +1.0\text{ m}
\end{align*}
Ramon has moved from $+10.0\text{ m}$ to $+1.0\text{ m}$, which is $\eqbox{9.0\text{ m}\text{ toward the mug}}$.
\end{example}

\begin{example}{8.15}{Acceleration of a rocket\bk{p.\ 264}}
\textbf{Problem:} A rocket in outer space ejects $\frac{1}{120}$ of its initial mass $m_0$ in the first $1.0\text{ s}$ at constant rate with relative exhaust speed $v_{\text{ex}} = 2400\text{ m/s}$. Find the initial acceleration.
\tcblower
\textbf{Solution:} Here $\frac{\diff m}{\diff t} = -\frac{m_0}{120\text{ s}}$. From Eq.\ \eqref{eq:rocket_thrust}:
\begin{equation*}
  a = -\frac{v_{\text{ex}}}{m_0}\frac{\diff m}{\diff t} = -\frac{2400\text{ m/s}}{m_0}\!\left(-\frac{m_0}{120\text{ s}}\right) = \eqbox{20\text{ m/s}^2}
\end{equation*}
\end{example}

\begin{example}{8.16}{Speed of a rocket\bk{p.\ 265}}
\textbf{Problem:} If the rocket in Example 8.15 starts from rest ($v_0=0$) and consumes $\frac{3}{4}$ of its initial mass as fuel ($m = m_0/4$), find its final speed.
\tcblower
\textbf{Solution:} From Eq.\ \eqref{eq:rocket_vel}:
\begin{equation*}
  v = v_0 + v_{\text{ex}}\ln\!\left(\frac{m_0}{m}\right) = 0 + (2400\text{ m/s})\ln 4 = \eqbox{3327\text{ m/s}}
\end{equation*}
\end{example}

\begin{example}{Ch.\ 8}{Bridging Problem: One Collision After Another\bk{p.\ 267, 277}}
\textbf{Problem:} Sphere $A$ ($0.600\text{ kg}$, $v_{A1x} = 4.00\text{ m/s}$) hits sphere $B$ ($1.80\text{ kg}$, $v_{B1x} = 2.00\text{ m/s}$). After collision 1, $B$ moves at $v_{B2x} = 3.00\text{ m/s}$ along $+x$. Then $B$ strikes stationary sphere $C$ ($1.20\text{ kg}$) off-center; after collision 2, $B$ moves at $v_{B3} = 2.00\text{ m/s}$ at $+19.0^\circ$ to $+x$. Find: (a) $\vect{v}_{A2}$, (b) whether collision 1 is elastic, (c) $\vect{v}_{C3}$, (d) impulse $\vect{J}_{\text{on }B}$ in collision 2, (e) whether collision 2 is elastic, and (f) $\vect{v}_{\text{cm}}$ of the three-sphere system.
\tcblower
\textbf{Solution:}
\begin{itemize}
  \item \textbf{(a)} $P_{1x} = 0.600(4.00) + 1.80(2.00) = 6.00\text{ kg}\cdot\text{m/s} = 0.600v_{A2x} + 1.80(3.00) \implies \eqbox{v_{A2x} = +1.00\text{ m/s}\text{ (right)}}$.
  \item \textbf{(b)} $K_1 = \frac{1}{2}(0.600)(4.00)^2 + \frac{1}{2}(1.80)(2.00)^2 = 8.40\text{ J}$. $K_2 = \frac{1}{2}(0.600)(1.00)^2 + \frac{1}{2}(1.80)(3.00)^2 = 8.40\text{ J} \implies \eqbox{\text{Elastic}}$.
  \item \textbf{(c)} In collision 2, $p_{B2x} = 1.80(3.00) = 5.40\text{ kg}\cdot\text{m/s}$, $p_{B2y} = 0$. After collision 2:
  \begin{align*}
    v_{C3x} &= \frac{5.40 - 1.80(2.00)\cos 19.0^\circ}{1.20} = 1.664\text{ m/s}, \qquad v_{C3y} = \frac{0 - 1.80(2.00)\sin 19.0^\circ}{1.20} = -0.9767\text{ m/s}\\
    v_{C3} &= \sqrt{(1.664)^2 + (-0.9767)^2} = \eqbox{1.93\text{ m/s}\text{ at }\theta_C = -30.4^\circ}
  \end{align*}
  \item \textbf{(d)} $\vect{J}_{\text{on }B} = \Delta\vect{p}_B = -\Delta\vect{p}_C = -m_C\vect{v}_{C3}$, so $J = 1.20(1.93) = \eqbox{2.31\text{ kg}\cdot\text{m/s}\text{ at }180^\circ - 30.4^\circ = 149.6^\circ}$.
  \item \textbf{(e)} $K_{\text{before 2}} = \frac{1}{2}(1.80)(3.00)^2 = 8.10\text{ J}$; $K_{\text{after 2}} = \frac{1}{2}(1.80)(2.00)^2 + \frac{1}{2}(1.20)(1.93)^2 = 5.83\text{ J} < 8.10\text{ J} \implies \eqbox{\text{Inelastic}}$.
  \item \textbf{(f)} $v_{\text{cm-}x} = \frac{P_{\text{tot},x}}{m_A+m_B+m_C} = \frac{6.00\text{ kg}\cdot\text{m/s}}{0.600+1.80+1.20\text{ kg}} = \eqbox{1.67\text{ m/s}\text{ in the }+x\text{-direction}}$.
\end{itemize}
\end{example}

\begin{example}{6.1}{Braking a rocket in interplanetary space\bk{Analytical Mechanics p.\ 152}}
\textbf{Problem:} A rocket of total mass $m_0$ coasts at speed $V_0$ and fires its engine forward (exhaust speed $U$ relative to rocket) to slow down to $V_0/2$. Find the mass of fuel burned.
\tcblower
\textbf{Solution:} Because fuel is ejected forward to decelerate the rocket, $v_{\text{fuel}} = V + U$. Momentum conservation $M\dd V + (V+U)\dd M_F - V\dd M_F = 0$ with $\dd M_F = -\dd M$ gives $\dd V = U\frac{\diff M}{M}$. Integrating from $(m_0, V_0)$ to $(m_f, V_0/2)$:
\begin{equation*}
  \frac{1}{U}\int_{V_0}^{V_0/2}\diff V = \int_{m_0}^{m_f}\frac{\diff M}{M} \implies -\frac{V_0}{2U} = \ln\!\left(\frac{m_f}{m_0}\right) \implies m_f = m_0 e^{-V_0/2U}, \quad \Delta m_{\text{fuel}} = m_0 - m_f = \eqbox{m_0\!\left(1 - e^{-V_0/2U}\right)}
\end{equation*}
\end{example}

\begin{example}{6.2}{Rocket equation from total system momentum\bk{Analytical Mechanics p.\ 152}}
\textbf{Problem:} Derive $M_R\frac{\diff V_R}{\diff t} = -U\frac{\diff M_R}{\diff t}$ by differentiating the integral expression for total momentum at time $\tau$.
\tcblower
\textbf{Solution:} Let burn rate be $\dot{m} = -\frac{\diff M_R}{\diff t} = \frac{\diff m_f}{\diff t}$. At time $\tau$, the momentum of all expelled fuel is $P_f(\tau) = \int_0^\tau \dot{m}(V_R(t)-U)\dd t$ and rocket momentum is $P_R(\tau) = M_R(\tau)V_R(\tau)$. Differentiating $P_{\text{tot}} = P_R + P_f$ with respect to $\tau$:
\begin{equation*}
  \frac{\diff P_{\text{tot}}}{\diff t} = \frac{\diff M_R}{\diff t}V_R + M_R\frac{\diff V_R}{\diff t} + \dot{m}(V_R - U) = 0
\end{equation*}
Since $\frac{\diff M_R}{\diff t} = -\dot{m}$, the first and third terms cancel, leaving $\eqbox{M_R\frac{\diff V_R}{\diff t} = \dot{m}U = -U\frac{\diff M_R}{\diff t}}$.
\end{example}

\begin{example}{6.3}{Spacecraft gravitational slingshot around Mars\bk{Analytical Mechanics p.\ 158}}
\textbf{Problem:} A spacecraft ($M_1 = 2000\text{ kg}$, $V_{\text{sc},i} = -12\text{ km/s}$) makes a 1D elastic slingshot around Mars ($M_2 = 6.4\times 10^{23}\text{ kg}$, $V_{\text{Mars},i} = +23.36\text{ km/s}$). Find (a) the final speed of the spacecraft and (b) the ratio $T_f/T_i$ of its kinetic energies.
\tcblower
\textbf{Solution:} Transform to the rest frame of Mars by subtracting $23.36\text{ km/s}$: $V_2 = 0$ and $V_1 = -12 - 23.36 = -35.36\text{ km/s}$. Since $M_1 \ll M_2$, Eq.\ \eqref{eq:elastic_rest} gives $V_1' = \frac{M_1 - M_2}{M_1 + M_2}V_1 \approx -V_1 = +35.36\text{ km/s}$. Transform back to the Sun frame by adding $23.36\text{ km/s}$:
\begin{equation*}
  \text{(a)}\quad v_{\text{sc}}' = 35.36 + 23.36 = \eqbox{+58.72\text{ km/s}}, \qquad \text{(b)}\quad \frac{T_f}{T_i} = \left(\frac{58.72}{12}\right)^2 = \eqbox{23.9}
\end{equation*}
\end{example}

\begin{example}{6.4}{90-degree deflection rule for equal masses\bk{Analytical Mechanics p.\ 160}}
\textbf{Problem:} Prove that in any non-head-on 2D elastic collision between two bodies of equal mass ($M_1 = M_2$), the sum of the deflection angles is $\theta_1 + \theta_2 = \pi/2$.
\tcblower
\textbf{Solution:} From Eqs.\ \eqref{eq:gl_ratio} and \eqref{eq:gl_E} with $M_1 = M_2$, $V_1' = V_1\cos\theta_1$ and $V_2' = V_1\sin\theta_1$. Substitute into $y$-momentum conservation $V_1'\sin\theta_1 = V_2'\sin\theta_2$:
\begin{equation*}
  V_1\cos\theta_1\sin\theta_1 = V_1\sin\theta_1\sin\theta_2 \implies \cos\theta_1 = \sin\theta_2 = \cos\!\left(\frac{\pi}{2}-\theta_2\right) \implies \eqbox{\theta_1 + \theta_2 = \frac{\pi}{2}}
\end{equation*}
\end{example}

% ============================================================================
% UNIQUE SOLVED PROBLEMS — MODULE 1
% ============================================================================
\section{Unique Solved Problems (Module 1)}

\subsubsection{Part A: Analytical Mechanics Chapter 6 Exercises (with Complete Derivations)}

\begin{enumerate}[label=\textbf{Ex.\ 6.\arabic*},leftmargin=3.5em]
  \item \textbf{(Bullet and Block):} A bullet of mass $m$ and speed $v$ embeds in a stationary block of mass $M$.
  \begin{itemize}
    \item \textbf{(a)} Final speed: $mv = (m+M)V \implies \eqbox{V = \frac{mv}{M+m}}$.
    \item \textbf{(b)} Kinetic energy is not conserved because nonconservative friction forces act during penetration.
    \item \textbf{(c)} Loss of kinetic energy: $\Delta T = \frac{1}{2}(m+M)V^2 - \frac{1}{2}mv^2 = \frac{1}{2}mv^2\!\left(\frac{m}{m+M}-1\right) = \eqbox{-\frac{1}{2}\frac{mM}{m+M}v^2}$.
  \end{itemize}

  \item \textbf{(Rocket Acceleration):} $m_0 = 500\text{ kg}$, burn rate $\dot{m} = 5\text{ kg/s}$, $U = 300\text{ m/s}$.
  \begin{itemize}
    \item Initial acceleration: $a_0 = \frac{\dot{m}U}{m_0} = \frac{5(300)}{500} = \eqbox{3\text{ m/s}^2}$.
    \item After $t = 60\text{ s}$: $m(60) = 500 - 5(60) = 200\text{ kg} \implies a(60) = \frac{5(300)}{200} = \eqbox{7.5\text{ m/s}^2}$.
  \end{itemize}

  \item \textbf{(Snowing Flatcars --- Joe and Bill):} Joe sweeps snow off the side perpendicular to the flatcar motion; Bill lets snow accumulate.
  \begin{itemize}
    \item \textbf{Conceptual \& Mathematical Solution:} Vertically falling snow has zero initial horizontal momentum. When snow lands on the car of mass $M$ and speed $V$, it acquires horizontal speed $V$. When Joe sweeps snow off perpendicular to the track, the snow leaves with horizontal velocity $V$ relative to the ground, so it carries away momentum $\diff m_{\text{snow}}V$ just like the snow remaining on Bill's car! In both cases, $P_x = M_0 V_0 = (M_0 + m_{\text{snow}}(t))V(t)$. Therefore, \textbf{both flatcars travel the exact same distance} in the same time interval.
  \end{itemize}

  \item \textbf{(Variable-Mass Rocket with External Force):} Generalize Eq.\ (6.2) to include external force $\vect{F}$.
  \begin{itemize}
    \item At time $t$, $\vect{P}_i = M\vect{V}$. At $t+\diff t$, $\vect{P}_f = (M+\diff M)(\vect{V}+\diff\vect{V}) + (-\diff M)(\vect{V}+\vect{u})$, where $\vect{u}$ is exhaust velocity relative to the rocket. Then $\diff\vect{P} = \vect{P}_f - \vect{P}_i = M\diff\vect{V} - \vect{u}\diff M = \vect{F}\diff t \implies \eqbox{M\frac{\diff\vect{V}}{\diff t} - \vect{u}\frac{\diff M}{\diff t} = \vect{F}}$ (or $M\frac{\diff\vect{V}}{\diff t} + \vect{u}\frac{\diff M}{\diff t} = \vect{F}$ if $\vect{u}$ is defined opposite to $\vect{V}$).
  \end{itemize}

  \item \textbf{(What Pushes a Rocket Upward?):} High-pressure gas inside the combustion chamber pushes in all directions. At the front wall of the chamber, the gas exerts a large forward pressure force; at the open rear nozzle, the gas escapes so there is no opposing backward wall force. Thus the internal pressure imbalance (reaction to the expelled gas) pushes the rocket upward.

  \item \textbf{(Head-on Limit of Glancing Collision):} Set $\theta_1 = 0 \implies \cos\theta_1 = 1$ in Eq.\ \eqref{eq:gl_ratio}:
  \begin{equation*}
    \frac{V_1'}{V_1} = \frac{M_1}{M_1+M_2}\!\left[1 \pm \sqrt{1 - 1 + \left(\frac{M_2}{M_1}\right)^2}\right] = \frac{M_1}{M_1+M_2}\!\left(1 \pm \frac{M_2}{M_1}\right) = \eqbox{\frac{M_1 \pm M_2}{M_1 + M_2}}
  \end{equation*}

  \item \textbf{(1D Elastic Collisions of $m$ and $3m$):}
  \begin{itemize}
    \item \textbf{(a)} $M_1 = m, M_2 = 3m$: $V_1' = \frac{m-3m}{m+3m}v = \eqbox{-\frac{v}{2}}$, $V_2' = \frac{2m}{m+3m}v = \eqbox{+\frac{v}{2}}$.
    \item \textbf{(b)} $M_1 = 3m, M_2 = m$: $V_1' = \frac{3m-m}{3m+m}v = \eqbox{+\frac{v}{2}}$, $V_2' = \frac{2(3m)}{3m+m}v = \eqbox{+\frac{3v}{2}}$.
  \end{itemize}

  \item \textbf{(Null Solution of Eq.\ 6.15):} For $M_1 = M_2$, the upper/lower solution $V_1' = 0$ in $\frac{1}{2}M_1 V_1^2 = \frac{1}{2}M_1 V_1'^2 + \frac{1}{2}M_2 V_2'^2$ gives $V_2' = V_1$. Momentum conservation $\vect{V}_1 = \vect{V}_2'$ then requires $\theta_2 = 0$, which is a head-on collision.

  \item \textbf{(Sign of $V_2' = \pm V_1\sin\theta_1$):} Since $\theta_1 \in (0, \pi/2)$ for a glancing collision with $M_1=M_2$, $\sin\theta_1 > 0$. Speed $V_2'$ is a magnitude and must be non-negative, so only $V_2' = +V_1\sin\theta_1$ is physical.

  \item \textbf{(Limit $M_1 \ll M_2$ in Glancing Collision):} Divide numerator and denominator of Eq.\ \eqref{eq:gl_ratio} by $M_2$:
  \begin{equation*}
    \frac{V_1'}{V_1} = \frac{M_1/M_2}{1 + M_1/M_2}\!\left[\cos\theta_1 + \sqrt{\cos^2\theta_1 - 1 + \left(\frac{M_2}{M_1}\right)^2}\right] \xrightarrow{M_1/M_2 \to 0} \frac{M_1}{M_2}\!\left(\frac{M_2}{M_1}\right) = \eqbox{1 \implies V_1' \approx V_1}
  \end{equation*}

  \item \textbf{(Billiard Ball Glancing Collision at $\theta_1 = 20^\circ$):} Since $M_1 = M_2$:
  \begin{equation*}
    \theta_2 = 90^\circ - \theta_1 = \eqbox{70^\circ}, \qquad V_1' = v\cos 20^\circ = \eqbox{0.940v}, \qquad V_2' = v\sin 20^\circ = \eqbox{0.342v}
  \end{equation*}

  \item \textbf{(Coefficient of Restitution $e$):} For a completely inelastic collision, $V_1' = V_2'$, so $e = \frac{|V_2'-V_1'|}{|V_2-V_1|} = \eqbox{0}$. For a 1D elastic collision, Eq.\ \eqref{eq:rel_vel_1D} gives $V_2' - V_1' = -(V_2 - V_1)$, so $e = \frac{|-(V_2-V_1)|}{|V_2-V_1|} = \eqbox{1}$.

  \item \textbf{($Q=0$ for 1D Elastic Collision):} Using Eqs.\ (6.13)--(6.14):
  \begin{equation*}
    T_f = \frac{1}{2}M_1 V_1^2\!\left(\frac{M_1-M_2}{M_1+M_2}\right)^2 + \frac{1}{2}M_2 V_1^2\!\left(\frac{2M_1}{M_1+M_2}\right)^2 = \frac{1}{2}M_1 V_1^2\frac{(M_1-M_2)^2 + 4M_1 M_2}{(M_1+M_2)^2} = \frac{1}{2}M_1 V_1^2 = T_i \implies \eqbox{Q = T_f - T_i = 0}
  \end{equation*}

  \item \textbf{(Impulse of Sinusoidal Force):} $F = 3\sin 5t\text{ N}$ on $m = 2\text{ kg}$ from $t=0$ to $t=\pi/10\text{ s}$ ($v_0=0$):
  \begin{equation*}
    J = \int_0^{\pi/10} 3\sin 5t\dd t = \left[-\frac{3}{5}\cos 5t\right]_0^{\pi/10} = -\frac{3}{5}(0 - 1) = 0.6\text{ N}\cdot\text{s} \implies v_f = \frac{J}{m} = \frac{0.6}{2} = \eqbox{0.3\text{ m/s}}
  \end{equation*}

  \item \textbf{(Piecewise Force Impulse):} $F(t) = 2t$ on $[0,2]$, $F(t)=4$ on $[2,5]$, $F(t)=-t$ on $[5,7]$:
  \begin{equation*}
    p_f = \int_0^2 2t\dd t + \int_2^5 4\dd t + \int_5^7 (-t)\dd t = [t^2]_0^2 + 4(3) - \left[\frac{t^2}{2}\right]_5^7 = 4 + 12 - \frac{49-25}{2} = 16 - 12 = \eqbox{4\text{ kg}\cdot\text{m/s}}
  \end{equation*}
\end{enumerate}

\subsubsection{Part B: Unique \& Challenge Problems from Young \& Freedman (PDF 1)}

\begin{enumerate}[label=\textbf{Prob.\ \arabic*},leftmargin=4.2em]
  \setcounter{enumi}{0}
  \item \textbf{[Challenge Problem 7.87, p.\ 239] Proton in 1D Potential $U(x) = \frac{\alpha}{x^2} - \frac{\beta}{x}$:}
  Released from rest at $x_0 = \alpha/\beta$.
  \begin{itemize}
    \item \textbf{(a)} Factor out $\frac{\alpha}{x_0^2}$ using $\beta = \frac{\alpha}{x_0}$:
    \begin{equation*}
      U(x) = \frac{\alpha}{x^2} - \frac{\alpha}{x_0 x} = \eqbox{\frac{\alpha}{x_0^2}\!\left[\left(\frac{x_0}{x}\right)^2 - \frac{x_0}{x}\right]}, \qquad U(x_0) = 0
    \end{equation*}
    \item \textbf{(b)} Since $K_1 + U(x_0) = 0 + 0 = 0$, conservation of energy $\frac{1}{2}mv^2 + U(x) = 0$ gives:
    \begin{equation*}
      v(x) = \sqrt{-\frac{2U(x)}{m}} = \eqbox{\sqrt{\frac{2\alpha}{m x_0^2}\!\left[\frac{x_0}{x} - \left(\frac{x_0}{x}\right)^2\right]}}
    \end{equation*}
    The proton moves from $x_0$ toward $+\infty$, accelerating at first and then slowing to $v\to 0$ as $x\to\infty$.
    \item \textbf{(c)--(d)} Maximum speed occurs where $U(x)$ is minimum, i.e., $F_x = -\frac{\diff U}{\diff x} = \frac{2\alpha}{x^3} - \frac{\beta}{x^2} = 0$:
    \begin{equation*}
      x_m = \frac{2\alpha}{\beta} = \eqbox{2x_0}, \qquad v_{\max} = \eqbox{\sqrt{\frac{\alpha}{2m x_0^2}} = \frac{\beta}{\sqrt{2m\alpha}}}, \qquad F_x(2x_0) = \eqbox{0}
    \end{equation*}
    \item \textbf{(e)--(f)} If released at $x_1 = \frac{3\alpha}{\beta} = 3x_0$, then $E = U(3x_0) = -\frac{2\alpha}{9x_0^2} < 0$. Setting $U(x) = E$ gives $\left(\frac{x_0}{x}\right)^2 - \frac{x_0}{x} + \frac{2}{9} = \left(\frac{x_0}{x} - \frac{1}{3}\right)\!\left(\frac{x_0}{x} - \frac{2}{3}\right) = 0$. Thus the proton oscillates periodically between $x_{\min} = \eqbox{\frac{3}{2}x_0 = \frac{3\alpha}{2\beta}}$ and $x_{\max} = \eqbox{3x_0 = \frac{3\alpha}{\beta}}$. For release at $x_0$, $x_{\min} = x_0$ and $x_{\max} = \infty$.
  \end{itemize}

  \item \textbf{[Prob.\ 8.66, p.\ 272] Girl Acted on by Time-Varying Force:} $m = 40.0\text{ kg}$, $p_{0x} = +90.0\text{ kg}\cdot\text{m/s}$ (east), $F_x(t) = -(8.20\text{ N/s})t$ (west).
  \begin{itemize}
    \item \textbf{(a)} Time when $p_x(t) = -60.0\text{ kg}\cdot\text{m/s}$:
    \begin{equation*}
      \Delta p_x = -60.0 - 90.0 = -150.0\text{ N}\cdot\text{s} = \int_0^t -(8.20)t'\dd t' = -4.10t^2 \implies t = \sqrt{\frac{150.0}{4.10}} = \eqbox{6.05\text{ s}}
    \end{equation*}
    \item \textbf{(b)} Work done by force: $W = K_2 - K_1 = \frac{p_2^2 - p_1^2}{2m} = \frac{(-60.0)^2 - (90.0)^2}{2(40.0)} = \eqbox{-56.3\text{ J}}$.
    \item \textbf{(c)} Acceleration magnitude at $t = 6.05\text{ s}$: $a = \frac{F(6.05)}{m} = \frac{8.20(6.05)}{40.0} = \eqbox{1.24\text{ m/s}^2}$.
  \end{itemize}

  \item \textbf{[Prob.\ 8.86, p.\ 274] Two Identical Masses in a Hemispherical Bowl:}
  Mass $A$ starts at the rim (height $R$ above bottom); mass $B$ starts at rest at the bottom.
  \begin{itemize}
    \item Speed of $A$ just before impact at the bottom: $\frac{1}{2}mv_1^2 = mgR \implies v_1 = \sqrt{2gR}$.
    \item Completely inelastic collision at the bottom: $mv_1 = (2m)v_2 \implies v_2 = \frac{1}{2}v_1$.
    \item Final height $h$: $\frac{1}{2}(2m)v_2^2 = (2m)gh \implies h = \frac{v_2^2}{2g} = \frac{v_1^2}{8g} = \eqbox{\frac{R}{4}}$.
  \end{itemize}

  \item \textbf{[Prob.\ 8.88, p.\ 274] Dart and Lead Sphere Complete Vertical Loop:}
  Sphere $M = 20.00\text{ kg}$ on a thin wire of length $L = 3.50\text{ m}$; dart $m = 5.00\text{ kg}$.
  \begin{itemize}
    \item For a flexible wire of length $L$, the combination must have minimum speed $v_{\text{top}} = \sqrt{gL}$ at the top of the circle ($y = 2L$) so tension does not go slack.
    \item Energy conservation from bottom ($v_2$) to top ($v_{\text{top}}$): $\frac{1}{2}M_{\text{tot}}v_2^2 = M_{\text{tot}}g(2L) + \frac{1}{2}M_{\text{tot}}(gL) \implies v_2 = \sqrt{5gL}$.
    \item Momentum conservation during impact: $mv_0 = (m+M)v_2 \implies v_0 = \frac{5.00+20.00}{5.00}\sqrt{5(9.80)(3.50)} = 5(13.09) = \eqbox{65.5\text{ m/s}}$.
  \end{itemize}

  \item \textbf{[Prob.\ 8.93 \& 8.94, p.\ 274] Energy Sharing in Head-On Elastic Collisions:}
  Incident mass $m$ with kinetic energy $K_0$ hits stationary mass $M$.
  \begin{itemize}
    \item \textbf{(a)} Final speed of $m$ is $v_{m2} = \frac{m-M}{m+M}v_0$, so $K_{m2} = \left(\frac{m-M}{m+M}\right)^2 K_0$. Kinetic energy lost by $m$ (and gained by $M$) is:
    \begin{equation*}
      \Delta K_{\text{lost}} = K_0 - K_{m2} = K_0\!\left[1 - \frac{(m-M)^2}{(m+M)^2}\right] = \eqbox{\frac{4mM}{(M+m)^2}K_0}
    \end{equation*}
    \item \textbf{(b)--(c)} Maximum energy loss (100\%) occurs when $\eqbox{M = m}$, where $\eqbox{v_{m2} = 0}$. Equal energy sharing ($\Delta K_{\text{lost}} = \frac{1}{2}K_0$) occurs when $\left(\frac{m_A-m_B}{m_A+m_B}\right)^2 = \frac{1}{2} \implies \frac{m_A}{m_B} = \eqbox{3 \pm 2\sqrt{2}}$ ($5.83$ or $0.172$).
  \end{itemize}

  \item \textbf{[Prob.\ 8.97, p.\ 275] Jack and Jill Jumping Off a Crate:}
  $m_J = 75.0\text{ kg}$, $m_{Ji} = 45.0\text{ kg}$, crate $m_C = 15.0\text{ kg}$. Each jumps at $v_{\text{rel}} = 4.00\text{ m/s}$ relative to the crate.
  \begin{itemize}
    \item \textbf{(a) Simultaneous jump:} $0 = (m_J+m_{Ji})(v_C - v_{\text{rel}}) + m_C v_C \implies v_C = \frac{120.0}{135.0}(4.00) = \eqbox{3.56\text{ m/s}}$.
    \item \textbf{(b) Jack first, then Jill:} After Jack jumps, $v_{C1} = \frac{m_J}{m_J+m_{Ji}+m_C}v_{\text{rel}} = \frac{75.0}{135.0}(4.00) = 2.222\text{ m/s}$. Next Jill jumps from the moving crate: $v_{C2} = v_{C1} + \frac{m_{Ji}}{m_{Ji}+m_C}v_{\text{rel}} = 2.222 + \frac{45.0}{60.0}(4.00) = \eqbox{5.22\text{ m/s}}$.
    \item \textbf{(c) Jill first, then Jack:} $v_C = \frac{45.0}{135.0}(4.00) + \frac{75.0}{90.0}(4.00) = 1.333 + 3.333 = \eqbox{4.67\text{ m/s}}$.
  \end{itemize}

  \item \textbf{[Prob.\ 8.98, p.\ 275] Superball Double Bounce:}
  Large ball ($M \gg m$) rebounds with $+\vect{v}$ and collides head-on with small ball ($m$) descending with $-\vect{v}$.
  \begin{itemize}
    \item \textbf{(a)} From Eq.\ \eqref{eq:elastic_gen} with $m \ll M$, $v_{m2} = \frac{m-M}{m+M}(-v) + \frac{2M}{m+M}(+v) \approx +v + 2v = \eqbox{3\vect{v}}$.
    \item \textbf{(b)} Since $h \propto v^2$, the rebound height ratio is $\left(\frac{3v}{v}\right)^2 = \eqbox{9}$.
  \end{itemize}

  \item \textbf{[Prob.\ 8.100 \& 8.101, p.\ 275] Explosion Energy Sharing:}
  Mass $m = m_A + m_B$ at rest explodes releasing energy $Q = K_A + K_B$.
  \begin{itemize}
    \item Since $p_A = p_B = p$, $K_A = \frac{p^2}{2m_A}$ and $K_B = \frac{p^2}{2m_B}$, so $\frac{K_A}{K_B} = \frac{m_B}{m_A}$:
    \begin{equation*}
      \eqbox{K_A = \frac{m_B}{m_A + m_B}Q, \qquad K_B = \frac{m_A}{m_A + m_B}Q}
    \end{equation*}
    For neutron decay ($m_p = 1836 m_e$), the proton gets $\frac{m_e}{m_p + m_e} = \frac{1}{1837} = \eqbox{5.44\times 10^{-4}\ (0.0544\%)}$ of $Q$.
  \end{itemize}

  \item \textbf{[Prob.\ 8.105 \& 8.106, p.\ 275] Walking on a Floating Log / Canoe:}
  In Prob.\ 8.106, woman ($m_w = 45.0\text{ kg}$) walks $d_{\text{rel}} = 3.00\text{ m}$ from $1.00\text{ m}$ from one end to $1.00\text{ m}$ from the other end of a $5.00\text{-m}$ canoe ($m_c = 60.0\text{ kg}$).
  \begin{itemize}
    \item Since $x_{\text{cm}}$ does not move: $m_w(\Delta x_{\text{canoe}} + d_{\text{rel}}) + m_c\Delta x_{\text{canoe}} = 0 \implies |\Delta x_{\text{canoe}}| = \frac{m_w d_{\text{rel}}}{m_w + m_c} = \frac{45.0(3.00)}{105.0} = \eqbox{1.29\text{ m}}$.
    \item Similarly, in Prob.\ 8.105 (Ernie $40.0\text{ kg}$ walks $3.0\text{ m}$ on a $20.0\text{-kg}$ log with Burt $30.0\text{ kg}$), $|\Delta x_{\text{log}}| = \frac{40.0(3.0)}{40.0+30.0+20.0} = \eqbox{1.3\text{ m}}$.
  \end{itemize}

  \item \textbf{[Prob.\ 8.113, p.\ 276] Multistage Rocket:}
  Stage 1: $m_1 = 12,000\text{ kg}$ ($9000\text{ kg}$ fuel, $3000\text{ kg}$ dry). Stage 2: $m_2 = 1000\text{ kg}$ ($700\text{ kg}$ fuel, $300\text{ kg}$ dry).
  \begin{itemize}
    \item \textbf{(a) Single-stage equivalent:} $m_0 = 13,000\text{ kg}$, $m_f = 13,000 - 9700 = 3300\text{ kg} \implies v = v_{\text{ex}}\ln\!\left(\frac{13,000}{3300}\right) = \eqbox{1.37v_{\text{ex}}}$.
    \item \textbf{(b) Speed after Stage 1 burnout:} $m_0 = 13,000\text{ kg}$, $m_f = 4000\text{ kg} \implies v_1 = v_{\text{ex}}\ln\!\left(\frac{13,000}{4000}\right) = \eqbox{1.18v_{\text{ex}}}$.
    \item \textbf{(c) Final speed of Stage 2:} $v_2 = v_1 + v_{\text{ex}}\ln\!\left(\frac{1000}{300}\right) = (1.18 + 1.20)v_{\text{ex}} = \eqbox{2.38v_{\text{ex}}}$.
    \item \textbf{(d)} For $v_2 = 7.00\text{ km/s}$: $v_{\text{ex}} = \frac{7.00\text{ km/s}}{2.38} = \eqbox{2.94\text{ km/s}}$.
  \end{itemize}

  \item \textbf{[Challenge Problem 8.114, p.\ 276] Variable-Mass Raindrop (Krane Problem):}
  Mass grows with distance fallen: $m = kx$, $\frac{\diff m}{\diff t} = kv$, and $F_{\text{ext}} = mg = \frac{\diff p}{\diff t} = m\frac{\diff v}{\diff t} + v\frac{\diff m}{\diff t}$, which simplifies to $xg = x\frac{\diff v}{\diff t} + v^2$.
  \begin{itemize}
    \item \textbf{(a)} Substitute $v = at$ and $x = \frac{1}{2}at^2$:
    \begin{equation*}
      \left(\frac{1}{2}at^2\right)g = \left(\frac{1}{2}at^2\right)a + (at)^2 = \frac{3}{2}a^2 t^2 \implies \eqbox{a = \frac{g}{3} = 3.27\text{ m/s}^2}
    \end{equation*}
    \item \textbf{(b)} In $t = 3.00\text{ s}$: $x = \frac{1}{2}\!\left(\frac{9.80}{3}\right)(3.00)^2 = \eqbox{14.7\text{ m}}$.
    \item \textbf{(c)} Given $k = 2.00\text{ g/m}$: $m = kx = (2.00\text{ g/m})(14.7\text{ m}) = \eqbox{29.4\text{ g}}$.
  \end{itemize}

  \item \textbf{[Challenge Problems 8.115 \& 8.116, p.\ 276--277] Continuous Center-of-Mass Integrals:}
  \begin{itemize}
    \item \textbf{8.115(a) Uniform rod ($\rho = M/V$, $\diff m = \frac{M}{L}\dd x$):} $x_{\text{cm}} = \frac{1}{M}\int_0^L x\frac{M}{L}\dd x = \eqbox{\frac{L}{2}}$.
    \item \textbf{8.115(b) Linearly varying rod ($\rho = \alpha x$, $\diff m = \alpha A x\dd x$):} $M = \int_0^L \alpha A x\dd x = \frac{1}{2}\alpha A L^2$, and $x_{\text{cm}} = \frac{1}{M}\int_0^L \alpha A x^2\dd x = \frac{\frac{1}{3}\alpha A L^3}{\frac{1}{2}\alpha A L^2} = \eqbox{\frac{2}{3}L}$.
    \item \textbf{8.116 Semicircular plate of radius $a$ in upper half-plane ($y\ge 0$):} By symmetry, $x_{\text{cm}} = 0$. Using polar coordinates ($y = r\sin\theta$, $\diff m = \rho t r\dd r\dd\theta$, $M = \frac{1}{2}\rho\pi a^2 t$):
    \begin{equation*}
      y_{\text{cm}} = \frac{\rho t}{M}\int_0^a r^2\dd r \int_0^\pi \sin\theta\dd\theta = \frac{\rho t}{\frac{1}{2}\rho\pi a^2 t}\!\left(\frac{a^3}{3}\right)(2) = \eqbox{\frac{4a}{3\pi} \approx 0.424a}
    \end{equation*}
  \end{itemize}
\end{enumerate}

% ============================================================================
% FORMULA SHEET — MODULE 1
% ============================================================================
\formulasheet{Module 1 Formula Sheet: Momentum, Collisions, and Rockets}

\begin{fs}{1. Momentum, Newton's Second Law, and Impulse}
\begin{align*}
  \vect{p} &= m\vect{v}, \qquad K = \frac{p^2}{2m}, \qquad \sum\vect{F} = \frac{\diff\vect{p}}{\diff t} = m\frac{\diff\vect{v}}{\diff t} + \vect{v}\frac{\diff m}{\diff t}\\
  \vect{J} &= \int_{t_1}^{t_2}\sum\vect{F}\dd t = \vect{F}_{\text{av}}\Delta t = \vect{p}_2 - \vect{p}_1
\end{align*}
\end{fs}

\begin{fs}{2. Collisions (1D and 2D) and Coefficient of Restitution}
\begin{align*}
  \text{Completely Inelastic:}\quad & \vect{v}_2 = \frac{m_A\vect{v}_{A1}+m_B\vect{v}_{B1}}{m_A+m_B}, \qquad \frac{K_2}{K_1} = \frac{m_A}{m_A+m_B}\quad (v_{B1}=0)\\
  \text{1D Elastic ($v_{B1x}=0$):}\quad & v_{A2x} = \frac{m_A-m_B}{m_A+m_B}v_{A1x}, \qquad v_{B2x} = \frac{2m_A}{m_A+m_B}v_{A1x}, \qquad v_{B2x}-v_{A2x} = -(v_{B1x}-v_{A1x})\\
  \text{2D Glancing Elastic:}\quad & \frac{V_1'}{V_1} = \frac{M_1}{M_1+M_2}\!\left[\cos\theta_1 \pm \sqrt{\cos^2\theta_1 - 1 + \left(\frac{M_2}{M_1}\right)^2}\right], \qquad \theta_1+\theta_2 = \frac{\pi}{2}\ (M_1=M_2)\\
  \text{Restitution \& $Q$-value:}\quad & e = \frac{|V_2'-V_1'|}{|V_2-V_1|}, \qquad Q = T_f - T_i
\end{align*}
\end{fs}

\begin{fs}{3. Center of Mass, Reduced Mass, and Rocket Propulsion}
\begin{align*}
  \vect{r}_{\text{cm}} &= \frac{\sum m_i\vect{r}_i}{M} = \frac{1}{M}\int\vect{r}\dd m, \qquad \vect{P}_{\text{tot}} = M\vect{v}_{\text{cm}}, \qquad \sum\vect{F}_{\text{ext}} = M\vect{a}_{\text{cm}}, \qquad \mu = \frac{m_1 m_2}{m_1+m_2}\\
  F_{\text{thrust}} &= -v_{\text{ex}}\frac{\diff m}{\diff t}, \qquad a = -\frac{v_{\text{ex}}}{m}\frac{\diff m}{\diff t} - g, \qquad v - v_0 = v_{\text{ex}}\ln\!\left(\frac{m_0}{m}\right) - gt
\end{align*}
\end{fs}

% ============================================================================
% MODULE 2: ROTATION OF RIGID BODIES (PDF 4: Y&F CH. 9)
% ============================================================================
\part{Rotation of Rigid Bodies: Kinematics and Moment of Inertia}

\chapter{Rotational Kinematics and Linear--Angular Relations}
\begin{syll}
\textbf{Syllabus Coverage (PDF 4):} Young \& Freedman Chapter 9 (Sections 9.1--9.3). Radian measure, angular velocity, angular acceleration, constant angular acceleration kinematics, and tangential/centripetal linear kinematics.
\end{syll}

\section{Angular Position, Velocity, and Acceleration\bk{9.1}}

\begin{goal}
Define radian angle $\theta$, angular velocity $\omega_z$, and angular acceleration $\alpha_z$ for a rigid body rotating about a fixed $z$-axis.
\end{goal}

\begin{pic}
\begin{tikzpicture}[scale=0.95]
  \draw[ax] (-1.8,0) -- (2.6,0) node[right,font=\small]{$x$};
  \draw[ax] (0,-1.6) -- (0,2.2) node[above,font=\small]{$y$};
  \draw[thick,teal,fill=tealsoft!50] (0,0) circle (1.4cm);
  \draw[thick,ink] (0,0) -- (1.4,0);
  \draw[thick,teal] (0,0) -- (0.99,0.99) node[midway,above left,font=\footnotesize]{$r$};
  \fill[rose] (0.99,0.99) circle (2.5pt) node[above right,font=\footnotesize]{$P$};
  \draw[ultra thick,amber] (1.4,0) arc (0:45:1.4) node[midway,right,font=\footnotesize]{$s = r\theta$};
  \draw[slate] (0.55,0) arc (0:45:0.55) node[midway,right,font=\footnotesize]{$\theta$};
\end{tikzpicture}
\figcap{Angle $\theta = s/r$ in radians on a circle of radius $r$.}
\end{pic}

\step{Step 1: Relate arc length $s$ and radius $r$.}
\begin{equation}\label{eq:rad_def}
  \eqbox{\theta = \frac{s}{r} \iff s = r\theta} \qquad \left(1\text{ rev} = 360^\circ = 2\pi\text{ rad},\quad 1\text{ rad} = \frac{360^\circ}{2\pi} \approx 57.3^\circ\right)
\end{equation}

\step{Step 2: Define average and instantaneous angular velocity $\omega_z$.}
\begin{equation}\label{eq:omega_def}
  \omega_{\text{av-}z} = \frac{\theta_2 - \theta_1}{t_2 - t_1} = \frac{\Delta\theta}{\Delta t}, \qquad \eqbox{\omega_z = \lim_{\Delta t\to 0}\frac{\Delta\theta}{\Delta t} = \frac{\diff\theta}{\diff t}} \qquad \left(1\text{ rpm} = \frac{2\pi}{60}\text{ rad/s}\right)
\end{equation}

\step{Step 3: Define average and instantaneous angular acceleration $\alpha_z$.}
\begin{equation}\label{eq:alpha_def}
  \alpha_{\text{av-}z} = \frac{\omega_{2z} - \omega_{1z}}{t_2 - t_1} = \frac{\Delta\omega_z}{\Delta t}, \qquad \eqbox{\alpha_z = \lim_{\Delta t\to 0}\frac{\Delta\omega_z}{\Delta t} = \frac{\diff\omega_z}{\diff t} = \frac{\diff^2\theta}{\diff t^2}}
\end{equation}
Using the chain rule (\tref{tool:deriv}), we can also eliminate $\diff t$:
\begin{equation}\label{eq:alpha_dtheta}
  \alpha_z = \frac{\diff\omega_z}{\diff\theta}\frac{\diff\theta}{\diff t} = \omega_z\frac{\diff\omega_z}{\diff\theta}
\end{equation}

\section{Rotation with Constant Angular Acceleration\bk{9.2}}

\begin{goal}
Derive the four kinematic equations for fixed-axis rotation when $\alpha_z = \text{constant}$.
\end{goal}

\step{Step 1: Integrate $\frac{\diff\omega_z}{\diff t} = \alpha_z$ from $t=0$ ($\omega_{0z}$) to $t$ ($\omega_z$).}
\begin{equation}\label{eq:rot_kin1}
  \int_{\omega_{0z}}^{\omega_z}\diff\omega_z' = \int_0^t \alpha_z\dd t' \implies \eqbox{\omega_z = \omega_{0z} + \alpha_z t}
\end{equation}

\step{Step 2: Integrate $\frac{\diff\theta}{\diff t} = \omega_z = \omega_{0z} + \alpha_z t$ from $t=0$ ($\theta_0$) to $t$ ($\theta$).}
\begin{equation}\label{eq:rot_kin2}
  \int_{\theta_0}^\theta \diff\theta' = \int_0^t (\omega_{0z} + \alpha_z t')\dd t' \implies \eqbox{\theta = \theta_0 + \omega_{0z}t + \frac{1}{2}\alpha_z t^2}
\end{equation}

\step{Step 3: Express $\theta - \theta_0$ in terms of average angular velocity.}
Since $\omega_z$ varies linearly with $t$, $\omega_{\text{av-}z} = \frac{1}{2}(\omega_{0z} + \omega_z) = \frac{\theta - \theta_0}{t}$:
\begin{equation}\label{eq:rot_kin3}
  \eqbox{\theta - \theta_0 = \frac{1}{2}(\omega_{0z} + \omega_z)t}
\end{equation}

\step{Step 4: Eliminate $t$ (or integrate $\omega_z\dd\omega_z = \alpha_z\dd\theta$).}
Substitute $t = \frac{\omega_z - \omega_{0z}}{\alpha_z}$ from Eq.\ \eqref{eq:rot_kin1} into Eq.\ \eqref{eq:rot_kin3}:
\begin{equation}\label{eq:rot_kin4}
  \theta - \theta_0 = \frac{1}{2}(\omega_{0z} + \omega_z)\!\left(\frac{\omega_z - \omega_{0z}}{\alpha_z}\right) \implies \eqbox{\omega_z^2 = \omega_{0z}^2 + 2\alpha_z(\theta - \theta_0)}
\end{equation}

\section{Relating Linear and Angular Kinematics\bk{9.3}}

\begin{goal}
Derive the linear speed $v$, tangential acceleration $a_{\text{tan}}$, and centripetal acceleration $a_{\text{rad}}$ of a point at distance $r$ from the rotation axis.
\end{goal}

\begin{pic}
\begin{tikzpicture}[scale=0.95]
  \draw[ax] (-1.8,0) -- (2.6,0) node[right,font=\small]{$x$};
  \draw[ax] (0,-1.6) -- (0,2.3) node[above,font=\small]{$y$};
  \draw[thick,teal,dashed] (0,0) circle (1.5cm);
  \coordinate (P) at (40:1.5);
  \draw[thick,ink] (0,0) -- (P) node[midway,below right,font=\footnotesize]{$r$};
  \fill[teal] (P) circle (3pt) node[right=2pt,font=\footnotesize]{$P$};
  \draw[vec,amber,thick] (P) -- ++(130:1.2) node[above,font=\footnotesize]{$\vect{a}_{\text{tan}} = r\alpha\uv{\theta}$};
  \draw[vec,rose,thick] (P) -- ++(220:0.95) node[above left,font=\footnotesize]{$\vect{a}_{\text{rad}} = \omega^2 r(-\uv{r})$};
  \draw[vec,bluec,ultra thick] (P) -- ($(P)+(130:1.2)+(220:0.95)$) node[left,font=\footnotesize]{$\vect{a}$};
\end{tikzpicture}
\figcap{Tangential acceleration $a_{\text{tan}} = r\alpha$ and centripetal acceleration $a_{\text{rad}} = \omega^2 r$ at point $P$.}
\end{pic}

\step{Step 1: Differentiate $s = r\theta$ with respect to time $t$ (with $r$ constant).}
\begin{equation}\label{eq:v_romega}
  \left|\frac{\diff s}{\diff t}\right| = r\left|\frac{\diff\theta}{\diff t}\right| \implies \eqbox{v = r\omega}
\end{equation}

\step{Step 2: Differentiate $v = r\omega$ with respect to time $t$ to obtain tangential acceleration $a_{\text{tan}}$.}
\begin{equation}\label{eq:a_tan}
  \eqbox{a_{\text{tan}} = \frac{\diff v}{\diff t} = r\frac{\diff\omega}{\diff t} = r\alpha}
\end{equation}

\step{Step 3: Substitute $v = r\omega$ into centripetal acceleration $a_{\text{rad}} = \frac{v^2}{r}$.}
\begin{equation}\label{eq:a_rad}
  \eqbox{a_{\text{rad}} = \frac{v^2}{r} = \frac{(r\omega)^2}{r} = \omega^2 r}, \qquad \eqbox{a = \sqrt{a_{\text{tan}}^2 + a_{\text{rad}}^2} = r\sqrt{\alpha^2 + \omega^4}}
\end{equation}

\begin{trap}
Equations \eqref{eq:rad_def}, \eqref{eq:v_romega}, \eqref{eq:a_tan}, and \eqref{eq:a_rad} are valid \textbf{only} when $\theta$ is in $\text{rad}$, $\omega$ is in $\text{rad/s}$, and $\alpha$ is in $\text{rad/s}^2$. Always convert $\text{rev}$ and $\text{rpm}$ to radians first!
\end{trap}

% ============================================================================
% CHAPTER 4: ROTATIONAL ENERGY AND MOMENT OF INERTIA
% ============================================================================
\chapter{Rotational Energy and Moment-of-Inertia Calculations}
\begin{syll}
\textbf{Syllabus Coverage (PDF 4):} Young \& Freedman Chapter 9 (Sections 9.4--9.6). Rotational kinetic energy, moment of inertia, gravitational potential energy of an extended body, parallel-axis and perpendicular-axis theorems, and continuous moment-of-inertia integrals.
\end{syll}

\section{Rotational Kinetic Energy and Moment of Inertia\bk{9.4}}

\begin{goal}
Derive the rotational kinetic energy $K = \frac{1}{2}I\omega^2$ and gravitational potential energy $U = Mgy_{\text{cm}}$ of a rigid body.
\end{goal}

\step{Step 1: Sum the kinetic energies of all particles $m_i$ in the rotating rigid body.}
Particle $i$ is at perpendicular distance $r_i$ from the rotation axis, so its speed is $v_i = r_i\omega$:
\begin{align}
  K &= \sum_i \frac{1}{2}m_i v_i^2 = \sum_i \frac{1}{2}m_i(r_i\omega)^2 = \frac{1}{2}\!\left(\sum_i m_i r_i^2\right)\omega^2\label{eq:K_rot_sum}
\end{align}
Define the moment of inertia $I$ about the rotation axis:
\begin{equation}\label{eq:I_def_rot}
  \eqbox{I = \sum_i m_i r_i^2 = \int r^2\dd m}, \qquad \eqbox{K = \frac{1}{2}I\omega^2}
\end{equation}

\step{Step 2: Sum the gravitational potential energies of all particles $m_i$.}
\begin{equation}\label{eq:U_grav_ext}
  U = \sum_i m_i g y_i = \left(\sum_i m_i y_i\right)g = \eqbox{M g y_{\text{cm}}}
\end{equation}

\section{Parallel-Axis and Perpendicular-Axis Theorems\bk{9.5, 7.5.1}}

\begin{goal}
Prove the parallel-axis theorem $I_P = I_{\text{cm}} + Md^2$ for any 3D body and the perpendicular-axis theorem $I_O = I_x + I_y$ for a thin planar lamina.
\end{goal}

\begin{pic}
\begin{tikzpicture}[scale=0.9]
  \draw[thick,teal,fill=tealsoft!60] plot[smooth cycle,tension=0.8] coordinates {(-1.5,-1.0) (1.2,-1.3) (2.6,0.2) (1.8,1.8) (-0.8,1.6) (-2.0,0.4)};
  \draw[ax] (-2.4,0) -- (3.0,0) node[right,font=\small]{$x$};
  \draw[ax] (0,-1.6) -- (0,2.2) node[above,font=\small]{$y$};
  \fill[ink] (0,0) circle (2.5pt) node[below left,font=\footnotesize]{$O\ (\text{cm})$};
  \coordinate (P) at (1.1,0.7);
  \fill[amber] (P) circle (2.5pt) node[below right,font=\footnotesize]{$P(a,b)$};
  \draw[thick,amber] (0,0) -- (P) node[midway,above left,font=\footnotesize]{$d$};
  \coordinate (M) at (1.8,1.4);
  \fill[rose] (M) circle (2.5pt) node[above,font=\footnotesize]{$m_i(x_i,y_i)$};
  \draw[dashed,slate] (P) -- (1.8,0.7) -- (M);
\end{tikzpicture}
\figcap{Geometry for the parallel-axis theorem: cm at origin $O(0,0)$ and parallel axis through $P(a,b)$ at distance $d = \sqrt{a^2+b^2}$.}
\end{pic}

\step{Step 1: Proof of the Parallel-Axis Theorem.}
Place the origin at the center of mass ($\text{cm}$), so $x_{\text{cm}} = \frac{1}{M}\sum m_i x_i = 0$ and $y_{\text{cm}} = \frac{1}{M}\sum m_i y_i = 0$. Let the parallel axis pass through $P(a,b)$, where $d^2 = a^2 + b^2$. The squared distance from $m_i(x_i,y_i)$ to $P(a,b)$ is $(x_i - a)^2 + (y_i - b)^2$:
\begin{align}
  I_P &= \sum_i m_i\!\left[(x_i - a)^2 + (y_i - b)^2\right]\nonumber\\
      &= \underbrace{\sum_i m_i(x_i^2 + y_i^2)}_{I_{\text{cm}}} - 2a\underbrace{\sum_i m_i x_i}_{0} - 2b\underbrace{\sum_i m_i y_i}_{0} + \underbrace{(a^2 + b^2)}_{d^2}\underbrace{\sum_i m_i}_{M} = \eqbox{I_{\text{cm}} + Md^2}\label{eq:parallel_axis}
\end{align}

\step{Step 2: Proof of the Perpendicular-Axis Theorem (for a thin flat lamina in the $xy$-plane).}
For a thin lamina in the $xy$-plane ($z_i = 0$), the perpendicular distance of $m_i$ from the $x$-axis is $y_i$, from the $y$-axis is $x_i$, and from the $z$-axis (through $O$) is $r_i = \sqrt{x_i^2 + y_i^2}$:
\begin{equation}\label{eq:perp_axis}
  I_x = \sum_i m_i y_i^2, \qquad I_y = \sum_i m_i x_i^2 \implies I_O = I_z = \sum_i m_i(x_i^2 + y_i^2) = \eqbox{I_x + I_y}
\end{equation}

\section{Moment-of-Inertia Calculations by Integration\bk{9.6, 7.5}}

\begin{goal}
Derive the standard moments of inertia in Table 9.2 using $I = \int r^2\dd m = \int r^2\rho\dd V$.
\end{goal}

\subsection{1. Uniform Slender Rod of Mass $M$ and Length $L$}
Let the rod lie along the $x$-axis with linear mass density $\lambda = \frac{M}{L}$, so $\diff m = \frac{M}{L}\dd x$:
\begin{itemize}
  \item \textbf{Axis through center ($x\in [-L/2, L/2]$):}
  \begin{equation}\label{eq:I_rod_cm}
    I_{\text{cm}} = \int_{-L/2}^{L/2} x^2\!\left(\frac{M}{L}\dd x\right) = \frac{M}{L}\!\left[\frac{x^3}{3}\right]_{-L/2}^{L/2} = \eqbox{\frac{1}{12}ML^2}
  \end{equation}
  \item \textbf{Axis through one end ($x\in [0, L]$):}
  \begin{equation}\label{eq:I_rod_end}
    I_{\text{end}} = \int_0^L x^2\!\left(\frac{M}{L}\dd x\right) = \eqbox{\frac{1}{3}ML^2} \qquad \left(\text{or by parallel axis: }\frac{1}{12}ML^2 + M\!\left(\frac{L}{2}\right)^2 = \frac{1}{3}ML^2\right)
  \end{equation}
\end{itemize}

\subsection{2. Uniform Hollow or Solid Cylinder (Example 9.10)}
Divide a cylinder of length $L$, inner radius $R_1$, outer radius $R_2$, and density $\rho$ into thin cylindrical shells of radius $r$, thickness $\diff r$, volume $\diff V = 2\pi r L\dd r$, and mass $\diff m = 2\pi\rho L r\dd r$. Since total mass is $M = \pi\rho L(R_2^2 - R_1^2)$:
\begin{align}\label{eq:I_cyl}
  I &= \int_{R_1}^{R_2} r^2(2\pi\rho L r\dd r) = 2\pi\rho L\frac{R_2^4 - R_1^4}{4} = \frac{1}{2}\!\left[\pi\rho L(R_2^2 - R_1^2)\right](R_1^2 + R_2^2) = \eqbox{\frac{1}{2}M(R_1^2 + R_2^2)}
\end{align}
\begin{itemize}
  \item \textbf{Solid cylinder/disk ($R_1 = 0, R_2 = R$):} $\eqbox{I = \frac{1}{2}MR^2}$.
  \item \textbf{Thin-walled hollow cylinder/hoop ($R_1 = R_2 = R$):} $\eqbox{I = MR^2}$.
\end{itemize}

\subsection{3. Uniform Solid Sphere of Mass $M$ and Radius $R$ (Example 9.11)}
Slice the sphere into thin disks of radius $r = \sqrt{R^2 - x^2}$ and thickness $\diff x$ at distance $x$ from the center. Each disk has mass $\diff m = \rho\pi(R^2 - x^2)\dd x$ and moment of inertia $\diff I = \frac{1}{2}r^2\dd m = \frac{1}{2}\pi\rho(R^2 - x^2)^2\dd x$. Since $M = \frac{4}{3}\pi R^3\rho$:
\begin{align}\label{eq:I_sphere}
  I &= 2\int_0^R \frac{1}{2}\pi\rho(R^4 - 2R^2 x^2 + x^4)\dd x = \pi\rho\!\left[R^5 - \frac{2}{3}R^5 + \frac{1}{5}R^5\right] = \frac{8\pi\rho R^5}{15} = \left(\frac{4}{3}\pi\rho R^3\right)\!\left(\frac{2}{5}R^2\right) = \eqbox{\frac{2}{5}MR^2}
\end{align}
For a \textbf{thin-walled hollow sphere} of radius $R$, $\diff I = \frac{2}{5}\diff(MR^2) \implies \eqbox{I = \frac{2}{3}MR^2}$.

% ============================================================================
% SOLVED BOOK EXAMPLES — MODULE 2
% ============================================================================
\section{Solved Book Examples (Module 2)}

\begin{example}{9.1}{Calculating angular velocity\bk{p.\ 281}}
\textbf{Problem:} The angular position of a flywheel of diameter $0.36\text{ m}$ ($r = 0.18\text{ m}$) is $\theta = (2.0\text{ rad/s}^3)t^3$. Find: (a) $\theta_1$ at $t_1 = 2.0\text{ s}$ and $\theta_2$ at $t_2 = 5.0\text{ s}$ in rad and degrees, (b) distance $s$ moved by a rim point, (c) $\omega_{\text{av-}z}$ on $[2.0\text{ s}, 5.0\text{ s}]$, and (d) $\omega_z$ at $t_1$ and $t_2$.
\tcblower
\textbf{Solution:}
\begin{itemize}
  \item \textbf{(a)} $\theta_1 = 2.0(2.0)^3 = \eqbox{16\text{ rad} = 920^\circ}$, $\theta_2 = 2.0(5.0)^3 = \eqbox{250\text{ rad} = 14,000^\circ}$.
  \item \textbf{(b)} $\Delta\theta = 250 - 16 = 234\text{ rad} \implies s = r\Delta\theta = (0.18\text{ m})(234\text{ rad}) = \eqbox{42\text{ m}}$.
  \item \textbf{(c)} $\omega_{\text{av-}z} = \frac{234\text{ rad}}{5.0 - 2.0\text{ s}} = \eqbox{78\text{ rad/s} = 740\text{ rev/min}}$.
  \item \textbf{(d)} $\omega_z = \frac{\diff\theta}{\diff t} = (6.0\text{ rad/s}^3)t^2 \implies \omega_{1z} = 6.0(2.0)^2 = \eqbox{24\text{ rad/s}}$, $\omega_{2z} = 6.0(5.0)^2 = \eqbox{150\text{ rad/s}}$.
\end{itemize}
\end{example}

\begin{example}{9.2}{Calculating angular acceleration\bk{p.\ 283}}
\textbf{Problem:} For the flywheel of Example 9.1, find (a) $\alpha_{\text{av-}z}$ between $t_1 = 2.0\text{ s}$ and $t_2 = 5.0\text{ s}$, and (b) $\alpha_z$ at $t_1$ and $t_2$.
\tcblower
\textbf{Solution:}
\begin{itemize}
  \item \textbf{(a)} $\alpha_{\text{av-}z} = \frac{\omega_{2z} - \omega_{1z}}{t_2 - t_1} = \frac{150 - 24}{5.0 - 2.0} = \eqbox{42\text{ rad/s}^2}$.
  \item \textbf{(b)} $\alpha_z = \frac{\diff\omega_z}{\diff t} = (12\text{ rad/s}^3)t \implies \alpha_{1z} = 12(2.0) = \eqbox{24\text{ rad/s}^2}$, $\alpha_{2z} = 12(5.0) = \eqbox{60\text{ rad/s}^2}$.
\end{itemize}
\end{example}

\begin{example}{9.3}{Rotation with constant angular acceleration\bk{p.\ 285}}
\textbf{Problem:} A Blu-ray disc has $\omega_{0z} = 27.5\text{ rad/s}$, $\theta_0 = 0$, and constant $\alpha_z = -10.0\text{ rad/s}^2$. At $t = 0.300\text{ s}$, find (a) $\omega_z$ and (b) the angle a reference line $PQ$ makes with the $+x$-axis.
\tcblower
\textbf{Solution:}
\begin{itemize}
  \item \textbf{(a)} $\omega_z = \omega_{0z} + \alpha_z t = 27.5 + (-10.0)(0.300) = \eqbox{24.5\text{ rad/s}}$.
  \item \textbf{(b)} $\theta = 0 + (27.5)(0.300) + \frac{1}{2}(-10.0)(0.300)^2 = \eqbox{7.80\text{ rad} = 1.24\text{ rev}}$. Subtracting $1\text{ rev} = 360^\circ$ leaves $(0.24\text{ rev})(360^\circ/\text{rev}) = \eqbox{87^\circ}$.
\end{itemize}
\end{example}

\begin{example}{9.4}{Throwing a discus\bk{p.\ 287}}
\textbf{Problem:} A discus moves in a circle of radius $r = 0.800\text{ m}$. At an instant when $\omega = 10.0\text{ rad/s}$ and $\alpha = 50.0\text{ rad/s}^2$, find $a_{\text{tan}}$, $a_{\text{rad}}$, and total acceleration magnitude $a$.
\tcblower
\textbf{Solution:}
\begin{align*}
  a_{\text{tan}} &= r\alpha = (0.800\text{ m})(50.0\text{ rad/s}^2) = \eqbox{40.0\text{ m/s}^2}\\
  a_{\text{rad}} &= \omega^2 r = (10.0\text{ rad/s})^2(0.800\text{ m}) = \eqbox{80.0\text{ m/s}^2}\\
  a &= \sqrt{a_{\text{tan}}^2 + a_{\text{rad}}^2} = \sqrt{(40.0)^2 + (80.0)^2} = \eqbox{89.4\text{ m/s}^2}
\end{align*}
\end{example}

\begin{example}{9.5}{Designing a propeller\bk{p.\ 288}}
\textbf{Problem:} An airplane propeller turns at $2400\text{ rpm}$ ($\omega = 251\text{ rad/s}$) while the plane flies forward at $v_{\text{plane}} = 75.0\text{ m/s}$. The tip speed through the air is $v_{\text{tip}} = \sqrt{v_{\text{plane}}^2 + r^2\omega^2} \le 270\text{ m/s}$. Find (a) $r_{\max}$ and (b) $a_{\text{rad}}$ of the tip.
\tcblower
\textbf{Solution:}
\begin{align*}
  \text{(a)}\quad r &= \frac{\sqrt{v_{\text{tip}}^2 - v_{\text{plane}}^2}}{\omega} = \frac{\sqrt{(270)^2 - (75.0)^2}}{251\text{ rad/s}} = \eqbox{1.03\text{ m}}\\
  \text{(b)}\quad a_{\text{rad}} &= \omega^2 r = (251\text{ rad/s})^2(1.03\text{ m}) = \eqbox{6.5\times 10^4\text{ m/s}^2 = 6600g}
\end{align*}
\end{example}

\begin{example}{9.6}{Moments of inertia for different rotation axes\bk{p.\ 290}}
\textbf{Problem:} Three point-like disks ($m_A = 0.30\text{ kg}$, $m_B = 0.10\text{ kg}$, $m_C = 0.20\text{ kg}$) are joined by massless struts ($r_{AB} = 0.50\text{ m}$, $r_{AC} = 0.40\text{ m}$, $r_{BC} = 0.30\text{ m}$, forming a right triangle with right angle at $C$). Find: (a) $I_A$ about an axis through $A$ perpendicular to the plane, (b) $I_{BC}$ about line $BC$, and (c) $K_A$ when $\omega = 4.0\text{ rad/s}$ about $A$.
\tcblower
\textbf{Solution:}
\begin{align*}
  \text{(a)}\quad I_A &= m_B r_{AB}^2 + m_C r_{AC}^2 = (0.10)(0.50)^2 + (0.20)(0.40)^2 = \eqbox{0.057\text{ kg}\cdot\text{m}^2}\\
  \text{(b)}\quad I_{BC} &= m_A r_{AC}^2 = (0.30\text{ kg})(0.40\text{ m})^2 = \eqbox{0.048\text{ kg}\cdot\text{m}^2}\\
  \text{(c)}\quad K_A &= \frac{1}{2}I_A\omega^2 = \frac{1}{2}(0.057)(4.0)^2 = \eqbox{0.46\text{ J}}
\end{align*}
\end{example}

\begin{example}{9.7}{An unwinding cable I\bk{p.\ 291}}
\textbf{Problem:} A light cable wrapped around a solid cylinder ($M = 50\text{ kg}$, $R = 0.060\text{ m}$) is pulled with constant force $F = 9.0\text{ N}$ over distance $s = 2.0\text{ m}$ starting from rest. Find the final angular speed $\omega$ and cable speed $v$.
\tcblower
\textbf{Solution:}
\begin{align*}
  W_{\text{other}} &= Fs = (9.0\text{ N})(2.0\text{ m}) = 18\text{ J}, \qquad I = \frac{1}{2}MR^2 = \frac{1}{2}(50)(0.060)^2 = 0.090\text{ kg}\cdot\text{m}^2\\
  \omega &= \sqrt{\frac{2W_{\text{other}}}{I}} = \sqrt{\frac{2(18\text{ J})}{0.090\text{ kg}\cdot\text{m}^2}} = \eqbox{20\text{ rad/s}}, \qquad v = R\omega = (0.060)(20) = \eqbox{1.2\text{ m/s}}
\end{align*}
\end{example}

\begin{example}{9.8}{An unwinding cable II\bk{p.\ 292}}
\textbf{Problem:} A block of mass $m$ hangs from a light cable wrapped around a solid cylinder of mass $M$ and radius $R$. Released from rest at height $h$, find the block's speed $v$ after falling distance $h$.
\tcblower
\textbf{Solution:} Conservation of mechanical energy $K_1 + U_1 = K_2 + U_2$ with $I = \frac{1}{2}MR^2$ and $\omega = v/R$:
\begin{equation*}
  0 + mgh = \frac{1}{2}mv^2 + \frac{1}{2}\!\left(\frac{1}{2}MR^2\right)\!\left(\frac{v}{R}\right)^2 = \frac{1}{2}\!\left(m + \frac{1}{2}M\right)v^2 \implies \eqbox{v = \sqrt{\frac{2gh}{1 + M/2m}}}
\end{equation*}
\end{example}

\begin{example}{9.9}{Using the parallel-axis theorem\bk{p.\ 294}}
\textbf{Problem:} A linkage part of mass $M = 3.6\text{ kg}$ has $I_P = 0.132\text{ kg}\cdot\text{m}^2$ about an axis $d = 0.15\text{ m}$ from its center of mass. Find $I_{\text{cm}}$.
\tcblower
\textbf{Solution:}
\begin{equation*}
  I_{\text{cm}} = I_P - Md^2 = 0.132 - (3.6)(0.15)^2 = \eqbox{0.051\text{ kg}\cdot\text{m}^2}
\end{equation*}
\end{example}

\begin{example}{Ch.\ 9}{Bridging Problem: A Rotating, Uniform Thin Rod\bk{p.\ 298, 307}}
\textbf{Problem:} A slender uniform rod of mass $M$ and length $L$ pivots about an axis through $O$ at distance $h$ from one end. (a) Integrate to find $I$ about $O$. (b) Starting from rest with constant angular acceleration $\alpha$, find the work $W$ done in time $t$. (c) Find the magnitude of the linear acceleration $a$ of the far end (distance $L-h$) at time $t$.
\tcblower
\textbf{Solution:}
\begin{itemize}
  \item \textbf{(a)} With $\diff m = \frac{M}{L}\dd x$ from $x = -h$ to $x = L-h$:
  \begin{equation*}
    I = \int_{-h}^{L-h} x^2\!\left(\frac{M}{L}\dd x\right) = \frac{M}{3L}\!\left[(L-h)^3 - (-h)^3\right] = \eqbox{\frac{1}{3}M(L^2 - 3Lh + 3h^2)}
  \end{equation*}
  \item \textbf{(b)} At time $t$, $\omega = \alpha t$, so $W = K = \frac{1}{2}I\omega^2 = \eqbox{\frac{1}{6}M(L^2 - 3Lh + 3h^2)\alpha^2 t^2}$.
  \item \textbf{(c)} At $r = L-h$, $a_{\text{tan}} = (L-h)\alpha$ and $a_{\text{rad}} = \omega^2 r = (L-h)\alpha^2 t^2$, so:
  \begin{equation*}
    a = \sqrt{a_{\text{tan}}^2 + a_{\text{rad}}^2} = \eqbox{(L-h)\alpha\sqrt{1 + \alpha^2 t^4}}
  \end{equation*}
\end{itemize}
\end{example}

% ============================================================================
% UNIQUE SOLVED PROBLEMS — MODULE 2
% ============================================================================
\section{Unique Solved Problems (Module 2)}

\begin{enumerate}[label=\textbf{Prob.\ 9.\arabic*},leftmargin=4.2em]
  \setcounter{enumi}{59}
  \item \textbf{[Prob.\ 9.60, p.\ 303] Non-Uniform Slender Rod ($\diff m/\diff x = \gamma x$):}
  \begin{itemize}
    \item \textbf{(a) Total mass:} $M = \int_0^L \gamma x\dd x = \eqbox{\frac{1}{2}\gamma L^2} \implies \gamma = \frac{2M}{L^2}$.
    \item \textbf{(b) Axis at light end ($x=0$):} $I_{\text{left}} = \int_0^L x^2(\gamma x\dd x) = \frac{1}{4}\gamma L^4 = \eqbox{\frac{1}{2}ML^2}$ (greater than $\frac{1}{3}ML^2$ because more mass is near $x=L$).
    \item \textbf{(c) Axis at heavy end ($x=L$):} $I_{\text{right}} = \int_0^L (L-x)^2(\gamma x\dd x) = \gamma L^4\!\left(\frac{1}{2} - \frac{2}{3} + \frac{1}{4}\right) = \frac{1}{12}\gamma L^4 = \eqbox{\frac{1}{6}ML^2}$ (smaller than $\frac{1}{3}ML^2$ because more mass is close to the axis at $x=L$).
  \end{itemize}

  \item \textbf{[Prob.\ 9.78, p.\ 304] Disk with Point Mass on Rim:}
  Uniform disk ($m, R$, $I_{\text{disk}} = \frac{1}{2}mR^2$) pivoted at center has a point mass $m$ glued to its rim at horizontal level ($y=0$) and is released from rest. When the point mass reaches the bottom ($y = -R$):
  \begin{equation*}
    -\Delta U = mgR = \frac{1}{2}(I_{\text{disk}} + I_{\text{point}})\omega^2 = \frac{1}{2}\!\left(\frac{1}{2}mR^2 + mR^2\right)\omega^2 = \frac{3}{4}mR^2\omega^2 \implies \eqbox{\omega = \sqrt{\frac{4g}{3R}}}
  \end{equation*}

  \item \textbf{[Prob.\ 9.79, p.\ 304] Right-Triangular Plate of Base $b$, Height $h$, and Mass $M$:}
  Divide the triangle into thin vertical strips parallel to side $h$ at distance $x\in [0,b]$ from side $h$. Strip height is $y(x) = h\!\left(1 - \frac{x}{b}\right)$ and area is $A = \frac{1}{2}bh$, so $\diff m = \frac{M}{A}y(x)\dd x = \frac{2M}{b}\!\left(1 - \frac{x}{b}\right)\dd x$:
  \begin{equation*}
    I_h = \int_0^b x^2\dd m = \frac{2M}{b}\int_0^b \left(x^2 - \frac{x^3}{b}\right)\dd x = \frac{2M}{b}\!\left(\frac{b^3}{3} - \frac{b^3}{4}\right) = \eqbox{\frac{1}{6}Mb^2}
  \end{equation*}

  \item \textbf{[Prob.\ 9.82, p.\ 305] Flexible Rope Unwrapping from a Cylinder:}
  One turn of rope (mass $m$, length $L = 2\pi R$) unwraps from cylinder ($M, R$, initial $\omega_0$) and hangs vertically tangent to the cylinder. Its center of mass drops by $y_{\text{cm}} = \frac{L}{2} = \pi R$:
  \begin{align*}
    \frac{1}{2}\!\left(\frac{1}{2}MR^2 + mR^2\right)\omega_0^2 + mg(\pi R) &= \frac{1}{2}\!\left(\frac{1}{2}MR^2 + mR^2\right)\omega^2\\
    \eqbox{\omega = \sqrt{\omega_0^2 + \frac{4\pi mg}{(M+2m)R}}}, \qquad v &= R\omega
  \end{align*}

  \item \textbf{[Prob.\ 9.83, p.\ 305] Block--Pulley--Table System with Friction:}
  Block $B$ ($m_B$) descends distance $d$, pulling block $A$ ($m_A$, friction $\mu_k$) across a horizontal table over a pulley ($I, R$).
  \begin{align*}
    K_1 + U_1 + W_{\text{fric}} = K_2 + U_2 &\implies 0 + m_B gd - \mu_k m_A gd = \frac{1}{2}(m_A + m_B)v^2 + \frac{1}{2}I\!\left(\frac{v}{R}\right)^2\\
    &\implies \eqbox{v = \sqrt{\frac{2(m_B - \mu_k m_A)gd}{m_A + m_B + I/R^2}}}
  \end{align*}

  \item \textbf{[Prob.\ 9.91, p.\ 305] Disk with Off-Center Circular Hole:}
  Full disk has mass $M_{\text{orig}}$ and radius $R$. A hole of radius $r_h = R/4$ (area $= \frac{1}{16}$ of full disk) is punched at distance $d = R/2$ from center $O$. Let $M$ be the mass of the original full disk: $m_{\text{hole}} = \frac{M}{16}$.
  \begin{itemize}
    \item \textbf{(a) Axis through original center $O$:} By parallel-axis theorem, $I_{\text{hole},O} = \frac{1}{2}m_{\text{hole}}r_h^2 + m_{\text{hole}}d^2 = \frac{1}{2}\!\left(\frac{M}{16}\right)\!\left(\frac{R}{4}\right)^2 + \left(\frac{M}{16}\right)\!\left(\frac{R}{2}\right)^2 = \frac{9}{512}MR^2$. Subtracting from $I_{\text{full},O} = \frac{1}{2}MR^2$:
    \begin{equation*}
      I_O = \frac{1}{2}MR^2 - \frac{9}{512}MR^2 = \eqbox{\frac{247}{512}MR^2 \approx 0.482MR^2}
    \end{equation*}
    \item \textbf{(b) Axis through center of hole $C$:} $I_{\text{full},C} = \frac{1}{2}MR^2 + M\!\left(\frac{R}{2}\right)^2 = \frac{3}{4}MR^2$, and $I_{\text{hole},C} = \frac{1}{2}\!\left(\frac{M}{16}\right)\!\left(\frac{R}{4}\right)^2 = \frac{1}{512}MR^2$:
    \begin{equation*}
      I_C = \frac{3}{4}MR^2 - \frac{1}{512}MR^2 = \eqbox{\frac{383}{512}MR^2 \approx 0.748MR^2}
    \end{equation*}
  \end{itemize}

  \item \textbf{[Prob.\ 9.96, p.\ 306] Thin Rod Bent into a Square of Side $a$ (Total Mass $M$, $L=4a$):}
  Each side has mass $m_s = M/4$, length $a$, and distance $d = a/2$ from the square's center:
  \begin{equation*}
    I = 4\!\left(\frac{1}{12}m_s a^2 + m_s d^2\right) = 4\!\left(\frac{M}{4}\right)\!\left(\frac{a^2}{12} + \frac{a^2}{4}\right) = \eqbox{\frac{1}{3}Ma^2}
  \end{equation*}

  \item \textbf{[Prob.\ 9.97, p.\ 306] Cylinder with Radially Increasing Density $\rho(r) = \alpha r$:}
  \begin{itemize}
    \item \textbf{(a)} Mass $M = \int_0^R (\alpha r)(2\pi r L\dd r) = \frac{2}{3}\pi\alpha L R^3$. Moment of inertia:
    \begin{equation*}
      I = \int_0^R r^2(\alpha r)(2\pi r L\dd r) = \frac{2}{5}\pi\alpha L R^5 = \left(\frac{2}{3}\pi\alpha L R^3\right)\!\left(\frac{3}{5}R^2\right) = \eqbox{\frac{3}{5}MR^2 = 0.60MR^2}
    \end{equation*}
    \item \textbf{(b)} This is larger than $\frac{1}{2}MR^2$ for a uniform cylinder because more mass is concentrated near the outer rim.
  \end{itemize}

  \item \textbf{[Challenge Problem 9.100, p.\ 306] Moment of Inertia of a Uniform Solid Cone:}
  Slice the cone (mass $M$, base radius $R$, height $h$, volume $V = \frac{1}{3}\pi R^2 h$) into thin disks at distance $y\in [0,h]$ from the apex, where disk radius is $r(y) = \frac{R}{h}y$ and $\diff m = \rho\pi r^2\dd y$:
  \begin{equation*}
    I = \int_0^h \frac{1}{2}r^2\dd m = \frac{1}{2}\pi\rho\frac{R^4}{h^4}\int_0^h y^4\dd y = \frac{1}{10}\pi\rho R^4 h = \frac{3}{10}\!\left(\frac{1}{3}\pi\rho R^2 h\right)R^2 = \eqbox{\frac{3}{10}MR^2}
  \end{equation*}

  \item \textbf{[Challenge Problem 9.101, p.\ 306--307] Compact Disc Spiral Track Kinematics:}
  Given $r(\theta) = r_0 + \beta\theta$ and constant linear speed $v$:
  \begin{itemize}
    \item \textbf{(a)} Integrate $\diff s = r(\theta)\dd\theta$ from $0$ to $\theta$:
    \begin{equation*}
      s(\theta) = \int_0^\theta (r_0 + \beta\theta')\dd\theta' = \eqbox{r_0\theta + \frac{1}{2}\beta\theta^2}
    \end{equation*}
    \item \textbf{(b)} Set $s(\theta) = vt \implies \frac{1}{2}\beta\theta^2 + r_0\theta - vt = 0$. Taking the positive root (since $\theta \ge 0$ for $t \ge 0$):
    \begin{equation*}
      \eqbox{\theta(t) = \frac{\sqrt{r_0^2 + 2\beta vt} - r_0}{\beta}}
    \end{equation*}
    \item \textbf{(c)} Differentiate $\theta(t)$ with respect to $t$:
    \begin{equation*}
      \eqbox{\omega_z(t) = \frac{\diff\theta}{\diff t} = \frac{v}{\sqrt{r_0^2 + 2\beta vt}}, \qquad \alpha_z(t) = \frac{\diff\omega_z}{\diff t} = -\frac{\beta v^2}{(r_0^2 + 2\beta vt)^{3/2}}} \quad(\text{not constant})
    \end{equation*}
  \end{itemize}
\end{enumerate}

% ============================================================================
% FORMULA SHEET — MODULE 2
% ============================================================================
\formulasheet{Module 2 Formula Sheet: Rotation of Rigid Bodies}

\begin{fs}{1. Rotational and Linear Kinematics}
\begin{align*}
  s &= r\theta, \qquad \omega_z = \frac{\diff\theta}{\diff t}, \qquad \alpha_z = \frac{\diff\omega_z}{\diff t} = \frac{\diff^2\theta}{\diff t^2} = \omega_z\frac{\diff\omega_z}{\diff\theta}\\
  \text{Constant }\alpha_z:\quad & \omega_z = \omega_{0z} + \alpha_z t, \qquad \theta - \theta_0 = \omega_{0z}t + \frac{1}{2}\alpha_z t^2 = \frac{1}{2}(\omega_{0z}+\omega_z)t, \qquad \omega_z^2 = \omega_{0z}^2 + 2\alpha_z(\theta-\theta_0)\\
  \text{Linear--Angular:}\quad & v = r\omega, \qquad a_{\text{tan}} = r\alpha, \qquad a_{\text{rad}} = \frac{v^2}{r} = \omega^2 r, \qquad a = \sqrt{a_{\text{tan}}^2 + a_{\text{rad}}^2}
\end{align*}
\end{fs}

\begin{fs}{2. Rotational Energy, Parallel/Perpendicular-Axis Theorems, and Standard Moments of Inertia}
\begin{align*}
  I &= \sum_i m_i r_i^2 = \int r^2\dd m, \qquad K_{\text{rot}} = \frac{1}{2}I\omega^2, \qquad U_{\text{grav}} = Mgy_{\text{cm}}, \qquad I_P = I_{\text{cm}} + Md^2, \qquad I_z = I_x + I_y\ (\text{lamina})\\
  I_{\text{rod,cm}} &= \frac{1}{12}ML^2, \quad I_{\text{rod,end}} = \frac{1}{3}ML^2, \quad I_{\text{cyl}} = \frac{1}{2}M(R_1^2+R_2^2), \quad I_{\text{disk}} = \frac{1}{2}MR^2, \quad I_{\text{hoop}} = MR^2, \quad I_{\text{sphere}} = \frac{2}{5}MR^2, \quad I_{\text{shell}} = \frac{2}{3}MR^2
\end{align*}
\end{fs}

% ============================================================================
% MODULE 3: DYNAMICS OF ROTATIONAL MOTION AND ANGULAR MOMENTUM (PDF 3)
% ============================================================================
\part{Dynamics of Rotational Motion and Angular Momentum}

\chapter{Torque, Rotational Dynamics, and Rolling Motion}
\begin{syll}
\textbf{Syllabus Coverage (PDF 3):} Young \& Freedman Chapter 10 (Sections 10.1--10.4) and Analytical Mechanics Chapter 7 (Section 7.2.1). Torque vector, rotational analog of Newton's second law, combined translation and rotation (Koenig's theorem), rolling without slipping, and rotational work and power.
\end{syll}

\section{Torque and Rotational Equilibrium\bk{10.1, 7.2.1}}

\begin{goal}
Define the torque vector $\bm{\tau} = \vect{r}\times\vect{F}$ and prove that if a body is in translational and rotational equilibrium about one point $O$, the net torque about any other point $O'$ is also zero.
\end{goal}

\begin{pic}
\begin{tikzpicture}[scale=0.95]
  \fill[ink] (0,0) circle (2.5pt) node[below left,font=\small]{$O$};
  \coordinate (P) at (2.8,0.8);
  \draw[vec,teal] (0,0) -- (P) node[midway,below,font=\small]{$\vect{r}$};
  \draw[dashed,slate] (P) -- ++(16.0:1.2);
  \draw[vec,rose,ultra thick] (P) -- ++(115:1.6) node[above,font=\small]{$\vect{F}$};
  \draw[dashed,amber,thick] (0,0) -- ($(P)!(0,0)!($(P)+(115:1.6)$)$) node[midway,above left,font=\footnotesize]{$l = r\sin\phi$};
  \draw[slate] ($(P)+(16.0:0.55)$) arc (16.0:115:0.55) node[midway,right,font=\footnotesize]{$\phi$};
\end{tikzpicture}
\figcap{Torque $\bm{\tau} = \vect{r}\times\vect{F}$ has magnitude $\tau = Fl = rF\sin\phi = F_{\text{tan}}r$ and points out of the page.}
\end{pic}

\step{Step 1: Define torque magnitude and vector.}
\begin{equation}\label{eq:torque_def}
  \eqbox{\bm{\tau} = \vect{r}\times\vect{F} \quad (\text{or }\vect{N} = \vect{r}\times\vect{F})}, \qquad \eqbox{\tau = Fl = rF\sin\phi = F_{\text{tan}}r}
\end{equation}

\step{Step 2: Origin independence of equilibrium (Worked Example 7.1, p.\ 178--179).}
Let $\vect{r}_{iO}$ be the position vector of force $\vect{F}_i$ from $O$, and let $O'$ be displaced by $\vect{d}$ from $O$, so $\vect{r}_{iO} = \vect{d} + \vect{r}_{iO'}$ ($\vect{r}_{iO'} = \vect{r}_{iO} - \vect{d}$):
\begin{equation}\label{eq:torque_shift}
  \sum\vect{N}_{iO'} = \sum \vect{r}_{iO'}\times\vect{F}_i = \sum(\vect{r}_{iO} - \vect{d})\times\vect{F}_i = \sum\vect{N}_{iO} - \vect{d}\times\sum\vect{F}_i
\end{equation}
If the body is in equilibrium ($\sum\vect{F}_i = \vect{0}$ and $\sum\vect{N}_{iO} = \vect{0}$), then $\eqbox{\sum\vect{N}_{iO'} = \vect{0}}$ about \textbf{every} point $O'$.

\section{Torque and Angular Acceleration for a Rigid Body\bk{10.2}}

\begin{goal}
Derive $\sum\tau_z = I\alpha_z$ for a rigid body rotating about a fixed $z$-axis.
\end{goal}

\step{Step 1: Apply Newton's second law to the tangential motion of particle $i$.}
Particle $i$ (mass $m_i$, distance $r_i$ from the $z$-axis) has tangential acceleration $a_{i,\text{tan}} = r_i\alpha_z$:
\begin{align}
  F_{i,\text{tan}} &= m_i a_{i,\text{tan}} = m_i r_i\alpha_z &&\why{multiply both sides by $r_i$}\\
  \tau_{iz} &= F_{i,\text{tan}}r_i = m_i r_i^2\alpha_z = I_i\alpha_z &&\why{torque on particle $i$}
\end{align}

\step{Step 2: Sum over all particles in the rigid body.}
Internal forces between any pair of particles act along the line joining them with equal lever arms and opposite directions, so internal torques cancel in pairs:
\begin{equation}\label{eq:tau_Ialpha}
  \sum\tau_{iz} = \left(\sum_i m_i r_i^2\right)\alpha_z \implies \eqbox{\sum\tau_z = I\alpha_z}
\end{equation}

\section{Rigid-Body Rotation About a Moving Axis and Rolling Motion\bk{10.3}}

\begin{goal}
Prove that the total kinetic energy of a translating and rotating rigid body is $K = \frac{1}{2}Mv_{\text{cm}}^2 + \frac{1}{2}I_{\text{cm}}\omega^2$, and derive the dynamics of rolling without slipping.
\end{goal}

\step{Step 1: Prove Koenig's Kinetic Energy Theorem.}
Write the velocity of particle $i$ as $\vect{v}_i = \vect{v}_{\text{cm}} + \vect{v}_i'$, where $\vect{v}_i'$ is the velocity relative to the center of mass:
\begin{align}
  K &= \sum_i \frac{1}{2}m_i(\vect{v}_{\text{cm}} + \vect{v}_i')\cdot(\vect{v}_{\text{cm}} + \vect{v}_i')\nonumber\\
    &= \frac{1}{2}\!\left(\sum_i m_i\right)v_{\text{cm}}^2 + \vect{v}_{\text{cm}}\cdot\!\underbrace{\left(\sum_i m_i\vect{v}_i'\right)}_{\vect{0}\text{ (by CM def.)}} + \sum_i \frac{1}{2}m_i v_i'^2 = \eqbox{\frac{1}{2}Mv_{\text{cm}}^2 + \frac{1}{2}I_{\text{cm}}\omega^2}\label{eq:K_koenig}
\end{align}

\step{Step 2: Rolling without slipping condition and dynamics.}
For a wheel of radius $R$ rolling without slipping on a stationary surface, the contact point is instantaneously at rest:
\begin{equation}\label{eq:roll_cond}
  \eqbox{v_{\text{cm}} = R\omega, \qquad a_{\text{cm}} = R\alpha}
\end{equation}
If the axis through the center of mass is a symmetry axis that does not change direction, the coupled equations of motion are:
\begin{equation}\label{eq:trans_rot_dyn}
  \eqbox{\sum\vect{F}_{\text{ext}} = M\vect{a}_{\text{cm}}, \qquad \sum\tau_z = I_{\text{cm}}\alpha_z}
\end{equation}

\step{Step 3: General formulas for a round body ($I_{\text{cm}} = cMR^2$) rolling without slipping down an incline $\beta$.}
Along the incline: $Mg\sin\beta - f_s = Ma_{\text{cm-}x}$. About the center of mass: $f_s R = I_{\text{cm}}\alpha_z = (cMR^2)\!\left(\frac{a_{\text{cm-}x}}{R}\right) \implies f_s = cMa_{\text{cm-}x}$. Substitute $f_s$ into the force equation:
\begin{equation}\label{eq:roll_incline}
  \eqbox{a_{\text{cm-}x} = \frac{g\sin\beta}{1 + c}, \qquad f_s = \frac{c}{1 + c}Mg\sin\beta, \qquad \mu_{s,\min} = \frac{f_s}{Mg\cos\beta} = \frac{c}{1 + c}\tan\beta, \qquad v_{\text{cm}} = \sqrt{\frac{2gh}{1 + c}}}
\end{equation}

\section{Work and Power in Rotational Motion\bk{10.4}}

\begin{goal}
Derive the work done by a torque, the rotational work--energy theorem, and rotational power $P = \tau_z\omega_z$.
\end{goal}

\step{Step 1: Infinitesimal work $\diff W$ of a tangential force.}
\begin{equation}\label{eq:work_rot}
  \diff W = F_{\text{tan}}\dd s = F_{\text{tan}}(R\dd\theta) = \tau_z\dd\theta \implies \eqbox{W = \int_{\theta_1}^{\theta_2}\tau_z\dd\theta} \quad(W = \tau_z\Delta\theta\text{ for constant }\tau_z)
\end{equation}

\step{Step 2: Rotational Work--Energy Theorem.}
Substitute $\tau_z\dd\theta = (I\alpha_z)\dd\theta = I\frac{\diff\omega_z}{\diff t}\dd\theta = I\omega_z\dd\omega_z$:
\begin{equation}\label{eq:rot_WE}
  \eqbox{W_{\text{tot}} = \int_{\omega_1}^{\omega_2} I\omega_z\dd\omega_z = \frac{1}{2}I\omega_2^2 - \frac{1}{2}I\omega_1^2}
\end{equation}

\step{Step 3: Rotational Power.}
Divide $\diff W = \tau_z\dd\theta$ by $\diff t$:
\begin{equation}\label{eq:rot_power}
  \eqbox{P = \frac{\diff W}{\diff t} = \tau_z\frac{\diff\theta}{\diff t} = \tau_z\omega_z}
\end{equation}

% ============================================================================
% CHAPTER 6: ANGULAR MOMENTUM, GYROSCOPES, AND AXIAL VECTORS
% ============================================================================
\chapter{Angular Momentum, Gyroscopes, and Axial Vectors}
\begin{syll}
\textbf{Syllabus Coverage (PDF 3):} Young \& Freedman Chapter 10 (Sections 10.5--10.7) and Analytical Mechanics Chapter 7 (Sections 7.1--7.8). Angular momentum of a particle and system of particles, center-of-mass decomposition $\vect{L} = \vect{L}_c + \vect{L}'$, BAC--CAB derivation of $\vect{L} = I\bm{\omega}$, conservation of angular momentum, gyroscope precession, and axial vectors (pseudovectors).
\end{syll}

\section{Angular Momentum of a Particle and a System of Particles\bk{10.5, 7.1--7.3}}

\begin{goal}
Derive $\frac{\diff\vect{l}}{\diff t} = \bm{\tau}$ for a single particle, $\frac{\diff\vect{L}}{\diff t} = \vect{N}_{\text{tot}}^{(e)}$ for an $N$-particle system, and the center-of-mass decomposition $\vect{L} = \vect{L}_c + \vect{L}'$.
\end{goal}

\step{Step 1: Single particle angular momentum $\vect{l}$ and its time derivative.}
\begin{align}
  \vect{l} &= \vect{r}\times\vect{p} = \vect{r}\times m\vect{v} \qquad (l = mvr\sin\phi = mvl)\label{eq:l_particle}\\
  \frac{\diff\vect{l}}{\diff t} &= \frac{\diff\vect{r}}{\diff t}\times m\vect{v} + \vect{r}\times m\frac{\diff\vect{v}}{\diff t} = \underbrace{\vect{v}\times m\vect{v}}_{\vect{0}} + \vect{r}\times\vect{F} = \eqbox{\bm{\tau} = \vect{N}}\label{eq:dl_dt}
\end{align}

\step{Step 2: System of $N$ particles and pairwise cancellation of internal torques (Sec.\ 7.3).}
For particle $i$, $\frac{\diff\vect{l}_i}{\diff t} = \vect{r}_i\times\vect{F}_i^{(e)} + \sum_{j\neq i}\vect{r}_i\times\vect{F}_{ij}$. Summing over all $i = 1,\dots,N$ with $\vect{L} = \sum_i\vect{l}_i$:
\begin{equation}
  \frac{\diff\vect{L}}{\diff t} = \sum_i \vect{r}_i\times\vect{F}_i^{(e)} + \sum_i\sum_{j\neq i}\vect{r}_i\times\vect{F}_{ij}
\end{equation}
For each pair $(i,j)$, Newton's third law gives $\vect{F}_{ji} = -\vect{F}_{ij}$, so the pair contributes:
\begin{equation}
  \vect{r}_i\times\vect{F}_{ij} + \vect{r}_j\times\vect{F}_{ji} = (\vect{r}_i - \vect{r}_j)\times\vect{F}_{ij} = \vect{r}_{ij}\times\vect{F}_{ij} = \vect{0}
\end{equation}
because central internal forces act along the line $\vect{r}_{ij}$ joining the particles ($\vect{r}_{ij}\parallel\vect{F}_{ij}$). Therefore:
\begin{equation}\label{eq:dL_dt_sys}
  \eqbox{\frac{\diff\vect{L}}{\diff t} = \sum\bm{\tau}_{\text{ext}} = \vect{N}_{\text{tot}}^{(e)}} \implies \text{If }\vect{N}_{\text{tot}}^{(e)} = \vect{0},\text{ then }\eqbox{\vect{L} = \text{constant}}
\end{equation}

\step{Step 3: Angular Momentum Relative to the Center of Mass (Sec.\ 7.3.1, p.\ 184).}
Write $\vect{r}_i = \vect{r}_c + \vect{r}_i'$ and $\dot{\vect{r}}_i = \dot{\vect{r}}_c + \dot{\vect{r}}_i'$, where $\sum_i m_i\vect{r}_i' = \vect{0}$ and $\sum_i m_i\dot{\vect{r}}_i' = \vect{0}$:
\begin{align}
  \vect{L} &= \sum_i m_i(\vect{r}_c + \vect{r}_i')\times(\dot{\vect{r}}_c + \dot{\vect{r}}_i')\nonumber\\
           &= \left(\sum_i m_i\right)\vect{r}_c\times\dot{\vect{r}}_c + \vect{r}_c\times\!\underbrace{\left(\sum_i m_i\dot{\vect{r}}_i'\right)}_{\vect{0}} - \dot{\vect{r}}_c\times\!\underbrace{\left(\sum_i m_i\vect{r}_i'\right)}_{\vect{0}} + \sum_i m_i(\vect{r}_i'\times\dot{\vect{r}}_i')\nonumber\\
  \vect{L} &= \eqbox{M(\vect{r}_c\times\dot{\vect{r}}_c) + \sum_i m_i(\vect{r}_i'\times\dot{\vect{r}}_i') = \vect{L}_c + \vect{L}'}\label{eq:L_cm_decomp}
\end{align}

\section{Angular Momentum of a Rigid Body (BAC--CAB Derivation)\bk{10.5, 7.4--7.5}}

\begin{goal}
Derive $\vect{l}_i = m_i r_i^2\omega\sin^2\theta_i\khat - m_i r_i^2\omega\sin\theta_i\cos\theta_i\uv{e}_i$ using the BAC--CAB identity, and show that $\vect{L} = I\bm{\omega}$ for a symmetric rigid body.
\end{goal}

\begin{pic}
\begin{tikzpicture}[scale=0.95]
  \draw[ax] (0,-1.4) -- (0,2.4) node[above,font=\small]{$\bm{\omega} = \omega\khat$};
  \draw[thick,teal,dashed] (0,1.2) ellipse (1.6cm and 0.38cm);
  \coordinate (M1) at (1.6,1.2);
  \coordinate (M2) at (-1.6,1.2);
  \fill[teal] (M1) circle (3pt) node[right,font=\footnotesize]{$m_i$};
  \fill[amber] (M2) circle (3pt) node[left,font=\footnotesize]{$m_j = m_i$};
  \draw[vec,ink] (0,0) -- (M1) node[midway,below right,font=\footnotesize]{$\vect{r}_i$};
  \draw[vec,ink] (0,0) -- (M2) node[midway,below left,font=\footnotesize]{$\vect{r}_j$};
  \draw[vec,rose,thick] (0,0) -- (-0.75,1.0) node[left,font=\footnotesize]{$\vect{l}_i$};
  \draw[vec,rose,thick] (0,0) -- (0.75,1.0) node[right,font=\footnotesize]{$\vect{l}_j$};
  \draw[vec,plum,ultra thick] (0,0) -- (0,2.0) node[right,font=\footnotesize]{$\vect{l}_i + \vect{l}_j$};
  \draw[slate] (0,0.55) arc (90:36.9:0.55) node[midway,above right,font=\scriptsize]{$\theta_i$};
\end{tikzpicture}
\figcap{For a symmetric rigid body, the radial components of $\vect{l}_i$ and $\vect{l}_j$ cancel, leaving $\vect{L} = I\bm{\omega}$ along the rotation axis.}
\end{pic}

\step{Step 1: Expand $\vect{l}_i = m_i\vect{r}_i\times(\bm{\omega}\times\vect{r}_i)$ using the BAC--CAB rule (\tref{tool:cross}) (Exercise 7.8, p.\ 186--187).}
Since $\vect{v}_i = \bm{\omega}\times\vect{r}_{i\perp} = \bm{\omega}\times\vect{r}_i$:
\begin{equation}\label{eq:li_baccab}
  \vect{l}_i = m_i\vect{r}_i\times(\bm{\omega}\times\vect{r}_i) = m_i\!\left[r_i^2\bm{\omega} - (\vect{r}_i\cdot\bm{\omega})\vect{r}_i\right]
\end{equation}

\step{Step 2: Resolve $\vect{r}_i$ along $\khat$ and perpendicular unit vector $\uv{e}_i$.}
Let $\bm{\omega} = \omega\khat$ and $\vect{r}_i = r_i\cos\theta_i\khat + r_i\sin\theta_i\uv{e}_i$, so $\vect{r}_i\cdot\bm{\omega} = r_i\omega\cos\theta_i$. Substitute into Eq.\ \eqref{eq:li_baccab}:
\begin{align}
  \vect{l}_i &= m_i\!\left[r_i^2\omega\khat - (r_i\omega\cos\theta_i)(r_i\cos\theta_i\khat + r_i\sin\theta_i\uv{e}_i)\right]\nonumber\\
             &= m_i r_i^2\omega(1 - \cos^2\theta_i)\khat - m_i r_i^2\omega\sin\theta_i\cos\theta_i\uv{e}_i\nonumber\\
  \vect{l}_i &= \eqbox{m_i r_i^2\omega\sin^2\theta_i\khat - m_i r_i^2\omega\sin\theta_i\cos\theta_i\uv{e}_i}\label{eq:li_components}
\end{align}

\step{Step 3: Sum over all particles for a symmetric rigid body.}
For a body symmetric about the rotation axis, every particle $i$ at $+\uv{e}_i$ has an identical partner $j$ at $-\uv{e}_i$, so the $\uv{e}_i$ terms cancel in pairs. With $r_{i\perp} = r_i\sin\theta_i$:
\begin{equation}\label{eq:L_Iomega}
  \vect{L} = \left(\sum_i m_i r_i^2\sin^2\theta_i\right)\omega\khat = \left(\sum_i m_i r_{i\perp}^2\right)\omega\khat \implies \eqbox{\vect{L} = I\bm{\omega}}
\end{equation}
When $I$ can vary with zero net external torque ($\sum\tau_{\text{ext},z} = 0$), angular momentum is conserved:
\begin{equation}\label{eq:L_conserv}
  \eqbox{I_1\omega_{1z} = I_2\omega_{2z}}
\end{equation}

\section{Gyroscopes and Precession\bk{10.7, 7.6}}

\begin{goal}
Derive the precession angular speed $\Omega_p = \frac{Mgr}{I\omega}$ for both a horizontal gyroscope and a gyroscope tilted at an arbitrary angle $\theta$ to the vertical.
\end{goal}

\begin{pic}
\begin{tikzpicture}[scale=0.95]
  % Left: Side view tilted at angle theta
  \draw[ax] (0,0) -- (0,2.5) node[above,font=\footnotesize]{$z$};
  \draw[vec,teal,ultra thick] (0,0) -- (1.4,2.0) node[right,font=\footnotesize]{$\vect{L}$};
  \draw[dashed,slate] (0,2.0) -- (1.4,2.0) node[midway,above,font=\scriptsize]{$L\sin\theta$};
  \draw[slate] (0,0.7) arc (90:55:0.7) node[midway,above right,font=\scriptsize]{$\theta$};
  \node[font=\footnotesize,slate] at (0.7,-0.4) {Side view};
  % Right: Top view of horizontal circle
  \begin{scope}[xshift=4.8cm,yshift=1.1cm]
    \draw[thick,teal,dashed] (0,0) circle (1.2cm);
    \draw[vec,teal,thick] (0,0) -- (-10:1.2) node[midway,below,font=\scriptsize]{$L\sin\theta$};
    \draw[vec,teal,thick] (0,0) -- (25:1.2);
    \draw[vec,rose,thick] (-10:1.2) -- (25:1.2) node[midway,right,font=\scriptsize]{$\Delta\vect{L} = \bm{\tau}\dd t$};
    \draw[slate] (-10:0.45) arc (-10:25:0.45) node[midway,right,font=\scriptsize]{$\Delta\phi$};
    \node[font=\footnotesize,slate] at (0,-1.5) {Top view};
  \end{scope}
\end{tikzpicture}
\figcap{Precession of a gyroscope tilted at angle $\theta$ from the vertical: $\Delta L = (L\sin\theta)\Delta\phi$.}
\end{pic}

\step{Step 1: Torque of gravity about the pivot $Q$.}
Let the center of mass be at distance $r$ (or $R$) along the axle at angle $\theta$ from the vertical $z$-axis:
\begin{equation}
  \bm{\tau} = \vect{r}\times M\vect{g} \implies \tau = Mgr\sin\theta = wr\sin\theta \quad(\text{directed horizontally in the }\uv{\phi}\text{-direction})
\end{equation}

\step{Step 2: Relate $\diff L = \tau\dd t$ to the precession angle $\diff\phi$.}
The horizontal radius of the circle traced by the tip of $\vect{L}$ is $L_{\text{horiz}} = L\sin\theta$. In time $\diff t$, the tip of $\vect{L}$ sweeps through azimuthal angle $\diff\phi$:
\begin{align}
  \diff L &= (L\sin\theta)\dd\phi \implies \frac{\diff L}{\diff t} = L\sin\theta\frac{\diff\phi}{\diff t} = L\sin\theta\,\Omega_p\\
  Mgr\sin\theta &= (I\omega)\sin\theta\,\Omega_p \implies \eqbox{\Omega_p = \frac{\diff\phi}{\diff t} = \frac{Mgr}{I\omega} = \frac{wr}{I\omega}}\label{eq:precession}
\end{align}
Notice that $\sin\theta$ cancels: $\Omega_p$ is independent of the tilt angle $\theta$ (for fast spin $\omega \gg \Omega_p$).

\section{Polar Vectors versus Axial Vectors (Pseudovectors)\bk{7.7}}

\begin{goal}
Distinguish polar (true) vectors from axial vectors (pseudovectors) under an improper coordinate transformation (mirror reflection / parity inversion).
\end{goal}

\begin{itemize}
  \item \textbf{Polar (True) Vectors ($\vect{r}, \vect{v}, \vect{a}, \vect{p}, \vect{F}$):} Under space inversion ($\vect{r} \to -\vect{r}$), all polar vectors reverse sign: $\vect{v} \to -\vect{v}$, $\vect{p} \to -\vect{p}$, $\vect{F} \to -\vect{F}$.
  \item \textbf{Axial Vectors / Pseudovectors ($\bm{\omega}, \vect{l} = \vect{r}\times\vect{p}, \bm{\tau} = \vect{r}\times\vect{F}$):} Because an axial vector is the cross product of two polar vectors, under inversion ($\vect{r}\to -\vect{r}$, $\vect{p}\to -\vect{p}$) it \textbf{does not} change sign:
  \begin{equation}\label{eq:axial_parity}
    \eqbox{\vect{l}' = (-\vect{r})\times(-\vect{p}) = +(\vect{r}\times\vect{p}) = +\vect{l}, \qquad \bm{\tau}' = (-\vect{r})\times(-\vect{F}) = +\bm{\tau}}
  \end{equation}
\end{itemize}

% ============================================================================
% SOLVED BOOK EXAMPLES — MODULE 3
% ============================================================================
\section{Solved Book Examples (Module 3)}

\begin{example}{10.1}{Applying a torque\bk{p.\ 310}}
\textbf{Problem:} A plumber applies a downward force $F = 900\text{ N}$ at the end of a wrench handle of length $r = 0.80\text{ m}$ inclined at $19^\circ$ above the horizontal ($\phi = 90^\circ + 19^\circ = 109^\circ$). Find the magnitude and direction of the torque about the center of the fitting $O$.
\tcblower
\textbf{Solution:}
\begin{align*}
  l &= r\sin 109^\circ = (0.80\text{ m})\cos 19^\circ = 0.76\text{ m}\\
  \tau &= Fl = rF\sin 109^\circ = (0.80\text{ m})(900\text{ N})\sin 109^\circ = \eqbox{680\text{ N}\cdot\text{m}\text{ (out of page, counterclockwise)}}
\end{align*}
\end{example}

\begin{example}{10.2}{An unwinding cable I (Dynamics)\bk{p.\ 313}}
\textbf{Problem:} For the solid cylinder ($M = 50\text{ kg}$, $R = 0.060\text{ m}$) pulled with force $F = 9.0\text{ N}$ in Example 9.7, find the angular acceleration $\alpha_z$ of the cylinder and the linear acceleration $a_{\text{tan}}$ of the cable.
\tcblower
\textbf{Solution:}
\begin{align*}
  \alpha_z &= \frac{\tau_z}{I} = \frac{FR}{\frac{1}{2}MR^2} = \frac{2F}{MR} = \frac{2(9.0\text{ N})}{(50\text{ kg})(0.060\text{ m})} = \eqbox{6.0\text{ rad/s}^2}\\
  a_{\text{tan}} &= R\alpha_z = (0.060\text{ m})(6.0\text{ rad/s}^2) = \eqbox{0.36\text{ m/s}^2}
\end{align*}
\end{example}

\begin{example}{10.3}{An unwinding cable II (Dynamics)\bk{p.\ 314}}
\textbf{Problem:} For the falling block ($m$) and solid cylinder ($M, R$) of Example 9.8, find the downward acceleration $a_y$ of the block and the cable tension $T$.
\tcblower
\textbf{Solution:} For the block: $mg - T = ma_y$. For the cylinder: $RT = I\alpha_z = \left(\frac{1}{2}MR^2\right)\!\left(\frac{a_y}{R}\right) \implies T = \frac{1}{2}Ma_y$. Substitute $T$ into the block equation:
\begin{equation*}
  mg - \frac{1}{2}Ma_y = ma_y \implies \eqbox{a_y = \frac{g}{1 + M/2m}}, \qquad \eqbox{T = \frac{1}{2}Ma_y = \frac{mg}{1 + 2m/M}}
\end{equation*}
\end{example}

\begin{example}{10.4}{Speed of a primitive yo-yo\bk{p.\ 317}}
\textbf{Problem:} A solid cylinder ($M, R$, $I_{\text{cm}} = \frac{1}{2}MR^2$) unwinds from rest down a massless string held fixed at the top. Find $v_{\text{cm}}$ after descending distance $h$.
\tcblower
\textbf{Solution:} Using Eq.\ \eqref{eq:K_koenig} with $\omega = v_{\text{cm}}/R$:
\begin{equation*}
  K_2 = \frac{1}{2}Mv_{\text{cm}}^2 + \frac{1}{2}\!\left(\frac{1}{2}MR^2\right)\!\left(\frac{v_{\text{cm}}}{R}\right)^2 = \frac{3}{4}Mv_{\text{cm}}^2 = Mgh \implies \eqbox{v_{\text{cm}} = \sqrt{\frac{4}{3}gh}}
\end{equation*}
\end{example}

\begin{example}{10.5}{Race of the rolling bodies\bk{p.\ 317}}
\textbf{Problem:} Various round bodies with $I_{\text{cm}} = cMR^2$ roll without slipping from rest down an incline of vertical height $h$. Find $v_{\text{cm}}$ at the bottom and rank the bodies.
\tcblower
\textbf{Solution:} From energy conservation $Mgh = \frac{1}{2}Mv_{\text{cm}}^2 + \frac{1}{2}(cMR^2)(v_{\text{cm}}/R)^2 = \frac{1}{2}(1+c)Mv_{\text{cm}}^2$:
\begin{equation*}
  \eqbox{v_{\text{cm}} = \sqrt{\frac{2gh}{1+c}}} \implies \text{Order of finish: Solid sphere }\left(c=\frac{2}{5}\right),\text{ Solid cylinder }\left(c=\frac{1}{2}\right),\text{ Hollow sphere }\left(c=\frac{2}{3}\right),\text{ Hoop }(c=1).
\end{equation*}
\end{example}

\begin{example}{10.6}{Acceleration of a primitive yo-yo\bk{p.\ 319}}
\textbf{Problem:} Find the downward acceleration $a_{\text{cm-}y}$ and string tension $T$ for the yo-yo of Example 10.4.
\tcblower
\textbf{Solution:} Translation: $Mg - T = Ma_{\text{cm-}y}$. Rotation: $TR = \left(\frac{1}{2}MR^2\right)\!\left(\frac{a_{\text{cm-}y}}{R}\right) \implies T = \frac{1}{2}Ma_{\text{cm-}y}$:
\begin{equation*}
  Mg - \frac{1}{2}Ma_{\text{cm-}y} = Ma_{\text{cm-}y} \implies \eqbox{a_{\text{cm-}y} = \frac{2}{3}g}, \qquad \eqbox{T = \frac{1}{3}Mg}
\end{equation*}
\end{example}

\begin{example}{10.7}{Acceleration of a rolling sphere\bk{p.\ 319}}
\textbf{Problem:} A uniform solid sphere ($c = 2/5$) rolls without slipping down a ramp inclined at angle $\beta$. Find $a_{\text{cm-}x}$ and the static friction force $f$.
\tcblower
\textbf{Solution:} Using Eq.\ \eqref{eq:roll_incline} with $c = 2/5$:
\begin{equation*}
  \eqbox{a_{\text{cm-}x} = \frac{g\sin\beta}{1 + 2/5} = \frac{5}{7}g\sin\beta}, \qquad \eqbox{f = \frac{2/5}{1 + 2/5}Mg\sin\beta = \frac{2}{7}Mg\sin\beta}
\end{equation*}
\end{example}

\begin{example}{10.8}{Calculating power from torque\bk{p.\ 322}}
\textbf{Problem:} A motor applies constant torque $\tau_z = 10\text{ N}\cdot\text{m}$ to a grindstone ($I = 2.0\text{ kg}\cdot\text{m}^2$) starting from rest for $t = 8.0\text{ s}$. Find $W$, $K$, and $P_{\text{av}}$.
\tcblower
\textbf{Solution:}
\begin{align*}
  \alpha_z &= \frac{\tau_z}{I} = \frac{10}{2.0} = 5.0\text{ rad/s}^2, \qquad \Delta\theta = \frac{1}{2}\alpha_z t^2 = \frac{1}{2}(5.0)(8.0)^2 = 160\text{ rad}, \qquad \omega_z = \alpha_z t = 40\text{ rad/s}\\
  W &= \tau_z\Delta\theta = (10)(160) = \eqbox{1600\text{ J}} = K, \qquad P_{\text{av}} = \frac{W}{\Delta t} = \frac{1600\text{ J}}{8.0\text{ s}} = \eqbox{200\text{ W}}
\end{align*}
\end{example}

\begin{example}{10.9}{Angular momentum and torque\bk{p.\ 324}}
\textbf{Problem:} A turbine fan ($I = 2.5\text{ kg}\cdot\text{m}^2$) has $\omega_z = (40\text{ rad/s}^3)t^2$. Find (a) $L_z(t)$ and $L_z(3.0\text{ s})$, and (b) $\tau_z(t)$ and $\tau_z(3.0\text{ s})$.
\tcblower
\textbf{Solution:}
\begin{align*}
  \text{(a)}\quad L_z(t) &= I\omega_z = (2.5)(40)t^2 = \eqbox{(100\text{ kg}\cdot\text{m}^2/\text{s}^3)t^2}, \qquad L_z(3.0\text{ s}) = \eqbox{900\text{ kg}\cdot\text{m}^2/\text{s}}\\
  \text{(b)}\quad \tau_z(t) &= \frac{\diff L_z}{\diff t} = \eqbox{(200\text{ kg}\cdot\text{m}^2/\text{s}^3)t}, \qquad \tau_z(3.0\text{ s}) = \eqbox{600\text{ N}\cdot\text{m}}
\end{align*}
\end{example}

\begin{example}{10.10}{Anyone can be a ballerina\bk{p.\ 326}}
\textbf{Problem:} A professor on a frictionless turntable ($\omega_{1z} = 1\text{ rev}/2.0\text{ s} = 0.50\text{ rev/s} = 3.14\text{ rad/s}$) holds two $5.0\text{-kg}$ dumbbells at $r_1 = 1.0\text{ m}$ ($I_{\text{prof},1} = 3.0\text{ kg}\cdot\text{m}^2$) and pulls them in to $r_2 = 0.20\text{ m}$ ($I_{\text{prof},2} = 2.2\text{ kg}\cdot\text{m}^2$). Find $\omega_{2z}$, $K_1$, and $K_2$.
\tcblower
\textbf{Solution:}
\begin{align*}
  I_1 &= 3.0 + 2(5.0)(1.0)^2 = 13\text{ kg}\cdot\text{m}^2, \qquad I_2 = 2.2 + 2(5.0)(0.20)^2 = 2.6\text{ kg}\cdot\text{m}^2\\
  \omega_{2z} &= \frac{I_1}{I_2}\omega_{1z} = \frac{13}{2.6}(0.50\text{ rev/s}) = \eqbox{2.5\text{ rev/s} = 15.7\text{ rad/s}}\\
  K_1 &= \frac{1}{2}(13)(3.14)^2 = \eqbox{64\text{ J}}, \qquad K_2 = \frac{1}{2}(2.6)(15.7)^2 = \eqbox{320\text{ J}}
\end{align*}
\end{example}

\begin{example}{10.11}{A rotational ``collision''\bk{p.\ 326}}
\textbf{Problem:} Two disks with moments of inertia $I_A, I_B$ and initial angular speeds $\omega_A, \omega_B$ are pushed together along their common axis until they reach a common angular speed $\omega$. Find $\omega$.
\tcblower
\textbf{Solution:} Conservation of angular momentum gives:
\begin{equation*}
  I_A\omega_A + I_B\omega_B = (I_A + I_B)\omega \implies \eqbox{\omega = \frac{I_A\omega_A + I_B\omega_B}{I_A + I_B}}
\end{equation*}
\end{example}

\begin{example}{10.12}{Angular momentum in a crime bust\bk{p.\ 327}}
\textbf{Problem:} A bullet ($m = 0.010\text{ kg}$, $v = 400\text{ m/s}$) strikes the center ($l = 0.50\text{ m}$) of a door ($M = 15\text{ kg}$, width $d = 1.00\text{ m}$) perpendicular to the door and embeds in it. Find $\omega$, $K_1$, and $K_2$.
\tcblower
\textbf{Solution:}
\begin{align*}
  L &= mvl = (0.010)(400)(0.50) = 2.0\text{ kg}\cdot\text{m}^2/\text{s}, \qquad I = \frac{1}{3}Md^2 + ml^2 = 5.0 + 0.0025 = 5.0025\text{ kg}\cdot\text{m}^2\\
  \omega &= \frac{L}{I} = \frac{2.0}{5.0025} = \eqbox{0.40\text{ rad/s}}, \qquad K_1 = \frac{1}{2}mv^2 = \eqbox{800\text{ J}}, \qquad K_2 = \frac{1}{2}I\omega^2 = \eqbox{0.40\text{ J}}
\end{align*}
\end{example}

\begin{example}{10.13}{A precessing gyroscope\bk{p.\ 330}}
\textbf{Problem:} A solid cylindrical gyroscope wheel ($R = 3.0\text{ cm} = 0.030\text{ m}$, pivot distance $r = 2.0\text{ cm} = 0.020\text{ m}$) has spin vector $\bm{\omega}$ pointing left as seen from above and takes $4.0\text{ s}$ per revolution of precession ($\Omega = 2\pi/4.0 = 1.57\text{ rad/s}$). Find (a) the sense of precession and (b) spin angular speed $\omega$.
\tcblower
\textbf{Solution:}
\begin{itemize}
  \item \textbf{(a)} $\vect{r}$ points left and $\vect{w}$ points into the page, so $\bm{\tau} = \vect{r}\times\vect{w}$ points toward the top of the page. Thus $\vect{L}$ turns \eqbox{\text{clockwise}} as seen from above.
  \item \textbf{(b)} With $I = \frac{1}{2}mR^2$, Eq.\ \eqref{eq:precession} gives:
  \begin{equation*}
    \omega = \frac{mgr}{\left(\frac{1}{2}mR^2\right)\Omega} = \frac{2gr}{R^2\Omega} = \frac{2(9.8)(0.020)}{(0.030)^2(1.57)} = \eqbox{280\text{ rad/s} = 2600\text{ rev/min}}
  \end{equation*}
\end{itemize}
\end{example}

\begin{example}{Ch.\ 10}{Bridging Problem: Billiard Physics\bk{p.\ 332}}
\textbf{Problem:} A uniform solid sphere ($m, R$, $I_{\text{cm}} = \frac{2}{5}mR^2$) is struck horizontally with force $F$ for a short time $\Delta t$ at height $h$ above its center. (a) For what height $h$ does the ball roll without slipping immediately after the hit? (b) If struck dead center ($h=0$) giving initial speed $v_0 = F\Delta t/m$ and $\omega_0 = 0$, find the final center-of-mass speed $v_{\text{cm}}$ when rolling without slipping begins.
\tcblower
\textbf{Solution:}
\begin{itemize}
  \item \textbf{(a)} Linear impulse: $F\Delta t = mv_{\text{cm}}$. Angular impulse about cm: $(Fh)\Delta t = I_{\text{cm}}\omega = \left(\frac{2}{5}mR^2\right)\!\left(\frac{v_{\text{cm}}}{R}\right) = \frac{2}{5}mRv_{\text{cm}}$. Dividing the two equations gives:
  \begin{equation*}
    \eqbox{h = \frac{2}{5}R}
  \end{equation*}
  \item \textbf{(b)} When hit at $h=0$, $v_{\text{cm}}(0) = v_0$ and $\omega(0) = 0$. While skidding, kinetic friction $f_k = \mu_k mg$ decelerates translation ($a_{\text{cm}} = -\mu_k g$) and accelerates rotation ($\alpha = \frac{f_k R}{\frac{2}{5}mR^2} = \frac{5\mu_k g}{2R}$):
  \begin{align*}
    v_{\text{cm}}(t) &= v_0 - \mu_k gt, \qquad R\omega(t) = \frac{5}{2}\mu_k gt
  \end{align*}
  Rolling without slipping starts when $v_{\text{cm}}(t) = R\omega(t) \implies v_0 - \mu_k gt = \frac{5}{2}\mu_k gt \implies \mu_k gt = \frac{2}{7}v_0$. Thus:
  \begin{equation*}
    \eqbox{v_{\text{cm,final}} = v_0 - \frac{2}{7}v_0 = \frac{5}{7}v_0}
  \end{equation*}
\end{itemize}
\end{example}

\begin{example}{7.2}{Puck on a String Pulled Through a Hole\bk{Analytical Mechanics p.\ 183}}
\textbf{Problem:} A puck of mass $m$ revolves on a frictionless table at initial radius $d_0$ and angular velocity $\omega_0$. The string is pulled from below to reduce the radius to $d_f = d_0/2$. Find: (a) $\Delta l$, (b) $\omega_f$, (c) $\Delta T$, and (d) the work $W$ done in pulling the string.
\tcblower
\textbf{Solution:}
\begin{itemize}
  \item \textbf{(a)} Since string tension is central ($\vect{r}\parallel\vect{F}$), $\vect{N} = \vect{r}\times\vect{F} = \vect{0}$, so $\eqbox{\Delta l = 0}$.
  \item \textbf{(b)} $md_f^2\omega_f = md_0^2\omega_0 \implies \omega_f = \left(\frac{d_0}{d_0/2}\right)^2\omega_0 = \eqbox{4\omega_0}$.
  \item \textbf{(c)--(d)} With $T_0 = \frac{1}{2}md_0^2\omega_0^2$, $T_f = \frac{1}{2}m(d_0/2)^2(4\omega_0)^2 = 4T_0$, so $\eqbox{\Delta T = W = 3T_0 = \frac{3}{2}md_0^2\omega_0^2}$.
\end{itemize}
\end{example}

% ============================================================================
% UNIQUE SOLVED PROBLEMS — MODULE 3
% ============================================================================
\section{Unique Solved Problems (Module 3)}

\subsubsection{Part A: Analytical Mechanics Chapter 7 Exercises (with Complete Derivations)}

\begin{enumerate}[label=\textbf{Ex.\ 7.\arabic*},leftmargin=3.5em]
  \item \textbf{(Automobile Angular Momentum):} $m = 1000\text{ kg}$, $v = 100\text{ km/h} = 27.78\text{ m/s}$, perpendicular distance $b = 20\text{ m}$:
  \begin{equation*}
    l = mvb = (1000\text{ kg})(27.78\text{ m/s})(20\text{ m}) = \eqbox{5.56\times 10^5\text{ kg}\cdot\text{m}^2/\text{s}} \quad(\text{downward for one side, upward for the other})
  \end{equation*}

  \item \textbf{(Second Derivative of $\vect{l}$):} Particle ($m$) at origin at $t=0$ with $\vect{v} = v_0\ihat$ under constant force $\vect{F} = f_0\jhat$:
  \begin{equation*}
    \frac{\diff\vect{l}}{\diff t} = \vect{r}\times\vect{F} \implies \frac{\diff^2\vect{l}}{\diff t^2} = \frac{\diff\vect{r}}{\diff t}\times\vect{F} + \vect{r}\times\underbrace{\frac{\diff\vect{F}}{\diff t}}_{\vect{0}} = \vect{v}\times\vect{F}\Big|_{t=0} = (v_0\ihat)\times(f_0\jhat) = \eqbox{v_0 f_0\khat}
  \end{equation*}

  \item \textbf{(Girl Jumping onto Merry-Go-Round):} Girl ($m = 20\text{ kg}$, $v = 3\text{ m/s}$) jumps tangentially onto rim of stationary disk ($M = 50\text{ kg}$, $R = 1\text{ m}$, $I_{\text{disk}} = \frac{1}{2}MR^2 = 25\text{ kg}\cdot\text{m}^2$):
  \begin{equation*}
    L_i = mvR = (20)(3)(1) = 60\text{ kg}\cdot\text{m}^2/\text{s} = (I_{\text{disk}} + mR^2)\omega_f = (25 + 20)\omega_f \implies \eqbox{\omega_f = \frac{60}{45} = \frac{4}{3}\text{ rad/s}}
  \end{equation*}

  \item \textbf{(Atwood's Machine in CGS Units):} $m_1 = 500\text{ g}$, $m_2 = 700\text{ g}$, pulley $r = 2\text{ cm}$, $I = 100\text{ g}\cdot\text{cm}^2$, $g = 980\text{ cm/s}^2$.
  \begin{itemize}
    \item Equations of motion: $m_2 g - T_2 = m_2 a$, $T_1 - m_1 g = m_1 a$, and $\frac{\diff L}{\diff t} = (T_2 - T_1)r = I\frac{a}{r}$.
    \item Adding the three equations gives $(m_2 - m_1)g = \left(m_1 + m_2 + \frac{I}{r^2}\right)a \implies a = \frac{200(980)}{500 + 700 + 100/4} = \frac{196,000}{1225} = 160\text{ cm/s}^2$.
    \item Rate of change of pulley angular momentum: $\frac{\diff L}{\diff t} = I\alpha = I\frac{a}{r} = 100\!\left(\frac{160}{2}\right) = \eqbox{8000\text{ dyne}\cdot\text{cm}}$.
  \end{itemize}

  \item \textbf{(Spinning Top Torque):} Top of mass $M$, CM at distance $R$ from tip, tilted at angle $\theta$ from vertical:
  \begin{equation*}
    \vect{N} = \vect{R}\times M\vect{g} \implies \eqbox{N = RMg\sin\theta} \quad(\text{horizontal, perpendicular to }\bm{\omega}\text{ and }\vect{g})
  \end{equation*}

  \item \textbf{(Dumbbell Mass Flying Off):} Two point masses $M$ on massless rod of length $d$ rotate about CM at $\omega_0$. Each mass has speed $v_0 = \frac{d}{2}\omega_0$. When one mass comes loose, no force acts on the remaining mass, so it moves in a straight line at constant speed $v_0$ and its angular velocity about the original pivot is $\eqbox{\omega_0}$ (or $0$ about its own CM).

  \item \textbf{(Proof that $\sum_i m_i\vect{r}_i' = \vect{0}$):} Since $\vect{r}_i = \vect{r}_c + \vect{r}_i'$ and $M\vect{r}_c = \sum_i m_i\vect{r}_i$, summing $m_i\vect{r}_i'$ gives $\sum_i m_i(\vect{r}_i - \vect{r}_c) = \sum_i m_i\vect{r}_i - M\vect{r}_c = \eqbox{\vect{0}}$.

  \item \textbf{(Exercise 7.6[9] on p.\ 188 --- Gum Falling on Spinning Disk):} Disk ($R = 10\text{ cm}$, $M = 15\text{ g}$, $I_{\text{disk}} = \frac{1}{2}MR^2 = 750\text{ g}\cdot\text{cm}^2$) spins at $\omega_i = 50\text{ rad/s}$. Gum ($m = 3\text{ g}$) lands on rim ($I_{\text{gum}} = mR^2 = 300\text{ g}\cdot\text{cm}^2$):
  \begin{equation*}
    \omega_f = \frac{I_{\text{disk}}}{I_{\text{disk}} + I_{\text{gum}}}\omega_i = \frac{750}{750 + 300}(50\text{ rad/s}) = \frac{5}{7}(50) = \eqbox{35.7\text{ rad/s}}
  \end{equation*}

  \setcounter{enumi}{9}
  \item \textbf{(Rod about Axis at Distance $a$ Beyond End):} Here $d = \frac{l}{2} + a$ from the center of mass:
  \begin{equation*}
    I = I_c + Md^2 = \frac{1}{12}Ml^2 + M\!\left(\frac{l}{2} + a\right)^2 = \eqbox{Ma^2 + Mal + \frac{1}{3}Ml^2}
  \end{equation*}

  \item \textbf{(Three Exercises Numbered 7.11 on p.\ 189):}
  \begin{itemize}
    \item \textbf{7.11(a) Disk Angular Acceleration:} $F = 2\text{ N}$, $M = 15\text{ kg}$, $R = 0.20\text{ m} \implies \alpha = \frac{FR}{\frac{1}{2}MR^2} = \frac{2(2)}{15(0.20)} = \eqbox{\frac{4}{3}\text{ rad/s}^2}$.
    \item \textbf{7.11(b) Square Plate In-Plane Axis:} Given $I_z = \frac{1}{6}Ma^2$, by symmetry $I_x = I_y$ and $I_x + I_y = I_z \implies \eqbox{I_x = \frac{1}{12}Ma^2}$.
    \item \textbf{7.11(c) Ring/Hoop from Definition:} All mass elements $\diff m$ are at $r_\perp = R$, so $I = \int R^2\dd m = R^2\int\dd m = \eqbox{MR^2}$.
  \end{itemize}

  \setcounter{enumi}{13}
  \item \textbf{(Meter Stick Pivoted at $20\text{ cm}$ Mark):} $M = 0.25\text{ kg}$, $l = 1.00\text{ m}$, pivot at $0.20\text{ m}$ ($d = 0.30\text{ m}$ from CM at $0.50\text{ m}$), force $F = 3\text{ N}$ at $0.80\text{ m}$ (lever arm $r_F = 0.80 - 0.20 = 0.60\text{ m}$):
  \begin{equation*}
    I = \frac{1}{12}Ml^2 + Md^2 = 0.25\!\left(\frac{1}{12} + 0.09\right) = 0.04333\text{ kg}\cdot\text{m}^2, \qquad \alpha = \frac{F r_F}{I} = \frac{3(0.60)}{0.04333} = \eqbox{41.5\text{ rad/s}^2}
  \end{equation*}

  \item \textbf{(Ring Gyroscope Precession Rate):} Ring of mass $M$ and radius $a$ has $I = Ma^2$ and CM at distance $d$ from pivot:
  \begin{equation*}
    \Omega_p = \frac{Mgd}{I\omega} = \frac{Mgd}{Ma^2\omega} = \eqbox{\frac{gd}{a^2\omega}}
  \end{equation*}
\end{enumerate}

\subsubsection{Part B: Unique \& Challenge Problems from Young \& Freedman (PDF 3)}

\begin{enumerate}[label=\textbf{Prob.\ 10.\arabic*},leftmargin=4.5em]
  \setcounter{enumi}{65}
  \item \textbf{[Prob.\ 10.66, p.\ 338--339] Balancing a Rod with a Clay Blob (Pool Cue Physics):}
  Uniform rod ($L, M$) with clay blob ($M$) at one end.
  \begin{itemize}
    \item \textbf{(a)} Center of mass is at distance $\frac{M(L/2) + M(L)}{2M} = \eqbox{\frac{3}{4}L}$ from the clay-free end (or $\frac{1}{4}L$ from the clay end).
    \item \textbf{(b) Clay at top, pivoted at bottom:} $I_{\text{bottom}} = \frac{1}{3}ML^2 + ML^2 = \frac{4}{3}ML^2$, $\tau = (2M)g\!\left(\frac{3}{4}L\right)\sin\theta \approx \frac{3}{2}MgL\theta \implies \alpha = \frac{\tau}{I} = \eqbox{\frac{9g}{8L}\theta}$.
    \item \textbf{(c) Clay at bottom, pivoted at bottom:} $I_{\text{bottom}} = \frac{1}{3}ML^2 + 0 = \frac{1}{3}ML^2$, $\tau = Mg\!\left(\frac{L}{2}\right)\sin\theta \approx \frac{1}{2}MgL\theta \implies \alpha = \eqbox{\frac{3g}{2L}\theta}$ (larger angular acceleration, so it tips over faster and is harder to balance!).
  \end{itemize}

  \setcounter{enumi}{74}
  \item \textbf{[Prob.\ 10.75, p.\ 339] Two-Disk Yo-Yo with Axle of Radius $b$:}
  Total mass $M_{\text{tot}} = 2m$, moment of inertia $I_{\text{cm}} = 2\!\left(\frac{1}{2}mR^2\right) = mR^2$. Unwinding without slipping on axle of radius $b$ gives $a_y = b\alpha_z$.
  \begin{align*}
    2mg - T &= 2ma_y, \qquad Tb = I_{\text{cm}}\alpha_z = mR^2\!\left(\frac{a_y}{b}\right) \implies T = m\frac{R^2}{b^2}a_y\\
    \eqbox{a_y = \frac{g}{1 + R^2/(2b^2)}, \qquad \alpha_z = \frac{g}{b + R^2/(2b)}, \qquad T = \frac{2mg}{1 + 2b^2/R^2}}
  \end{align*}

  \item \textbf{[Prob.\ 10.76, p.\ 339--340] Hollow Spherical Shell Rolling in a Loop-the-Loop:}
  For a thin-walled hollow sphere ($c = 2/3$), $K = \frac{1}{2}(1+2/3)mv^2 = \frac{5}{6}mv^2$.
  \begin{itemize}
    \item \textbf{(a)} At the top of the loop ($y = 2R$), minimum speed for contact is $v_{\text{top}}^2 = gR$. Conservation of energy from $h_0$ to $2R$:
    \begin{equation*}
      mgh_0 = mg(2R) + \frac{5}{6}mv_{\text{top}}^2 = 2mgR + \frac{5}{6}mgR \implies \eqbox{h_0 = \frac{17}{6}R}
    \end{equation*}
    \item \textbf{(b)} At point $B$ ($y = R$): $mg(h_0 - R) = \frac{11}{6}mgR = \frac{5}{6}mv_B^2 \implies v_B^2 = \frac{11}{5}gR \implies N_B = \frac{mv_B^2}{R} = \eqbox{\frac{11}{5}mg = 2.20mg}$.
  \end{itemize}

  \setcounter{enumi}{98}
  \item \textbf{[Prob.\ 10.99, p.\ 342] Center of Percussion of a Baseball Bat:}
  Bat has mass $M = 0.800\text{ kg}$, length $L = 0.900\text{ m}$, CM at $d = 0.600\text{ m}$ from handle end, and $I_{\text{cm}} = 0.0530\text{ kg}\cdot\text{m}^2$. Impact impulse $J$ at distance $x$ from handle gives $v_{\text{cm}} = \frac{J}{M}$ and $\omega = \frac{J(x - d)}{I_{\text{cm}}}$. For the handle end to remain at rest ($v_{\text{handle}} = v_{\text{cm}} - d\omega = 0$):
  \begin{equation*}
    \frac{J}{M} = d\frac{J(x - d)}{I_{\text{cm}}} \implies \eqbox{x = d + \frac{I_{\text{cm}}}{Md}} = 0.600\text{ m} + \frac{0.0530\text{ kg}\cdot\text{m}^2}{(0.800\text{ kg})(0.600\text{ m})} = \eqbox{0.710\text{ m}}
  \end{equation*}

  \item \textbf{[Challenge Problem 10.100, p.\ 342] Ball Rolling Between Two Parallel Rails:}
  \begin{itemize}
    \item \textbf{(a)} The contact points on the two rails (separated by $d$) are at perpendicular distance $r_{\text{eff}} = \sqrt{R^2 - (d/2)^2} = \sqrt{R^2 - d^2/4}$ from the ball's rotation axis, so instantaneous rolling without slipping gives $\eqbox{v_{\text{cm}} = \omega\sqrt{R^2 - d^2/4}}$.
    \item \textbf{(b)} Energy conservation $Mgh = \frac{1}{2}Mv_{\text{cm}}^2 + \frac{1}{2}\!\left(\frac{2}{5}MR^2\right)\omega^2$ with $\omega^2 = \frac{v_{\text{cm}}^2}{R^2 - d^2/4}$ yields:
    \begin{equation*}
      gh = \frac{1}{2}v_{\text{cm}}^2\!\left[1 + \frac{2}{5(1 - d^2/4R^2)}\right] \implies \eqbox{v_{\text{cm}} = \sqrt{\frac{10gh}{5 + 2/(1 - d^2/4R^2)}}}
    \end{equation*}
  \end{itemize}

  \item \textbf{[Challenge Problem 10.101, p.\ 342] Skidding Cylinder Transition to Rolling Without Slipping:}
  Solid cylinder ($M, R, I_{\text{cm}} = \frac{1}{2}MR^2$) is set down with $v_{0x} = 0$ and initial spin $\omega_0$.
  \begin{itemize}
    \item \textbf{(a)} The bottom point slips backward, so kinetic friction $f_k = \mu_k Mg$ acts forward ($+x$):
    \begin{equation*}
      \eqbox{a_x = \frac{f_k}{M} = \mu_k g, \qquad \alpha_z = -\frac{f_k R}{\frac{1}{2}MR^2} = -\frac{2\mu_k g}{R}}
    \end{equation*}
    \item \textbf{(b)} $v_x(t) = \mu_k gt$ and $\omega_z(t) = \omega_0 - \frac{2\mu_k g}{R}t$. Setting $v_x(t^*) = R\omega_z(t^*)$ gives $\mu_k gt^* = R\omega_0 - 2\mu_k gt^* \implies t^* = \frac{R\omega_0}{3\mu_k g}$. The distance traveled before slipping stops is:
    \begin{equation*}
      d = \frac{1}{2}a_x(t^*)^2 = \frac{1}{2}(\mu_k g)\!\left(\frac{R\omega_0}{3\mu_k g}\right)^2 = \eqbox{\frac{R^2\omega_0^2}{18\mu_k g}}
    \end{equation*}
    \item \textbf{(c)} Final angular speed is $\omega_f = \frac{1}{3}\omega_0$ and $v_f = \frac{1}{3}R\omega_0$. Work done by friction on the cylinder is:
    \begin{equation*}
      W_{\text{fric}} = K_f - K_i = \left[\frac{1}{2}M\!\left(\frac{R\omega_0}{3}\right)^2 + \frac{1}{2}\!\left(\frac{1}{2}MR^2\right)\!\left(\frac{\omega_0}{3}\right)^2\right] - \frac{1}{2}\!\left(\frac{1}{2}MR^2\right)\omega_0^2 = \eqbox{-\frac{1}{6}MR^2\omega_0^2}
    \end{equation*}
  \end{itemize}

  \setcounter{enumi}{102}
  \item \textbf{[Challenge Problem 10.103, p.\ 342] Work Integral for Block Pulled Inward on a String:}
  \begin{itemize}
    \item \textbf{(a)} Angular momentum $L = mv_1 r_1 = mv(r)r \implies v(r) = \frac{v_1 r_1}{r}$. String tension provides centripetal force:
    \begin{equation*}
      \eqbox{T(r) = \frac{mv(r)^2}{r} = \frac{mv_1^2 r_1^2}{r^3}}
    \end{equation*}
    \item \textbf{(b)--(c)} Since $\vect{T}$ points radially inward ($-\uv{r}$), $\vect{T}(r)\cdot\diff\vect{r} = -T(r)\dd r$:
    \begin{equation*}
      W = \int_{r_1}^{r_2} -\frac{mv_1^2 r_1^2}{r^3}\dd r = \left[\frac{1}{2}\frac{mv_1^2 r_1^2}{r^2}\right]_{r_1}^{r_2} = \eqbox{\frac{1}{2}mv_1^2\!\left(\frac{r_1^2}{r_2^2} - 1\right) = \Delta K}
    \end{equation*}
  \end{itemize}
\end{enumerate}

% ============================================================================
% FORMULA SHEET — MODULE 3
% ============================================================================
\formulasheet{Module 3 Formula Sheet: Rotational Dynamics and Angular Momentum}

\begin{fs}{1. Torque, Rotational Dynamics, Rolling, Work, and Power}
\begin{align*}
  \bm{\tau} &= \vect{r}\times\vect{F}, \qquad \tau = Fl = rF\sin\phi = F_{\text{tan}}r, \qquad \sum\tau_z = I\alpha_z\\
  K_{\text{tot}} &= \frac{1}{2}Mv_{\text{cm}}^2 + \frac{1}{2}I_{\text{cm}}\omega^2, \qquad v_{\text{cm}} = R\omega, \qquad a_{\text{cm-}x} = \frac{g\sin\beta}{1+c}, \qquad f_s = \frac{c}{1+c}Mg\sin\beta\\
  W &= \int_{\theta_1}^{\theta_2}\tau_z\dd\theta = \frac{1}{2}I\omega_2^2 - \frac{1}{2}I\omega_1^2, \qquad P = \tau_z\omega_z
\end{align*}
\end{fs}

\begin{fs}{2. Angular Momentum, Conservation, and Gyroscope Precession}
\begin{align*}
  \vect{l} &= \vect{r}\times\vect{p}, \qquad \frac{\diff\vect{L}}{\diff t} = \sum\bm{\tau}_{\text{ext}}, \qquad \vect{L} = M(\vect{r}_c\times\dot{\vect{r}}_c) + \sum m_i(\vect{r}_i'\times\dot{\vect{r}}_i') = \vect{L}_c + \vect{L}'\\
  \vect{l}_i &= m_i r_i^2\omega\sin^2\theta_i\khat - m_i r_i^2\omega\sin\theta_i\cos\theta_i\uv{e}_i \implies \vect{L} = I\bm{\omega}\ (\text{symmetric}), \qquad I_1\omega_{1z} = I_2\omega_{2z}, \qquad \Omega_p = \frac{Mgr}{I\omega}
\end{align*}
\end{fs}

% ============================================================================
% MODULE 4: GRAVITATION AND THE GRAVITATIONAL FIELD (PDF 2)
% ============================================================================
\part{Gravitation and The Gravitational Field}

\chapter{Newton's Law of Gravitation, Satellites, and Kepler's Laws}
\begin{syll}
\textbf{Syllabus Coverage (PDF 2):} Young \& Freedman Chapter 13 (Sections 13.1--13.5, 13.7--13.8) and Analytical Mechanics Chapter 9 (Section 9.1). Universal gravitation in scalar and vector forms, weight and planetary density, gravitational potential energy, escape speed, satellite orbits, Kepler's laws, apparent weight on a rotating planet, and black holes.
\end{syll}

\section{Newton's Law of Universal Gravitation and Weight\bk{13.1, 13.2, 9.1}}

\begin{goal}
Express Newton's law of universal gravitation in scalar and vector forms, and relate surface gravity $g$ to planetary mass $m_E$, radius $R_E$, and mean density $\rho$.
\end{goal}

\begin{pic}
\begin{tikzpicture}[scale=0.95]
  \draw[ax] (0,0) -- (4.2,0);
  \draw[ax] (0,0) -- (0,2.5);
  \coordinate (M1) at (1.0,2.0);
  \coordinate (M2) at (3.4,0.9);
  \draw[vec,slate] (0,0) -- (M1) node[midway,above left,font=\footnotesize]{$\vect{r}_1$};
  \draw[vec,slate] (0,0) -- (M2) node[midway,below right,font=\footnotesize]{$\vect{r}_2$};
  \draw[vec,teal] (M1) -- (M2) node[midway,above,font=\footnotesize]{$\vect{r}_2 - \vect{r}_1$};
  \draw[vec,rose,ultra thick] (M2) -- ($(M2)!0.42!(M1)$) node[below left,font=\footnotesize]{$\vect{F}_{21}$};
  \fill[teal] (M1) circle (4pt) node[above left,font=\small]{$m_1$};
  \fill[amber] (M2) circle (4pt) node[right,font=\small]{$m_2$};
\end{tikzpicture}
\figcap{Vector form of Newton's law of gravitation: $\vect{F}_{21}$ is the attractive force on $m_2$ due to $m_1$.}
\end{pic}

\step{Step 1: Scalar and vector forms of the gravitational force (Eqs.\ 13.1, 9.1, 9.2).}
Let $\uv{r} = \frac{\vect{r}_2 - \vect{r}_1}{|\vect{r}_2 - \vect{r}_1|}$ be the unit vector pointing from $m_1$ to $m_2$:
\begin{equation}\label{eq:F_grav_vec}
  \eqbox{F_g = \frac{Gm_1 m_2}{r^2}}, \qquad \eqbox{\vect{F}_{21} = -G\frac{m_1 m_2}{|\vect{r}_2 - \vect{r}_1|^2}\uv{r} = -Gm_1 m_2\frac{\vect{r}_2 - \vect{r}_1}{|\vect{r}_2 - \vect{r}_1|^3}}
\end{equation}
where $G = 6.67\times 10^{-11}\text{ N}\cdot\text{m}^2/\text{kg}^2$.

\step{Step 2: Weight, surface acceleration due to gravity $g$, and average density $\rho$.}
At the surface of a spherically symmetric planet of mass $m_E$ and radius $R_E$, weight is $w = mg = \frac{Gm_E m}{R_E^2}$:
\begin{equation}\label{eq:g_and_rho}
  \eqbox{g = \frac{Gm_E}{R_E^2}}, \qquad \eqbox{w(r) = \frac{Gm_E m}{r^2}\quad (r\ge R_E)}, \qquad \eqbox{\rho = \frac{m_E}{\frac{4}{3}\pi R_E^3} = \frac{3g}{4\pi G R_E}}
\end{equation}

\section{Gravitational Potential Energy and Escape Speed\bk{13.3}}

\begin{goal}
Derive $U(r) = -\frac{Gm_E m}{r}$ by integrating the gravitational work from $r_1$ to $r_2$, and derive the escape speed $v_{\text{esc}}$.
\end{goal}

\step{Step 1: Integrate radial gravitational work from $r_1$ to $r_2$.}
The radial component of gravitational force is inward, $F_r = -\frac{Gm_E m}{r^2}$:
\begin{align}
  W_{\text{grav}} &= \int_{r_1}^{r_2} F_r\dd r = -Gm_E m\int_{r_1}^{r_2}\frac{\diff r}{r^2} = \left[\frac{Gm_E m}{r}\right]_{r_1}^{r_2} = \frac{Gm_E m}{r_2} - \frac{Gm_E m}{r_1} = U_1 - U_2\label{eq:W_grav_int}
\end{align}
Choosing $U(\infty) = 0$ gives the gravitational potential energy:
\begin{equation}\label{eq:U_grav_gen}
  \eqbox{U(r) = -\frac{Gm_E m}{r}}
\end{equation}
Near Earth's surface ($r_1, r_2 \approx R_E$), $W_{\text{grav}} = Gm_E m\frac{r_1 - r_2}{r_1 r_2} \approx \frac{Gm_E m}{R_E^2}(r_1 - r_2) = mg(y_1 - y_2)$, recovering $\Delta U = mg\Delta y$.

\step{Step 2: Derive the escape speed $v_{\text{esc}}$ from a spherical body of mass $M$ and radius $R$.}
Set total mechanical energy at $r = R$ equal to zero at $r = \infty$ ($K_2 = 0, U_2 = 0$):
\begin{equation}\label{eq:v_esc}
  \frac{1}{2}mv_{\text{esc}}^2 + \left(-\frac{GMm}{R}\right) = 0 \implies \eqbox{v_{\text{esc}} = \sqrt{\frac{2GM}{R}} = \sqrt{\frac{8\pi G\rho}{3}}R}
\end{equation}

\section{The Motion of Satellites\bk{13.4}}

\begin{goal}
Derive the circular orbital speed $v$, period $T$, and total mechanical energy $E$ of a satellite at radius $r$.
\end{goal}

\step{Step 1: Equate gravitational attraction to centripetal force.}
\begin{equation}\label{eq:v_orb}
  \frac{Gm_E m}{r^2} = \frac{mv^2}{r} \implies \eqbox{v = \sqrt{\frac{Gm_E}{r}}}
\end{equation}

\step{Step 2: Derive orbital period $T$ and mechanical energy $E$.}
Substitute $v = \frac{2\pi r}{T}$ into Eq.\ \eqref{eq:v_orb}:
\begin{align}
  T &= \frac{2\pi r}{v} = 2\pi r\sqrt{\frac{r}{Gm_E}} = \eqbox{\frac{2\pi r^{3/2}}{\sqrt{Gm_E}}}\label{eq:T_orb}\\
  E &= K + U = \frac{1}{2}mv^2 - \frac{Gm_E m}{r} = \frac{1}{2}m\!\left(\frac{Gm_E}{r}\right) - \frac{Gm_E m}{r} = \eqbox{-\frac{Gm_E m}{2r} = \frac{U}{2} = -K}\label{eq:E_orb}
\end{align}

\section{Kepler's Laws and Planetary Motion\bk{13.5}}

\begin{goal}
State Kepler's three laws and prove that Kepler's second law (constant sector velocity $\frac{\diff A}{\diff t}$) is a direct consequence of conservation of angular momentum.
\end{goal}

\begin{pic}
\begin{tikzpicture}[scale=0.95]
  \draw[thick,teal] (0,0) ellipse (3.0cm and 1.7cm);
  \draw[ax] (-3.5,0) -- (3.5,0) node[right,font=\footnotesize]{Major axis ($2a$)};
  \draw[ax] (0,-2.0) -- (0,2.0);
  \coordinate (S) at (-1.8,0);
  \coordinate (S2) at (1.8,0);
  \fill[amber] (S) circle (5pt) node[below=4pt,font=\footnotesize\bfseries]{Sun $S$};
  \fill[slate] (S2) circle (2pt) node[below=4pt,font=\footnotesize]{$S'$};
  \node[left,font=\footnotesize,teal] at (-3.0,0.3) {Perihelion};
  \node[right,font=\footnotesize,teal] at (3.0,0.3) {Aphelion};
  \draw[dim] (-3.0,-0.6) -- (-1.8,-0.6) node[midway,below,font=\scriptsize]{$r_p = a(1-e)$};
  \draw[dim] (-1.8,-0.6) -- (0,-0.6) node[midway,below,font=\scriptsize]{$ea$};
  \draw[dim] (-1.8,-1.25) -- (3.0,-1.25) node[midway,below,font=\scriptsize]{$r_a = a(1+e)$};
  % Sector area dA
  \fill[ambersoft,opacity=0.8] (S) -- (40:3.0 and 1.7) -- (58:3.0 and 1.7) -- cycle;
  \draw[thick,amber] (S) -- (40:3.0 and 1.7) node[midway,below right,font=\scriptsize]{$r$};
  \draw[thick,amber] (S) -- (58:3.0 and 1.7);
  \node[font=\scriptsize,amber!80!black] at (0.3,0.95) {$\diff A$};
\end{tikzpicture}
\figcap{Elliptical planetary orbit of semi-major axis $a$ and eccentricity $e$: $r_p = a(1-e)$, $r_a = a(1+e)$.}
\end{pic}

\step{Step 1: Kepler's First Law (Elliptical Orbits).}
Each planet moves in an ellipse with the Sun $S$ at one focus. The perihelion and aphelion distances are:
\begin{equation}\label{eq:kepler1}
  \eqbox{r_p = a(1 - e), \qquad r_a = a(1 + e), \qquad r_p + r_a = 2a, \qquad e = \frac{r_a - r_p}{r_a + r_p}}
\end{equation}

\step{Step 2: Kepler's Second Law (Equal Areas in Equal Times).}
In time $\diff t$, the radius vector $SP$ turns through angle $\diff\theta$, sweeping out triangular area $\diff A = \frac{1}{2}r(r\dd\theta) = \frac{1}{2}r^2\dd\theta$. Since $v_\perp = v\sin\phi = r\frac{\diff\theta}{\diff t}$, the sector velocity is:
\begin{equation}\label{eq:kepler2}
  \eqbox{\frac{\diff A}{\diff t} = \frac{1}{2}r^2\frac{\diff\theta}{\diff t} = \frac{1}{2}rv\sin\phi = \frac{1}{2m}|\vect{r}\times m\vect{v}| = \frac{L}{2m} = \text{constant}}
\end{equation}
Because gravitational force is central ($\vect{F}\parallel\vect{r}$), $\frac{\diff\vect{L}}{\diff t} = \vect{r}\times\vect{F} = \vect{0}$, so $\vect{L}$ and $\frac{\diff A}{\diff t}$ are strictly constant. At perihelion and aphelion ($\phi = 90^\circ$), $L = mr_p v_p = mr_a v_a \implies r_p v_p = r_a v_a$.

\step{Step 3: Kepler's Third Law (Harmonic Law).}
\begin{equation}\label{eq:kepler3}
  \eqbox{T = \frac{2\pi a^{3/2}}{\sqrt{Gm_S}} \iff T^2 = \left(\frac{4\pi^2}{Gm_S}\right)a^3}
\end{equation}

\section{Apparent Weight and Black Holes\bk{13.7, 13.8}}

\begin{goal}
Derive the apparent weight on a rotating planet and the Schwarzschild radius $R_S$ of a black hole.
\end{goal}

\step{Step 1: Apparent weight on a rotating Earth.}
In general, $\vect{w} = \vect{w}_0 - m\vect{a}_{\text{rad}}$. At the equator, centripetal acceleration is $a_{\text{rad}} = \frac{v^2}{R_E} = \omega^2 R_E$ inward:
\begin{equation}\label{eq:app_weight}
  \eqbox{w = w_0 - \frac{mv^2}{R_E}, \qquad g = g_0 - \frac{v^2}{R_E} = g_0 - \omega^2 R_E\quad (\text{at equator})}
\end{equation}

\step{Step 2: Schwarzschild radius $R_S$ of a nonrotating black hole.}
Setting the escape speed $v_{\text{esc}} = \sqrt{\frac{2GM}{R_S}}$ equal to the speed of light $c$ yields Karl Schwarzschild's exact general-relativistic event-horizon radius:
\begin{equation}\label{eq:R_schwarz}
  c = \sqrt{\frac{2GM}{R_S}} \implies \eqbox{R_S = \frac{2GM}{c^2}}
\end{equation}

% ============================================================================
% CHAPTER 8: THE GRAVITATIONAL FIELD, GAUSS'S LAW, AND GENERAL RELATIVITY
% ============================================================================
\chapter{The Gravitational Field, Gauss's Law, and General Relativity}
\begin{syll}
\textbf{Syllabus Coverage (PDF 2):} Analytical Mechanics Chapter 9 (Sections 9.2--9.9) and Young \& Freedman Section 13.6. Gravitational field $\vect{g}(\vect{r})$ and potential $\Phi(\vect{r})$ of point and extended bodies, Spherical Shell Theorem, Gauss's law for gravity (integral and differential forms), Poisson's and Laplace's equations, and Einstein's field equations.
\end{syll}

\section{The Gravitational Field and Potential\bk{9.2--9.4}}

\begin{goal}
Define the gravitational field vector $\vect{g}(\vect{r})$ and scalar potential $\Phi(\vect{r})$ for discrete and continuous mass distributions, and show that $\vect{g} = -\nabla\Phi$.
\end{goal}

\step{Step 1: Gravitational field $\vect{g}(\vect{r})$ (force per unit test mass).}
For a point mass $M$ at source point $\vect{r}'$ and a continuous body with density $\rho(\vect{r}')$ and volume element $\diff\tau'$:
\begin{equation}\label{eq:g_field_def}
  \eqbox{\vect{g}(\vect{r}) = \lim_{m\to 0}\frac{\vect{F}}{m} = -GM\frac{\vect{r}-\vect{r}'}{|\vect{r}-\vect{r}'|^3}}, \qquad \eqbox{\vect{g}(\vect{r}) = -G\int_{\text{body}}\rho(\vect{r}')\frac{\vect{r}-\vect{r}'}{|\vect{r}-\vect{r}'|^3}\dd\tau'}
\end{equation}

\step{Step 2: Gravitational potential $\Phi(\vect{r})$ (potential energy per unit test mass).}
Since $\nabla\times\vect{g} = \vect{0}$, $\vect{g}$ is conservative and is the negative gradient of a scalar potential $\Phi(\vect{r}) = \lim_{m\to 0}\frac{V(\vect{r})}{m}$:
\begin{equation}\label{eq:Phi_def}
  \eqbox{\Phi(\vect{r}) = -\frac{GM}{r}\quad (\text{point mass at origin})}, \qquad \eqbox{\Phi(\vect{r}) = -G\int_{\text{body}}\frac{\rho(\vect{r}')\dd\tau'}{|\vect{r}-\vect{r}'|}}, \qquad \eqbox{\vect{g} = -\nabla\Phi}
\end{equation}

\section{Spherical Mass Distributions (The Shell Theorem)\bk{13.6, 9.4}}

\begin{goal}
Prove by integration that a uniform spherical shell of mass $M$ and radius $R$ (or $a$) produces potential $\Phi(r) = -\frac{GM}{r}$ and field $g_r = -\frac{GM}{r^2}$ outside ($r > R$), and constant potential $\Phi(r) = -\frac{GM}{R}$ and zero field ($g_r = 0$) inside ($r < R$).
\end{goal}

\begin{pic}
\begin{tikzpicture}[scale=0.95]
  \draw[thick,teal,fill=tealsoft!40] (0,0) circle (1.6cm);
  \draw[ax] (-2.0,0) -- (4.5,0);
  \fill[ink] (0,0) circle (2pt) node[below left,font=\footnotesize]{$O$};
  \coordinate (P) at (3.8,0);
  \fill[rose] (P) circle (2.5pt) node[below,font=\footnotesize]{$P\ (\text{mass }m)$};
  % Ring on sphere
  \draw[thick,amber,fill=ambersoft!80] ({1.6*cos(42)},-1.07) ellipse (0.15cm and 1.07cm);
  \coordinate (RingTop) at (42:1.6);
  \draw[thick,teal] (0,0) -- (RingTop) node[midway,above left,font=\footnotesize]{$R$};
  \draw[thick,rose] (RingTop) -- (P) node[midway,above,font=\footnotesize]{$s = |\vect{r}-\vect{r}'|$};
  \draw[dashed,slate] (RingTop) -- ({1.6*cos(42)},0) node[midway,left,font=\scriptsize]{$R\sin\phi$};
  \draw[slate] (0.5,0) arc (0:42:0.5) node[midway,right,font=\scriptsize]{$\phi$};
  \draw[dim] (0,-1.85) -- (3.8,-1.85) node[midway,below,font=\footnotesize]{$r$};
\end{tikzpicture}
\figcap{Integrating over thin rings of radius $R\sin\phi$ and width $R\dd\phi$ on a spherical shell of radius $R$ and mass $M$.}
\end{pic}

\step{Step 1: Express the ring mass $\diff M$ in terms of polar angle $\phi\in [0,\pi]$.}
The ring has radius $R\sin\phi$, circumference $2\pi R\sin\phi$, width $R\dd\phi$, and area $\diff A = 2\pi R^2\sin\phi\dd\phi$. Since the total shell area is $4\pi R^2$:
\begin{equation}\label{eq:shell_dM}
  \frac{\diff M}{M} = \frac{2\pi R^2\sin\phi\dd\phi}{4\pi R^2} = \frac{1}{2}\sin\phi\dd\phi \implies \diff U = -\frac{Gm\dd M}{s} = -\frac{GMm\sin\phi\dd\phi}{2s}
\end{equation}

\step{Step 2: Change integration variable from $\phi$ to $s$ using the Law of Cosines.}
In triangle $O\text{--Ring--}P$, the law of cosines gives:
\begin{equation}\label{eq:shell_cos}
  s^2 = r^2 + R^2 - 2rR\cos\phi \xrightarrow{\text{differentiate}} 2s\dd s = 2rR\sin\phi\dd\phi \implies \frac{\sin\phi\dd\phi}{s} = \frac{\dd s}{rR}
\end{equation}
Substitute Eq.\ \eqref{eq:shell_cos} into Eq.\ \eqref{eq:shell_dM}:
\begin{equation}\label{eq:shell_dU}
  \diff U = -\frac{GMm}{2rR}\dd s \qquad \left(\text{or }\diff\Phi = -\frac{GM}{2rR}\dd s\right)
\end{equation}

\step{Step 3: Case 1 --- Point $P$ Outside the Shell ($r > R$).}
As $\phi$ goes from $0$ to $\pi$, $s$ goes from $s_{\min} = r - R$ to $s_{\max} = r + R$:
\begin{align}
  U(r) &= -\frac{GMm}{2rR}\int_{r-R}^{r+R}\dd s = -\frac{GMm}{2rR}\!\left[(r+R) - (r-R)\right] = \eqbox{-\frac{GMm}{r}} \quad (r > R)\label{eq:U_shell_out}\\
  \Phi(r) &= \eqbox{-\frac{GM}{r}}, \qquad F_r = -\frac{\diff U}{\diff r} = \eqbox{-\frac{GMm}{r^2}}, \qquad g_r = -\frac{\diff\Phi}{\diff r} = \eqbox{-\frac{GM}{r^2}}\label{eq:g_shell_out}
\end{align}

\step{Step 4: Case 2 --- Point $P$ Inside the Shell ($r < R$).}
As $\phi$ goes from $0$ to $\pi$, $s$ goes from $s_{\min} = R - r$ to $s_{\max} = R + r$:
\begin{align}
  U(r) &= -\frac{GMm}{2rR}\int_{R-r}^{R+r}\dd s = -\frac{GMm}{2rR}\!\left[(R+r) - (R-r)\right] = \eqbox{-\frac{GMm}{R} = \text{constant}} \quad (r < R)\label{eq:U_shell_in}\\
  \Phi(r) &= \eqbox{-\frac{GM}{R}}, \qquad F_r = -\frac{\diff U}{\diff r} = \eqbox{0}, \qquad g_r = -\frac{\diff\Phi}{\diff r} = \eqbox{0}\label{eq:g_shell_in}
\end{align}

\section{Gauss's Law for Gravity, Poisson's Equation, and Laplace's Equation\bk{9.6, 9.7}}

\begin{goal}
Derive Gauss's law for gravity in integral form ($\oint_S\vect{g}\cdot\uv{n}\dd S = -4\pi G M_{\text{enc}}$) and differential form ($\nabla\cdot\vect{g} = -4\pi G\rho$), and obtain Poisson's and Laplace's equations for $\Phi$.
\end{goal}

\step{Step 1: Flux of $\vect{g}$ through a closed surface $S$ enclosing point mass $M$.}
Place mass $M$ at the origin. At surface patch $\diff S$ with outward normal $\uv{n}$, $\vect{g} = -\frac{GM}{r^2}\uv{r}$, so:
\begin{equation}
  \vect{g}\cdot\uv{n}\dd S = -\frac{GM}{r^2}(\uv{r}\cdot\uv{n})\dd S = -GM\frac{\diff S\cos\theta}{r^2} = -GM\dd\Omega
\end{equation}
where $\diff\Omega = \frac{\diff S_\perp}{r^2} = \frac{\diff S\cos\theta}{r^2}$ is the solid angle subtended by $\diff S$ at $M$ (\tref{tool:coords}). Since a closed surface surrounding an interior point subtends a total solid angle $\oint_S\dd\Omega = 4\pi\text{ sr}$:
\begin{equation}
  \oint_S \vect{g}\cdot\uv{n}\dd S = -GM\oint_S\dd\Omega = -4\pi GM
\end{equation}
By superposition over all enclosed masses $M_{\text{enc}} = \sum m_i = \int_V \rho\dd\tau$:
\begin{equation}\label{eq:gauss_int}
  \eqbox{\oint_S \vect{g}\cdot\uv{n}\dd S = -4\pi G M_{\text{enc}} = -\int_V 4\pi G\rho\dd\tau} \qquad \text{(Gauss's law for gravity, integral form)}
\end{equation}

\step{Step 2: Apply Gauss's Divergence Theorem (\tref{tool:vec_calc}) to obtain the differential field equations.}
Replace the surface integral in Eq.\ \eqref{eq:gauss_int} by $\int_V (\nabla\cdot\vect{g})\dd\tau$:
\begin{equation}\label{eq:gauss_diff}
  \int_V (\nabla\cdot\vect{g})\dd\tau = -\int_V 4\pi G\rho\dd\tau \implies \eqbox{\nabla\cdot\vect{g} = -4\pi G\rho} \qquad \text{(with }\eqbox{\nabla\times\vect{g} = \vect{0}}\text{)}
\end{equation}

\step{Step 3: Substitute $\vect{g} = -\nabla\Phi$ into $\nabla\cdot\vect{g} = -4\pi G\rho$.}
Since $\nabla\cdot\vect{g} = \nabla\cdot(-\nabla\Phi) = -\nabla^2\Phi$:
\begin{equation}\label{eq:poisson_laplace}
  \eqbox{\nabla^2\Phi = 4\pi G\rho \quad\text{(Poisson's equation)}}, \qquad \eqbox{\nabla^2\Phi = 0 \quad\text{(Laplace's equation in vacuum, }\rho = 0\text{)}}
\end{equation}

\section{Einstein's Theory of Gravitation (General Relativity)\bk{9.8}}

\begin{goal}
Summarize the 4D space-time metric tensor, Minkowski metric, and Einstein field equations.
\end{goal}

\begin{itemize}
  \item \textbf{4D Coordinates and Flat Minkowski Space-Time:} Let $(x^0, x^1, x^2, x^3) = (ct, x, y, z)$. The invariant line element in flat space-time is:
  \begin{equation}\label{eq:minkowski}
    \dd s^2 = \sum_{\mu,\nu=0}^3 \eta_{\mu\nu}\dd x^\mu\dd x^\nu = c^2\dd t^2 - \dd x^2 - \dd y^2 - \dd z^2, \qquad
    \eta_{\mu\nu} = \begin{pmatrix} 1 & 0 & 0 & 0 \\ 0 & -1 & 0 & 0 \\ 0 & 0 & -1 & 0 \\ 0 & 0 & 0 & -1 \end{pmatrix}
  \end{equation}
  \item \textbf{Curved Riemannian Space-Time and Einstein Field Equations:} In curved space-time, $\dd s^2 = \sum_{\mu,\nu=0}^3 g_{\mu\nu}\dd x^\mu\dd x^\nu$. Generalizing Poisson's equation $\nabla^2\Phi = 4\pi G\rho$ by replacing $\rho$ with the symmetric $4\times 4$ energy--momentum tensor $T^{\mu\nu}$ and $\nabla^2\Phi$ with the Einstein curvature tensor $G^{\mu\nu}$ gives:
  \begin{equation}\label{eq:einstein_eq}
    \eqbox{G^{\mu\nu} = -\frac{8\pi G}{c^2}T^{\mu\nu}}, \qquad \eqbox{G^{\mu\nu} = -\frac{8\pi G}{c^2}T^{\mu\nu} - \Lambda g^{\mu\nu}} \quad(\text{with cosmological constant }\Lambda)
  \end{equation}
\end{itemize}

% ============================================================================
% SOLVED BOOK EXAMPLES — MODULE 4
% ============================================================================
\section{Solved Book Examples (Module 4)}

\begin{example}{13.1--13.2}{Cavendish balance force and accelerations\bk{p.\ 405}}
\textbf{Problem:} Two spheres of masses $m_1 = 0.0100\text{ kg}$ and $m_2 = 0.500\text{ kg}$ have centers $r = 0.0500\text{ m}$ apart. Find (a) the gravitational force $F_g$ on each and (b) their accelerations $a_1, a_2$ in free space.
\tcblower
\textbf{Solution:}
\begin{align*}
  \text{(a)}\quad F_g &= \frac{Gm_1 m_2}{r^2} = \frac{(6.67\times 10^{-11})(0.0100)(0.500)}{(0.0500)^2} = \eqbox{1.33\times 10^{-10}\text{ N}}\\
  \text{(b)}\quad a_1 &= \frac{F_g}{m_1} = \eqbox{1.33\times 10^{-8}\text{ m/s}^2}, \qquad a_2 = \frac{F_g}{m_2} = \eqbox{2.66\times 10^{-10}\text{ m/s}^2}
\end{align*}
\end{example}

\begin{example}{13.3}{Superposition of gravitational forces\bk{p.\ 405}}
\textbf{Problem:} A small star ($m = 1.00\times 10^{30}\text{ kg}$) is at the origin $(0,0)$. Two large stars ($M_1 = M_2 = 8.00\times 10^{30}\text{ kg}$) are at $(2.00\times 10^{12}\text{ m}, 2.00\times 10^{12}\text{ m})$ and $(2.00\times 10^{12}\text{ m}, 0)$. Find the magnitude and direction of the total gravitational force $\vect{F}$ on the small star.
\tcblower
\textbf{Solution:}
\begin{align*}
  F_1 &= \frac{(6.67\times 10^{-11})(8.00\times 10^{30})(1.00\times 10^{30})}{(2.00\times 10^{12})^2 + (2.00\times 10^{12})^2} = 6.67\times 10^{25}\text{ N}\quad (\text{at }45^\circ)\\
  F_2 &= \frac{(6.67\times 10^{-11})(8.00\times 10^{30})(1.00\times 10^{30})}{(2.00\times 10^{12})^2} = 1.33\times 10^{26}\text{ N}\quad (\text{along }+x)\\
  F_x &= F_1\cos 45^\circ + F_2 = 1.81\times 10^{26}\text{ N}, \qquad F_y = F_1\sin 45^\circ = 4.72\times 10^{25}\text{ N}\\
  F &= \sqrt{F_x^2 + F_y^2} = \eqbox{1.87\times 10^{26}\text{ N}}, \qquad \theta = \arctan\!\left(\frac{F_y}{F_x}\right) = \eqbox{14.6^\circ}
\end{align*}
\end{example}

\begin{example}{13.4}{Gravity on Mars\bk{p.\ 408}}
\textbf{Problem:} A lander weighs $w = 3430\text{ N}$ on Earth ($m = w/g = 350\text{ kg}$). Find its weight $F_g$ and gravitational acceleration $g_M$ on Mars ($R_M = 3.40\times 10^6\text{ m}$, $m_M = 6.42\times 10^{23}\text{ kg}$).
\tcblower
\textbf{Solution:}
\begin{equation*}
  F_g = \frac{Gm_M m}{R_M^2} = \frac{(6.67\times 10^{-11})(6.42\times 10^{23})(350)}{(3.40\times 10^6)^2} = \eqbox{1.30\times 10^3\text{ N}}, \qquad g_M = \frac{F_g}{m} = \eqbox{3.7\text{ m/s}^2}
\end{equation*}
\end{example}

\begin{example}{13.5}{``From the earth to the moon''\bk{p.\ 410}}
\textbf{Problem:} Find the minimum launch speed from Earth ($R_E = 6.38\times 10^6\text{ m}$, $m_E = 5.97\times 10^{24}\text{ kg}$) to reach (a) $r_2 = 2R_E$ and (b) $r_2 = \infty$ (escape speed).
\tcblower
\textbf{Solution:}
\begin{align*}
  \text{(a)}\quad \frac{1}{2}mv_1^2 - \frac{Gm_E m}{R_E} &= -\frac{Gm_E m}{2R_E} \implies v_1 = \sqrt{\frac{Gm_E}{R_E}} = \eqbox{7900\text{ m/s}}\\
  \text{(b)}\quad \frac{1}{2}mv_1^2 - \frac{Gm_E m}{R_E} &= 0 \implies v_1 = \sqrt{\frac{2Gm_E}{R_E}} = \eqbox{1.12\times 10^4\text{ m/s} = 11.2\text{ km/s}}
\end{align*}
\end{example}

\begin{example}{13.6}{A satellite orbit\bk{p.\ 414}}
\textbf{Problem:} A $1000\text{-kg}$ satellite is in circular orbit $300\text{ km}$ above Earth's surface ($r = 6.68\times 10^6\text{ m}$). Find (a) $v$, $T$, and $a_{\text{rad}}$, (b) work $W_{\text{req}}$ to place it in orbit from rest on the launch pad, and (c) additional work to make it escape to $\infty$.
\tcblower
\textbf{Solution:}
\begin{align*}
  \text{(a)}\quad v &= \sqrt{\frac{Gm_E}{r}} = \eqbox{7720\text{ m/s}}, \qquad T = \frac{2\pi r}{v} = \eqbox{5440\text{ s} = 90.6\text{ min}}, \qquad a_{\text{rad}} = \frac{v^2}{r} = \eqbox{8.92\text{ m/s}^2}\\
  \text{(b)}\quad E_2 &= -\frac{Gm_E m}{2r} = -2.98\times 10^{10}\text{ J}, \qquad E_1 = -\frac{Gm_E m}{R_E} = -6.24\times 10^{10}\text{ J}\\
  W_{\text{req}} &= E_2 - E_1 = \eqbox{3.26\times 10^{10}\text{ J}}\\
  \text{(c)}\quad W_{\text{esc}} &= 0 - E_2 = \eqbox{2.98\times 10^{10}\text{ J}}
\end{align*}
\end{example}

\begin{example}{13.8--13.9}{Kepler's Third Law: Asteroid Pallas and Comet Halley\bk{p.\ 417}}
\textbf{Problem:} (a) Pallas has period $T = 4.62\text{ yr} = 1.46\times 10^8\text{ s}$ around the Sun ($m_S = 1.99\times 10^{30}\text{ kg}$); find $a$. (b) Comet Halley has $r_p = 8.75\times 10^7\text{ km}$ and $r_a = 5.26\times 10^9\text{ km}$; find $a$, $e$, and $T$.
\tcblower
\textbf{Solution:}
\begin{align*}
  \text{(a)}\quad a &= \left(\frac{Gm_S T^2}{4\pi^2}\right)^{1/3} = \eqbox{4.15\times 10^{11}\text{ m}}\\
  \text{(b)}\quad a &= \frac{r_p + r_a}{2} = \eqbox{2.67\times 10^9\text{ km}}, \qquad e = 1 - \frac{r_p}{a} = \eqbox{0.967}, \qquad T = \frac{2\pi a^{3/2}}{\sqrt{Gm_S}} = \eqbox{75.5\text{ yr}}
\end{align*}
\end{example}

\begin{example}{13.10}{``Journey to the center of the earth''\bk{p.\ 420}}
\textbf{Problem:} Assuming uniform Earth density, find the gravitational force $F_g(r)$ on a mass $m$ at distance $r < R_E$ from Earth's center.
\tcblower
\textbf{Solution:} Only the mass $M(r) = m_E\frac{r^3}{R_E^3}$ inside radius $r$ contributes to the force:
\begin{equation*}
  F_g(r) = \frac{GM(r)m}{r^2} = \frac{Gm}{r^2}\!\left(m_E\frac{r^3}{R_E^3}\right) = \eqbox{\frac{Gm_E m}{R_E^3}r = mg\frac{r}{R_E}}
\end{equation*}
\end{example}

\begin{example}{13.11}{Black hole calculations\bk{p.\ 424}}
\textbf{Problem:} Find the Schwarzschild radius $R_S$ and average density $\rho$ at the event horizon for a black hole of mass $M = 3M_\odot = 6.0\times 10^{30}\text{ kg}$.
\tcblower
\textbf{Solution:}
\begin{equation*}
  R_S = \frac{2GM}{c^2} = \frac{2(6.67\times 10^{-11})(6.0\times 10^{30})}{(3.00\times 10^8)^2} = \eqbox{8.9\times 10^3\text{ m} = 8.9\text{ km}}, \qquad \rho = \frac{M}{\frac{4}{3}\pi R_S^3} = \eqbox{2.0\times 10^{18}\text{ kg/m}^3}
\end{equation*}
\end{example}

\begin{example}{Ch.\ 13}{Bridging Problem: Speeds in an Elliptical Orbit\bk{p.\ 428}}
\textbf{Problem:} A comet orbits the Sun ($m_S$) in an ellipse of semi-major axis $a$ and eccentricity $e$. (a) Derive expressions for the speeds $v_p$ at perihelion and $v_a$ at aphelion. (b) Evaluate for Comet Halley ($a = 2.67\times 10^{12}\text{ m}$, $e = 0.967$).
\tcblower
\textbf{Solution:} At perihelion $r_p = a(1-e)$ and aphelion $r_a = a(1+e)$, angular momentum conservation ($\phi = 90^\circ$) gives $r_p v_p = r_a v_a \implies v_a = \frac{1-e}{1+e}v_p$. Energy conservation between $r_p$ and $r_a$:
\begin{align*}
  \frac{1}{2}m(v_p^2 - v_a^2) &= Gm_S m\!\left(\frac{1}{a(1-e)} - \frac{1}{a(1+e)}\right) = \frac{2Gm_S m e}{a(1-e^2)}\\
  v_p^2\!\left[1 - \left(\frac{1-e}{1+e}\right)^2\right] &= v_p^2\frac{4e}{(1+e)^2} = \frac{4Gm_S e}{a(1-e)(1+e)}\\
  \text{(a)}\quad \eqbox{v_p = \sqrt{\frac{Gm_S(1+e)}{a(1-e)}}}, &\qquad \eqbox{v_a = \sqrt{\frac{Gm_S(1-e)}{a(1+e)}}}\\
  \text{(b)}\quad \text{For Comet Halley: } v_p &= \eqbox{5.46\times 10^4\text{ m/s} = 54.6\text{ km/s}}, \qquad v_a = \eqbox{916\text{ m/s} = 0.916\text{ km/s}}
\end{align*}
\end{example}

\begin{example}{9.1}{Gravitational field of a uniform thin rod on its axis\bk{Analytical Mechanics p.\ 217}}
\textbf{Problem:} A uniform thin rod of length $L$ and linear mass density $\lambda$ lies along the $x$-axis from $x'=0$ to $x'=L$. Find the gravitational field $\vect{g}$ at $x = 4L$.
\tcblower
\textbf{Solution:} Element $\diff m = \lambda\dd x'$ at $x'$ produces field $\diff\vect{g} = -\frac{G\lambda\dd x'}{(4L-x')^2}\ihat$ at $x=4L$:
\begin{equation*}
  \vect{g} = -G\lambda\ihat\int_0^L \frac{\diff x'}{(4L-x')^2} = -G\lambda\ihat\!\left[\frac{1}{4L-x'}\right]_0^L = -G\lambda\ihat\!\left(\frac{1}{3L} - \frac{1}{4L}\right) = \eqbox{-\frac{1}{12L}\lambda G\ihat}
\end{equation*}
\end{example}

\begin{example}{9.2}{Gravitational field on the axis of a uniform thin disk\bk{Analytical Mechanics p.\ 218}}
\textbf{Problem:} Find the gravitational field $\vect{g}$ at distance $+z$ along the symmetry axis of a uniform thin disk of mass $M$, radius $R$, and surface mass density $\sigma = \frac{M}{\pi R^2}$.
\tcblower
\textbf{Solution:} By symmetry, horizontal components cancel. A ring of radius $r'$ and width $\diff r'$ has area $\diff A' = 2\pi r'\dd r'$, distance $D = \sqrt{r'^2+z^2}$ to $P$, and $\cos\theta = \frac{z}{\sqrt{r'^2+z^2}}$:
\begin{align*}
  \vect{g} &= -G\sigma\khat\int_0^R \frac{2\pi r'\dd r'}{r'^2+z^2}\frac{z}{\sqrt{r'^2+z^2}} = -2\pi\sigma G z\khat\int_0^R \frac{r'\dd r'}{(r'^2+z^2)^{3/2}}\\
           &= -2\pi\sigma G z\khat\!\left[-\frac{1}{\sqrt{r'^2+z^2}}\right]_0^R = \eqbox{-2\pi\sigma G\!\left(1 - \frac{z}{\sqrt{z^2+R^2}}\right)\khat = -\frac{2GM}{R^2}\!\left(1 - \frac{z}{\sqrt{z^2+R^2}}\right)\khat}
\end{align*}
\end{example}

\begin{example}{9.4}{Gravitational field of an infinite cylinder via Gauss's Law\bk{Analytical Mechanics p.\ 224--225}}
\textbf{Problem:} Use Gauss's law to find the gravitational field $\vect{g}$ outside ($r > a$) an infinitely long cylinder of radius $a$ and linear mass density $\lambda$.
\tcblower
\textbf{Solution:} Choose a coaxial cylindrical Gaussian surface of radius $r > a$ and length $L$. By symmetry, $\vect{g} = -g(r)\uv{\rho}$ is radial and zero flux passes through the end caps. Enclosed mass is $M_{\text{enc}} = \lambda L$:
\begin{equation*}
  \oint_S \vect{g}\cdot\uv{n}\dd S = -g(2\pi r L) = -4\pi G M_{\text{enc}} = -4\pi G(\lambda L) \implies \eqbox{\vect{g} = -\frac{2G\lambda}{r}\uv{\rho}}
\end{equation*}
\end{example}

% ============================================================================
% UNIQUE SOLVED PROBLEMS — MODULE 4
% ============================================================================
\section{Unique Solved Problems (Module 4)}

\subsubsection{Part A: Analytical Mechanics Chapter 9 Exercises (with Complete Derivations)}

\begin{enumerate}[label=\textbf{Ex.\ 9.\arabic*},leftmargin=3.5em]
  \item \textbf{(Mars--Sun Force and Accelerations):} $M_S = 1.99\times 10^{30}\text{ kg}$, $M_{\text{Mars}} = 6.42\times 10^{23}\text{ kg}$, $r = 2.28\times 10^{11}\text{ m}$:
  \begin{equation*}
    F = \frac{GM_S M_{\text{Mars}}}{r^2} = \eqbox{1.64\times 10^{21}\text{ N}}, \quad a_{\text{Mars}} = \frac{F}{M_{\text{Mars}}} = \eqbox{2.55\times 10^{-3}\text{ m/s}^2}, \quad a_S = \frac{F}{M_S} = \eqbox{8.24\times 10^{-10}\text{ m/s}^2}
  \end{equation*}

  \item \textbf{(Mass Independence of Free Fall):} $m\vect{a} = \vect{F}_g = m\vect{g} \implies \vect{a} = \vect{g} = -\frac{GM_E}{r^2}\uv{r}$, which is independent of $m$.

  \item \textbf{(Three Point Masses):} $m_1 = m$ at $(1,2)$, $m_2 = 2m$ at $(3,4)$, $m_3 = 3m$ at $(-2,-2)$. Force on $m_2$:
  \begin{align*}
    \vect{r}_2 - \vect{r}_1 &= 2\ihat + 2\jhat \quad(|\vect{r}_2-\vect{r}_1| = \sqrt{8}), \qquad \vect{r}_2 - \vect{r}_3 = 5\ihat + 6\jhat \quad(|\vect{r}_2-\vect{r}_3| = \sqrt{61})\\
    \vect{F}_2 &= -Gm(2m)\frac{2\ihat+2\jhat}{8\sqrt{8}} - G(3m)(2m)\frac{5\ihat+6\jhat}{61\sqrt{61}} = \eqbox{-Gm^2(0.24\ihat + 0.25\jhat)}
  \end{align*}

  \item \textbf{(Curl of Newton's Law of Gravitation):} For $\vect{F} = -GMm\frac{x\ihat+y\jhat+z\khat}{(x^2+y^2+z^2)^{3/2}}$, $(\nabla\times\vect{F})_z = \frac{\partial F_y}{\partial x} - \frac{\partial F_x}{\partial y} = -GMm\!\left[\frac{3xy}{r^5} - \frac{3yx}{r^5}\right] = 0$, and cyclic permutations give $\eqbox{\nabla\times\vect{F} = \vect{0}}$.

  \item \textbf{(Field on Perpendicular Bisector of Two Masses $M$ Separated by $a$):} At $(0,0,z)$ with sources at $(\pm a/2, 0, 0)$, distance is $D = \sqrt{(a/2)^2 + z^2}$ and $\cos\theta = \frac{z}{D}$:
  \begin{equation*}
    \vect{g} = -2\frac{GM}{D^2}\cos\theta\,\uv{z} = \eqbox{-\frac{2GMz}{\left[(a/2)^2 + z^2\right]^{3/2}}\uv{z}}
  \end{equation*}

  \item \textbf{(Two Masses $M$ at $(0,0)$ and $(0,4)$):}
  \begin{itemize}
    \item At $(0,3)$: $\vect{g} = -GM\!\left(\frac{\jhat}{3^2} - \frac{\jhat}{1^2}\right) = \eqbox{+\frac{8}{9}GM\jhat = 0.89GM\jhat}$.
    \item At $(3,0)$: $\vect{r}-\vect{r}_1 = 3\ihat$ ($r_1=3$), $\vect{r}-\vect{r}_2 = 3\ihat - 4\jhat$ ($r_2=5$):
    \begin{equation*}
      \vect{g}(3,0) = -GM\!\left[\frac{3\ihat}{27} + \frac{3\ihat - 4\jhat}{125}\right] = -GM\!\left[\left(\frac{1}{9}+\frac{3}{125}\right)\ihat - \frac{4}{125}\jhat\right] = \eqbox{GM(-0.14\ihat + 0.03\jhat)}
    \end{equation*}
  \end{itemize}

  \setcounter{enumi}{7}
  \item \textbf{(Field on the Axis of a Ring of Mass $M$ and Radius $R$):} Every element $\diff m$ is at distance $D = \sqrt{z^2+R^2}$ with $\cos\theta = \frac{z}{\sqrt{z^2+R^2}}$, so $g_z = -\int\frac{G\dd m}{z^2+R^2}\frac{z}{\sqrt{z^2+R^2}} = \eqbox{-\frac{GMz}{(z^2+R^2)^{3/2}}}$.

  \item \textbf{(Disk as Concentric Rings):} A ring of radius $r'$ and width $\diff r'$ has mass $\diff M = \sigma(2\pi r'\dd r')$. Using Exercise 9.8:
  \begin{equation*}
    g_z = -\int_0^R \frac{G(2\pi\sigma r'\dd r')z}{(z^2+r'^2)^{3/2}} = -2\pi\sigma G z\!\left[-\frac{1}{\sqrt{z^2+r'^2}}\right]_0^R = \eqbox{-2\pi\sigma G\!\left(1 - \frac{z}{\sqrt{z^2+R^2}}\right)}
  \end{equation*}

  \item \textbf{(Potential on the Axis of a Thin Disk of Radius $R$, Density $\rho$, and Thickness $t$):} With $\sigma = \rho t$ and ring mass $\diff m = 2\pi\rho t r'\dd r'$ at distance $\sqrt{z^2+r'^2}$:
  \begin{equation*}
    \Phi(z) = -G\int_0^R \frac{2\pi\rho t r'\dd r'}{\sqrt{z^2+r'^2}} = -2\pi\rho t G\!\left[\sqrt{z^2+r'^2}\right]_0^R = \eqbox{-2\pi\rho t G\!\left[(z^2+R^2)^{1/2} - z\right]}
  \end{equation*}

  \item \textbf{(Solid Angle of a Spherical Shell):} By Eq.\ (9.11), $\oint_S\dd\Omega = \int_0^{2\pi}\dd\phi\int_0^\pi\sin\theta\dd\theta = 2\pi(2) = \eqbox{4\pi\text{ sr}}$.

  \item \textbf{(Point Mass Field via Gauss's Law):} For a spherical Gaussian surface of radius $r$ centered on $M$, $\oint_S\vect{g}\cdot\uv{n}\dd S = g_r(4\pi r^2) = -4\pi GM \implies \eqbox{\vect{g} = -\frac{GM}{r^2}\uv{r}}$.

  \item \textbf{(Sources of $\vect{E}$ and $\vect{B}$ in Maxwell's Equations):} By Helmholtz's theorem, the sources of a vector field are its divergence and curl. Thus the sources of $\vect{E}$ are charge density $\rho$ and time-varying magnetic fields $-\frac{\partial\vect{B}}{\partial t}$; the sources of $\vect{B}$ are current density $\vect{J}$ and time-varying electric fields $\frac{\partial\vect{E}}{\partial t}$ (displacement current).
\end{enumerate}

\subsubsection{Part B: Section 9.10 End-of-Chapter Problems (PDF 2, pp.\ 231--234)}

\begin{enumerate}[label=\textbf{Prob.\ 9.\arabic*},leftmargin=4.2em]
  \setcounter{enumi}{3}
  \item \textbf{[Prob.\ 9.4, p.\ 231] Time for a Stopped Earth to Fall into the Sun:}
  An object falling from rest at $r_0 = a_{\text{orbit}}$ directly into the Sun follows a degenerate ellipse with semi-major axis $a_{\text{fall}} = \frac{r_0}{2}$ ($e\to 1$). By Kepler's third law ($T \propto a^{3/2}$), a full fall-and-return period is $T_{\text{fall}} = T_{\text{Earth}}\!\left(\frac{1}{2}\right)^{3/2}$, so the one-way collision time is half of $T_{\text{fall}}$:
  \begin{equation*}
    t_{\text{collide}} = \frac{1}{2}\!\left(\frac{1}{2}\right)^{3/2} T_{\text{Earth}} = \frac{T_{\text{Earth}}}{4\sqrt{2}} = \frac{365.25\text{ days}}{5.657} = \eqbox{64.6\text{ days}}
  \end{equation*}

  \setcounter{enumi}{6}
  \item \textbf{[Prob.\ 9.7, p.\ 231] Surface-Skimming Satellite Period Depends Only on Density $\rho$:}
  For a circular orbit at $r = R$ around a uniform spherical planet of mass $M = \frac{4}{3}\pi R^3\rho$:
  \begin{equation*}
    T = \frac{2\pi R^{3/2}}{\sqrt{GM}} = \frac{2\pi R^{3/2}}{\sqrt{G\!\left(\frac{4}{3}\pi R^3\rho\right)}} = \eqbox{\sqrt{\frac{3\pi}{G\rho}}}
  \end{equation*}
  Because $R$ cancels completely, a planet of radius $2R$ with the same density $\rho$ has the \textbf{exact same} surface-skimming orbital period $T$.

  \item \textbf{[Prob.\ 9.8 \& 9.18, p.\ 231--232] Infinite Flat Sheet of Surface Mass Density $\sigma$:}
  \begin{itemize}
    \item \textbf{(a) Via Gauss's Law (Prob.\ 9.18a) or Integration (Prob.\ 9.8a):} Using a cylindrical pillbox of cross-sectional area $A$ straddling the sheet ($M_{\text{enc}} = \sigma A$), flux through the two faces is $-2gA = -4\pi G(\sigma A) \implies \eqbox{g = 2\pi G\sigma}$ (uniform, pointing toward the sheet).
    \item \textbf{(b) Force on spherical shell of mass $M$ at height $b > a$:} $\eqbox{F = Mg = 2\pi G\sigma M}$.
    \item \textbf{(c) Comparison with spherical shell of surface density $\sigma$ (Prob.\ 9.18b--c):} Just outside a sphere of mass $M = 4\pi a^2\sigma$, $g_{\text{sphere}} = \frac{G(4\pi a^2\sigma)}{a^2} = \eqbox{4\pi G\sigma = 2g_{\text{sheet}}}$.
  \end{itemize}

  \item \textbf{[Prob.\ 9.9, p.\ 231] Gravitational Pressure on an Outer Concentric Spherical Shell:}
  Inner sphere has radius $a, m_a$; outer shell has radius $b > a$ and mass $m_b$ ($\sigma_b = \frac{m_b}{4\pi b^2}$). The field at radius $b$ is the sum of the field $g_a = \frac{Gm_a}{b^2}$ from the inner shell and the self-field $\frac{1}{2}g_b = \frac{Gm_b}{2b^2}$ of the outer shell on itself. Thus the inward force per unit area (pressure) on the outer sphere is:
  \begin{equation*}
    P = \sigma_b\!\left(\frac{Gm_a}{b^2} + \frac{Gm_b}{2b^2}\right) = \eqbox{\frac{Gm_b(m_a + m_b/2)}{4\pi b^4}} \qquad \left(\text{or }\frac{Gm_a m_b}{4\pi b^4}\text{ due to }m_a\text{ alone}\right)
  \end{equation*}

  \setcounter{enumi}{10}
  \item \textbf{[Prob.\ 9.11 \& 9.14, p.\ 231] Field and Potential on the Axis of a Right Circular Cylinder ($M, L, a$):}
  Center of cylinder is at $z'=0$ ($z'\in [-L/2, L/2]$); field point is at $z > L/2$. Distance to disk at $z'$ is $t = z - z'$ ($\diff t = -\diff z'$, $\diff m = \frac{M}{L}\dd z'$).
  \begin{itemize}
    \item \textbf{Prob.\ 9.11 (Field $g_z$):} Integrate Worked Example 9.2 over $t\in [z-L/2, z+L/2]$:
    \begin{align*}
      g_z &= -\frac{2GM}{La^2}\int_{z-L/2}^{z+L/2}\left(1 - \frac{t}{\sqrt{t^2+a^2}}\right)\dd t = \eqbox{-\frac{2GM}{La^2}\!\left[L - \sqrt{(z+L/2)^2+a^2} + \sqrt{(z-L/2)^2+a^2}\right]}
    \end{align*}
    \item \textbf{Prob.\ 9.14 (Potential $\Phi$):} Integrate $\diff\Phi_{\text{disk}} = \frac{2G\dd m}{a^2}\!\left[t - \sqrt{a^2+t^2}\right]$ from $t_1 = z - L/2$ to $t_2 = z + L/2$ using $\int\sqrt{a^2+t^2}\dd t = \frac{1}{2}\!\left[t\sqrt{a^2+t^2} + a^2\ln\!\left(t+\sqrt{a^2+t^2}\right)\right]$:
    \begin{equation*}
      \eqbox{\Phi(z) = \frac{2GM}{La^2}\!\left[zL - \frac{1}{2}\!\left(t\sqrt{a^2+t^2} + a^2\ln\!\left(t+\sqrt{a^2+t^2}\right)\right)_{t_1 = z-L/2}^{t_2 = z+L/2}\right]}
    \end{equation*}
  \end{itemize}

  \setcounter{enumi}{16}
  \item \textbf{[Prob.\ 9.17 \& 9.20, p.\ 232] Field and Potential Inside and Outside a Uniform Sphere ($M, R$):}
  \begin{itemize}
    \item \textbf{Outside ($r \ge R$):} $M_{\text{enc}} = M \implies \eqbox{\vect{g}(r) = -\frac{GM}{r^2}\uv{r}}$, and $\eqbox{\Phi(r) = -\frac{GM}{r}}$.
    \item \textbf{Inside ($r < R$):} $M_{\text{enc}} = M\frac{r^3}{R^3} \implies 4\pi r^2 g_r = -4\pi G M\frac{r^3}{R^3} \implies \eqbox{\vect{g}(r) = -\frac{GMr}{R^3}\uv{r}}$.
    \item Integrating $\Phi(r) - \Phi(R) = -\int_R^r g_{r'}\dd r' = \int_R^r \frac{GMr'}{R^3}\dd r' = \frac{GM}{2R^3}(r^2 - R^2)$ with $\Phi(R) = -\frac{GM}{R}$ gives:
    \begin{equation*}
      \eqbox{\Phi(r) = -\frac{GM}{2R^3}(3R^2 - r^2) \qquad (r < R)}
    \end{equation*}
  \end{itemize}

  \setcounter{enumi}{18}
  \item \textbf{[Prob.\ 9.19, p.\ 232] Intergalactic Gas Cloud with $\rho(r) = \frac{Mb^2}{2\pi r(r^2+b^2)^2}$:}
  \begin{itemize}
    \item Enclosed mass inside radius $r$:
    \begin{equation*}
      M_{\text{enc}}(r) = \int_0^r \frac{Mb^2}{2\pi r'(r'^2+b^2)^2}(4\pi r'^2\dd r') = 2Mb^2\int_0^r \frac{r'\dd r'}{(r'^2+b^2)^2} = 2Mb^2\!\left[-\frac{1}{2(r'^2+b^2)}\right]_0^r = \frac{Mr^2}{r^2+b^2}
    \end{equation*}
    \item By Gauss's law, $-4\pi r^2 g_r = -4\pi G M_{\text{enc}}(r)$:
    \begin{equation*}
      \eqbox{\vect{g}(r) = -\frac{GM}{r^2+b^2}\uv{r}}, \qquad \Phi(r) = \int_\infty^r g_{r'}\dd r' = \eqbox{-\frac{GM}{b}\!\left[\frac{\pi}{2} - \arctan\!\left(\frac{r}{b}\right)\right] = -\frac{GM}{b}\cot^{-1}\!\left(\frac{r}{b}\right)}
    \end{equation*}
  \end{itemize}

  \setcounter{enumi}{22}
  \item \textbf{[Prob.\ 9.23, p.\ 232] Sphere of Radius $R$ with Off-Center Spherical Cavity of Radius $R/2$:}
  Full sphere has mass $M_{\text{full}} = \frac{4}{3}\pi R^3\rho$ centered at $O$. Cavity has radius $R/2$, volume $\frac{1}{8}V_{\text{full}}$, mass $M_{\text{cav}} = \frac{1}{8}M_{\text{full}} = \frac{1}{6}\pi R^3\rho$, centered at $z_c = R/4$. At $z > R$ on the line of centers, superpose the full sphere at $z=0$ and negative cavity mass at $z = R/4$:
  \begin{equation*}
    F_z = -Gm\!\left[\frac{M_{\text{full}}}{z^2} - \frac{M_{\text{cav}}}{(z - R/4)^2}\right] = \eqbox{-\frac{4}{3}\pi G\rho R^3 m\!\left[\frac{1}{z^2} - \frac{1}{8(z - R/4)^2}\right]}
  \end{equation*}

  \item \textbf{[Prob.\ 9.24 \& 9.29, p.\ 232--233] Quadrupole Potential of a Ring of Mass $m$ and Radius $a$ at $(r,\theta)$ ($r \gg a$):}
  Place the ring in the $xy$-plane ($\theta_1 = \pi/2$, $\phi_1 = \phi'$). For a field point at $(r,\theta,0)$, $\cos\gamma = \sin\theta\cos\phi'$. The distance from ring element $\diff m = \frac{m}{2\pi}\dd\phi'$ to the field point is:
  \begin{align*}
    |\vect{r}-\vect{r}'|^{-1} &= r^{-1}\!\left(1 - \frac{2a}{r}\sin\theta\cos\phi' + \frac{a^2}{r^2}\right)^{-1/2}\\
    &\approx \frac{1}{r}\!\left[1 + \frac{a}{r}\sin\theta\cos\phi' + \frac{a^2}{2r^2}(3\sin^2\theta\cos^2\phi' - 1) + \mathcal{O}\!\left(\frac{a^3}{r^3}\right)\right]
  \end{align*}
  Integrate over $\phi'\in [0, 2\pi]$ using $\int_0^{2\pi}\cos\phi'\dd\phi' = 0$ and $\int_0^{2\pi}\cos^2\phi'\dd\phi' = \pi$:
  \begin{equation*}
    \eqbox{\Phi(r,\theta) = -\frac{Gm}{r}\!\left[1 + \frac{a^2}{4r^2}(3\sin^2\theta - 2)\right] = -\frac{Gm}{r}\!\left[1 - \frac{a^2}{2r^2}P_2(\cos\theta)\right]}
  \end{equation*}
  In the plane of the ring ($\theta = \pi/2$, $\sin\theta = 1$, Prob.\ 9.29 with mass $M$), this gives $\eqbox{\Phi(r,\pi/2) = -\frac{GM}{r}\!\left(1 + \frac{a^2}{4r^2}\right)}$.

  \item \textbf{[Prob.\ 9.25, p.\ 233] Simple Harmonic Motion in a Tunnel Through Earth:}
  From Example 13.10, at distance $r$ from Earth's center, the restoring force is $F_r = -\frac{GM_E m}{R_E^3}r = -kr$ where $k = \frac{GM_E m}{R_E^3} = \frac{mg}{R_E}$.
  \begin{itemize}
    \item \textbf{(a)} The equation of motion is:
    \begin{equation*}
      m\ddot{r} + \frac{mg}{R_E}r = 0 \implies \eqbox{\ddot{r} + \omega^2 r = 0, \qquad \omega = \sqrt{\frac{g}{R_E}}}
    \end{equation*}
    \item \textbf{(b)} The period of oscillation is:
    \begin{equation*}
      \eqbox{T = \frac{2\pi}{\omega} = 2\pi\sqrt{\frac{R_E}{g}}} = 2\pi\sqrt{\frac{6.38\times 10^6\text{ m}}{9.80\text{ m/s}^2}} = \eqbox{5070\text{ s} \approx 84.5\text{ min}}
    \end{equation*}
    This is identical to the period of a satellite in a surface-skimming circular orbit.
    \item \textbf{(c)} The one-way transit time through Earth is:
    \begin{equation*}
      \eqbox{t_{\text{transit}} = \frac{T}{2} \approx 42.2\text{ min}}
    \end{equation*}
  \end{itemize}

  \item \textbf{[Prob.\ 9.26, p.\ 233] Gravitational Self-Energy of a Uniform Solid Sphere:}
  Assemble a uniform solid sphere of mass $M$, radius $R$, and density $\rho = \frac{3M}{4\pi R^3}$ layer by layer from $r = 0$ to $r = R$. Bringing a spherical shell of radius $r$, thickness $\diff r$, and mass $\diff m = 4\pi\rho r^2\dd r$ from infinity to the surface of the core of mass $M(r) = \frac{4}{3}\pi\rho r^3$ requires gravitational work:
  \begin{equation*}
    \diff U = -\frac{GM(r)\dd m}{r} = -\frac{G\left(\frac{4}{3}\pi\rho r^3\right)(4\pi\rho r^2\dd r)}{r} = -\frac{16}{3}\pi^2 G\rho^2 r^4\dd r
  \end{equation*}
  Integrating from $r = 0$ to $r = R$:
  \begin{equation*}
    \eqbox{U_{\text{self}} = -\frac{16}{3}\pi^2 G\rho^2\int_0^R r^4\dd r = -\frac{16}{15}\pi^2 G\rho^2 R^5 = -\frac{3}{5}\frac{GM^2}{R}}
  \end{equation*}

  \item \textbf{[Young \& Freedman Prob.\ 13.71, p.\ 433] Secant Tunnel Through a Rotating Planet:}
  A straight frictionless tunnel is drilled through Earth along a chord at perpendicular distance $d$ from Earth's center.
  \begin{itemize}
    \item \textbf{(a)} When a body of mass $m$ is at distance $x$ along the tunnel from the midpoint, its distance from the center is $r = \sqrt{x^2 + d^2}$. The radial gravitational force points toward the center: $\vect{F}_g = -\frac{mg}{R_E}\vect{r}$.
    \item \textbf{(b)} The component of force along the tunnel is $F_x = F_g\cos\theta = -\left(\frac{mg}{R_E}r\right)\!\left(\frac{x}{r}\right) = \eqbox{-\frac{mg}{R_E}x}$.
    \item \textbf{(c)} Since $F_x = -kx$ with $k = \frac{mg}{R_E}$ independent of $d$, the oscillation period along \textbf{any} chord through Earth is:
    \begin{equation*}
      \eqbox{T = 2\pi\sqrt{\frac{R_E}{g}} \approx 84.5\text{ min}}
    \end{equation*}
  \end{itemize}

  \item \textbf{[Young \& Freedman Prob.\ 13.83, p.\ 434] Hohmann Transfer Orbit to Mars:}
  A spacecraft transfers from Earth's circular orbit ($r_E = 1.50\times 10^{11}\text{ m}$) to Mars's orbit ($r_M = 2.28\times 10^{11}\text{ m}$) along an elliptical orbit with perihelion at Earth and aphelion at Mars ($m_S = 1.99\times 10^{30}\text{ kg}$).
  \begin{itemize}
    \item \textbf{(a)} Semi-major axis: $a = \frac{r_E + r_M}{2} = \frac{1.50\times 10^{11} + 2.28\times 10^{11}}{2} = \eqbox{1.89\times 10^{11}\text{ m}}$.
    \item \textbf{(b)} The transfer time is half the orbital period:
    \begin{equation*}
      \eqbox{t_{\text{transfer}} = \frac{T}{2} = \frac{\pi a^{3/2}}{\sqrt{Gm_S}} = \frac{\pi(1.89\times 10^{11})^{3/2}}{\sqrt{(6.67\times 10^{-11})(1.99\times 10^{30})}} = 2.24\times 10^7\text{ s} \approx 259\text{ days} \approx 8.5\text{ months}}
    \end{equation*}
    \item \textbf{(c)} Speed needed at perihelion: $v_p = \sqrt{Gm_S\!\left(\frac{2}{r_E} - \frac{1}{a}\right)} = \eqbox{32.7\text{ km/s}}$ (an impulse of $\Delta v = 32.7 - 29.8 = 2.9\text{ km/s}$ relative to Earth's orbit).
  \end{itemize}
\end{enumerate}

% ============================================================================
% FORMULA SHEET — MODULE 4
% ============================================================================
\formulasheet{Module 4 Formula Sheet: Gravitation and Gravitational Fields}

\begin{fs}{1. Newton's Law of Gravitation, Satellites, and Kepler's Laws}
\begin{align*}
  \vect{F}_{21} &= -G\frac{m_1 m_2}{r^2}\uv{r}, \qquad g = \frac{GM_E}{R_E^2}, \qquad U(r) = -\frac{GMm}{r}, \qquad v_{\text{esc}} = \sqrt{\frac{2GM}{R}}\\
  v_{\text{orb}} &= \sqrt{\frac{GM}{r}}, \qquad T = \frac{2\pi r^{3/2}}{\sqrt{GM}}, \qquad E = -\frac{GMm}{2r} = -K = \frac{U}{2}\\
  \text{Kepler's Laws:}\quad & r_p = a(1-e), \quad r_a = a(1+e), \quad \frac{\diff A}{\diff t} = \frac{L}{2m} = \text{const}, \quad T^2 = \left(\frac{4\pi^2}{GM_S}\right)a^3, \quad R_S = \frac{2GM}{c^2}
\end{align*}
\end{fs}

\begin{fs}{2. Gravitational Field, Gauss's Law, Poisson's and Laplace's Equations}
\begin{align*}
  \vect{g} &= -\nabla\Phi, \qquad \Phi(\vect{r}) = -G\int\frac{\rho(\vect{r}')\dd\tau'}{|\vect{r}-\vect{r}'|}\\
  \text{Spherical Shell:}\quad & \Phi(r) = -\frac{GM}{r},\ \vect{g} = -\frac{GM}{r^2}\uv{r}\ (r > R); \qquad \Phi(r) = -\frac{GM}{R},\ \vect{g} = \vect{0}\ (r < R)\\
  \text{Solid Sphere:}\quad & g(r) = -\frac{GM}{r^2}\ (r \ge R); \qquad g(r) = -\frac{GMr}{R^3}\ (r < R)\\
  \text{Gauss's Law:}\quad & \oint_S\vect{g}\cdot\uv{n}\dd S = -4\pi GM_{\text{enc}} \iff \nabla\cdot\vect{g} = -4\pi G\rho, \qquad \nabla\times\vect{g} = \vect{0}\\
  \text{Field Equations:}\quad & \nabla^2\Phi = 4\pi G\rho\ (\text{Poisson}), \qquad \nabla^2\Phi = 0\ (\text{Laplace}), \qquad G^{\mu\nu} = -\frac{8\pi G}{c^2}T^{\mu\nu}\ (\text{Einstein})
\end{align*}
\end{fs}

\end{document}
